\input{preamble}

% OK, start here.
%
\begin{document}

\title{Cohomologie cristalline}


\maketitle

\phantomsection
\label{section-phantom}

\tableofcontents



\section{Introduction}
\label{section-introduction}

\noindent
Ce chapitre s'appuie sur une série de cours donnés par Johan de Jong
en 2012 à l'université Columbia.
L'objectif de ce chapitre est de donner une introduction rapide à la
cohomologie cristalline. Le livre \cite{Berthelot} fournit une référence.

\medskip\noindent
Nous avons placé l'étude purement algébrique et plus élémentaire des anneaux
à puissances divisées dans un chapitre préliminaire, car elle sert aussi
à étudier les résolutions de Tate en algèbre commutative.
Voir Algèbre à puissances divisées, section \ref{dpa-section-introduction}.













\section{Enveloppe à puissances divisées}
\label{section-divided-power-envelope}

\noindent
La construction du lemme suivant sera appelée l'enveloppe à
puissances divisées. Elle jouera plus loin un rôle important.

\begin{lemma}
\label{lemma-divided-power-envelope}
Soit $(A, I, \gamma)$ un anneau à puissances divisées.
Soit $A \to B$ un homomorphisme d'anneaux. Soit $J \subset B$ un idéal
tel que $IB \subset J$. Il existe un homomorphisme d'anneaux
à puissances divisées
$$
(A, I, \gamma) \longrightarrow (D, \bar J, \bar \gamma)
$$
tel que
$$
\Hom_{(A, I, \gamma)}((D, \bar J, \bar \gamma), (C, K, \delta)) =
\Hom_{(A, I)}((B, J), (C, K))
$$
fonctoriellement en l'algèbre à puissances divisées $(C, K, \delta)$ sur
$(A, I, \gamma)$. Ici, le membre de gauche désigne les morphismes d'anneaux
à puissances divisées sur $(A, I, \gamma)$, et celui de droite les morphismes
de couples (anneau, idéal) sur $(A, I)$.
\end{lemma}

\begin{proof}
Notons $\mathcal{C}$ la catégorie des anneaux à puissances divisées
$(C, K, \delta)$. Considérons le foncteur
$F : \mathcal{C} \longrightarrow \textit{Ens}$ défini par
$$
F(C, K, \delta) =
\left\{
(\varphi, \psi)
\middle|
\begin{matrix}
\varphi : (A, I, \gamma) \to (C, K, \delta)
\text{ homomorphisme d'anneaux à puissances divisées} \\
\psi : (B, J) \to (C, K)\text{ un homomorphisme de }
A\text{-algèbres tel que }\psi(J) \subset K
\end{matrix}
\right\}
$$
Nous allons montrer que
le lemme \ref{dpa-lemma-a-version-of-brown} d'Algèbre à puissances divisées
s'applique à ce foncteur, ce qui
démontrera le lemme. Supposons $(\varphi, \psi) \in F(C, K, \delta)$.
Soit $C' \subset C$ le sous-anneau engendré par $\varphi(A)$,
$\psi(B)$ et les $\delta_n(\psi(f))$ pour tout $f \in J$.
Soit $K' \subset K \cap C'$ l'idéal de $C'$ engendré par
$\varphi(I)$ et les $\delta_n(\psi(f))$, pour $f \in J$.
Alors $(C', K', \delta|_{K'})$ est un anneau à puissances divisées et
le cardinal de $C'$ est majoré par le cardinal
$\kappa = |A| \otimes |B|^{\aleph_0}$.
De plus, $\varphi$ se factorise sous la forme $A \to C' \to C$ et $\psi$
sous la forme $B \to C' \to C$. Cela démontre que l'hypothèse (1) du
lemme \ref{dpa-lemma-a-version-of-brown} d'Algèbre à puissances divisées
est satisfaite. L'hypothèse (2) est claire,
car le foncteur d'oubli de la catégorie des anneaux à puissances divisées,
$(C, K, \delta) \mapsto (C, K)$, commute aux limites ; voir le
lemme \ref{dpa-lemma-limits} d'Algèbre à puissances divisées et sa démonstration.
\end{proof}

\begin{definition}
\label{definition-divided-power-envelope}
Soit $(A, I, \gamma)$ un anneau à puissances divisées.
Soit $A \to B$ un homomorphisme d'anneaux. Soit $J \subset B$ un idéal
tel que $IB \subset J$. L'algèbre à puissances divisées $(D, \bar J, \bar\gamma)$
construite dans le lemme \ref{lemma-divided-power-envelope}
est appelée l'{\it enveloppe à puissances divisées de $J$ dans $B$
relative à $(A, I, \gamma)$} et se note $D_B(J)$ ou $D_{B, \gamma}(J)$.
\end{definition}

\noindent
Soit $(A, I, \gamma) \to (C, K, \delta)$ un homomorphisme d'anneaux à
puissances divisées. La propriété universelle de
$D_{B, \gamma}(J) = (D, \bar J, \bar \gamma)$ est
$$
\begin{matrix}
\text{homomorphismes d'anneaux }B \to C \\
\text{ qui envoient }J\text{ dans }K
\end{matrix}
\longleftrightarrow
\begin{matrix}
\text{homomorphismes à puissances divisées} \\
(D, \bar J, \bar \gamma) \to (C, K, \delta)
\end{matrix}
$$
et la correspondance s'obtient par précomposition avec l'application $B \to D$
qui correspond à $\text{id}_D$. Voici quelques propriétés de
$(D, \bar J, \bar \gamma)$ qui découlent directement de la propriété
universelle. Il existe des homomorphismes de $A$-algèbres
\begin{equation}
\label{equation-divided-power-envelope}
B \longrightarrow D \longrightarrow B/J
\end{equation}
La première flèche envoie $J$ dans $\bar J$, et $\bar J$ est le noyau
de la seconde. Les éléments $\bar\gamma_n(x)$, où $n > 0$
et où $x$ appartient à l'image de $J \to D$, engendrent $\bar J$
comme idéal de $D$ et engendrent $D$ comme $B$-algèbre.

\begin{lemma}
\label{lemma-divided-power-envelop-quotient}
Soit $(A, I, \gamma)$ un anneau à puissances divisées.
Soit $\varphi : B' \to B$ une surjection de $A$-algèbres de noyau $K$.
Soit $IB \subset J \subset B$ un idéal. Soit $J' \subset B'$
l'image réciproque de $J$. Écrivons
$D_{B', \gamma}(J') = (D', \bar J', \bar\gamma)$.
Alors $D_{B, \gamma}(J) = (D'/K', \bar J'/K', \bar\gamma)$
où $K'$ est l'idéal engendré par les éléments $\bar\gamma_n(k)$
pour $n \geq 1$ et $k \in K$.
\end{lemma}

\begin{proof}
Écrivons $D_{B, \gamma}(J) = (D, \bar J, \bar \gamma)$.
La propriété universelle de $D'$ fournit un homomorphisme $D' \to D$
d'algèbres à puissances divisées. Puisque $B' \to B$ et $J' \to J$ sont surjectifs,
l'homomorphisme $D' \to D$ est surjectif (voir les remarques ci-dessus). Il est clair que
$\bar\gamma_n(k)$ appartient au noyau pour $n \geq 1$ et $k \in K$ ; autrement dit,
nous obtenons un homomorphisme $D'/K' \to D$. Réciproquement, il existe une structure de
puissances divisées sur $\bar J'/K' \subset D'/K'$ ; voir le
lemme \ref{dpa-lemma-kernel} d'Algèbre à puissances divisées.
La propriété universelle de $D$ fournit donc un inverse $D \to D'/K'$, ce qui achève la démonstration.
\end{proof}

\noindent
Dans la situation de la définition \ref{definition-divided-power-envelope},
nous pouvons choisir une surjection $P \to B$, où $P$ est une algèbre
polynomiale sur $A$, et prendre pour $J' \subset P$ l'image réciproque de $J$.
Le lemme précédent décrit $D_{B, \gamma}(J)$ en fonction de
$D_{P, \gamma}(J')$. Notons que $\gamma$ se prolonge en une structure de puissances
divisées $\gamma'$ sur $IP$ d'après le
lemme \ref{dpa-lemma-gamma-extends} d'Algèbre à puissances divisées. Ainsi,
$D_{P, \gamma}(J') = D_{P, \gamma'}(J')$ est un exemple du cas particulier
d'enveloppes à puissances divisées que décrit le lemme suivant.

\begin{lemma}
\label{lemma-describe-divided-power-envelope}
Soit $(B, I, \gamma)$ une algèbre à puissances divisées. Soit $I \subset J \subset B$
un idéal. Soit $(D, \bar J, \bar \gamma)$ l'enveloppe à puissances divisées
de $J$ relativement à $\gamma$. Choisissons des éléments $f_t \in J$, $t \in T$, tels
que $J = I + (f_t)$. Il existe alors une surjection
$$
\Psi :
(B\langle x_t \rangle, IB\langle x_t \rangle + B\langle x_t \rangle_+, \delta)
\longrightarrow
(D, \bar J, \bar \gamma)
$$
d'anneaux à puissances divisées qui envoie $x_t$ sur l'image de $f_t$ dans $D$.
Le noyau de $\Psi$ est engendré par les éléments $x_t - f_t$ et par tous les
éléments
$$
\delta_n\left(\sum r_t x_t - r_0\right)
$$
dès que $\sum r_t f_t = r_0$ dans $B$ pour certains $r_t \in B$, $r_0 \in I$.
\end{lemma}

\begin{proof}
Dans l'énoncé du lemme, nous considérons $B\langle x_t \rangle$
comme un anneau à puissances divisées d'idéal
$J' = IB\langle x_t \rangle + B\langle x_t \rangle_{+}$ ; voir la
remarque \ref{dpa-remark-divided-power-polynomial-algebra} d'Algèbre à puissances divisées.
L'existence de $\Psi$ résulte de la propriété universelle des
algèbres polynomiales à puissances divisées. La surjectivité de $\Psi$ résulte
du fait que son image est un sous-anneau à puissances divisées de $D$, donc égale à $D$
par la propriété universelle de $D$. Il est clair que
$x_t - f_t$ appartient au noyau. Posons
$$
\mathcal{R} = \{(r_0, r_t) \in I \oplus \bigoplus\nolimits_{t \in T} B
\mid \sum r_t f_t = r_0 \text{ dans }B\}
$$
Si $(r_0, r_t) \in \mathcal{R}$, il est clair que
$\sum r_t x_t - r_0$ appartient au noyau.
Puisque $\Psi$ est un homomorphisme d'anneaux à puissances divisées
et que $\sum r_tx_t - r_0 \in J'$,
il s'ensuit que $\delta_n(\sum r_t x_t - r_0)$ appartient aussi au noyau.
Soit $K \subset B\langle x_t \rangle$ l'idéal engendré par
$x_t - f_t$ et les éléments $\delta_n(\sum r_t x_t - r_0)$ pour
$(r_0, r_t) \in \mathcal{R}$.
Pour montrer que $K = \Ker(\Psi)$, il suffit de montrer que
$\delta$ se prolonge à $B\langle x_t \rangle/K$. En effet, la propriété
universelle de $D$ fournit alors une application $D \to B\langle x_t \rangle/K$
inverse de $\Psi$. Nous devons donc montrer que $K \cap J'$ est
stable par $\delta_n$ ; voir le
lemme \ref{dpa-lemma-kernel} d'Algèbre à puissances divisées.
Soit $K' \subset B\langle x_t \rangle$ l'idéal
engendré par les éléments suivants :
\begin{enumerate}
\item $\delta_m(\sum r_t x_t - r_0)$, où $m > 0$ et
$(r_0, r_t) \in \mathcal{R}$,
\item $x_{t'}^{[m]}(x_t - f_t)$, où $m > 0$ et $t', t \in T$.
\end{enumerate}
Nous affirmons que $K' = K \cap J'$. Cette affirmation prouve que $K \cap J'$
est stable par $\delta_n$, $n > 0$, d'après le critère du
lemme \ref{dpa-lemma-kernel} (2)(c) d'Algèbre à puissances divisées
et le calcul de $\delta_n$ sur
les éléments énumérés, que nous laissons au lecteur.
Pour démontrer l'affirmation, notons que $K' \subset K \cap J'$.
Réciproquement, si $h \in K \cap J'$, nous pouvons écrire modulo $K'$
$$
h = \sum r_t (x_t - f_t)
$$
pour certains $r_t \in B$. Puisque $h \in K \cap J' \subset J'$,
nous avons $r_0 = \sum r_t f_t \in I$. Ainsi $(r_0, r_t) \in \mathcal{R}$,
et nous constatons que
$$
h = \sum r_t x_t - r_0
$$
appartient à $K'$, comme voulu.
\end{proof}

\begin{lemma}
\label{lemma-divided-power-envelope-add-variables}
Soit $(A, I, \gamma)$ un anneau à puissances divisées.
Soit $B$ une $A$-algèbre et $IB \subset J \subset B$ un idéal.
Soit $x_i$ une famille d'indéterminées. Alors
$$
D_{B[x_i], \gamma}(JB[x_i] + (x_i)) = D_{B, \gamma}(J) \langle x_i \rangle
$$
\end{lemma}

\begin{proof}
Une démonstration possible consiste à déduire ceci du
lemme \ref{lemma-describe-divided-power-envelope},
car toute relation entre les $x_i$ dans $B[x_i]$ est triviale.
D'autre part, le lemme résulte de la propriété universelle
de l'algèbre polynomiale à puissances divisées et de celle des
enveloppes à puissances divisées.
\end{proof}

\noindent
Les conditions (1) et (2) du lemme suivant sont satisfaites si $B \to B'$ est plat
en tout point de $V(IB') \subset \Spec(B')$ et sont très étroitement liées
à cette condition ; voir
Algèbre, lemme \ref{algebra-lemma-what-does-it-mean}.
Il affirme notamment que la formation de l'enveloppe à
puissances divisées commute à la localisation.

\begin{lemma}
\label{lemma-flat-base-change-divided-power-envelope}
Soit $(A, I, \gamma)$ un anneau à puissances divisées.
Soit $B \to B'$ un homomorphisme de $A$-algèbres.
Supposons que
\begin{enumerate}
\item $B/IB \to B'/IB'$ soit plat, et
\item $\text{Tor}_1^B(B', B/IB) = 0$.
\end{enumerate}
Alors, pour tout idéal $IB \subset J \subset B$, l'application canonique
$$
D_B(J) \otimes_B B' \longrightarrow D_{B'}(JB')
$$
est un isomorphisme.
\end{lemma}

\begin{proof}
Posons $D = D_B(J)$ et notons $\bar J \subset D$ son idéal à puissances divisées,
muni de la structure de puissances divisées $\bar\gamma$. La propriété universelle de
$D$ fournit un homomorphisme de $B$-algèbres $D \to D_{B'}(JB')$, donc l'application du
lemme. Il suffit de montrer que
les puissances divisées $\bar\gamma$ se prolongent à $D \otimes_B B'$, car alors
la propriété universelle de $D_{B'}(JB')$ fournira une application
$D_{B'}(JB') \to D \otimes_B B'$ inverse de celle du lemme.

\medskip\noindent
Choisissons une surjection $P \to B'$, où $P$ est une algèbre polynomiale sur $B$.
En particulier, $B \to P$ est plat, donc $D \to D \otimes_B P$ est plat d'après le
lemme \ref{algebra-lemma-flat-base-change} d'Algèbre.
Alors $\bar\gamma$ se prolonge à $D \otimes_B P$ d'après le
lemme \ref{dpa-lemma-gamma-extends} d'Algèbre à puissances divisées ; nous
noterons encore ce prolongement
$\bar\gamma$. Posons $\mathfrak a = \Ker(P \to B')$, de sorte que
nous avons la suite exacte courte
$$
0 \to \mathfrak a \to P \to B' \to 0
$$
Ainsi, $\text{Tor}_1^B(B', B/IB) = 0$ implique
$\mathfrak a \cap IP = I\mathfrak a$.
Nous avons maintenant le diagramme commutatif suivant :
$$
\xymatrix{
B/J \otimes_B \mathfrak a \ar[r]_\beta &
B/J \otimes_B P \ar[r] &
B/J \otimes_B B' \\
D \otimes_B \mathfrak a \ar[r]^\alpha \ar[u] &
D \otimes_B P \ar[r] \ar[u] &
D \otimes_B B' \ar[u] \\
\bar J \otimes_B \mathfrak a \ar[r] \ar[u] &
\bar J \otimes_B P \ar[r] \ar[u] &
\bar J \otimes_B B' \ar[u]
}
$$
Ce diagramme est exact, même après ajout de $0$ en haut et à droite.
Nous devons montrer que les puissances divisées sur l'idéal
$\bar J \otimes_B P$ préservent l'idéal
$\Im(\alpha) \cap \bar J \otimes_B P$ ; voir le
lemme \ref{dpa-lemma-kernel} d'Algèbre à puissances divisées.
Considérons la suite exacte
$$
0 \to \mathfrak a/I\mathfrak a \to P/IP \to B'/IB' \to 0
$$
(qui utilise $\mathfrak a \cap IP = I\mathfrak a$, comme on l'a vu ci-dessus).
Puisque $B'/IB'$ est plat sur $B/IB$, cette suite reste exacte après
application de $B/J \otimes_{B/IB} -$ ; voir le
lemme \ref{algebra-lemma-flat-tor-zero} d'Algèbre. Ainsi,
$$
\Ker(B/J \otimes_{B/IB} \mathfrak a/I\mathfrak a \to
B/J \otimes_{B/IB} P/IP) =
\Ker(\mathfrak a/J\mathfrak a \to P/JP)
$$
est nul. Ainsi $\beta$ est injectif. Il s'ensuit que
$\Im(\alpha) \cap \bar J \otimes_B P$ coïncide avec
l'image de $\bar J \otimes \mathfrak a$. Si maintenant
$f \in \bar J$ et $a \in \mathfrak a$, alors
$\bar\gamma_n(f \otimes a) = \bar\gamma_n(f) \otimes a^n$
et le résultat est donc clair.
\end{proof}

\noindent
Le lemme suivant est un cas particulier de
\cite[Proposition 2.1.7]{dJ-crystalline}, qui est elle-même une généralisation de
\cite[Proposition 2.8.2]{Berthelot}.

\begin{lemma}
\label{lemma-flat-extension-divided-power-envelope}
Soit $(B, I, \gamma) \to (B', I', \gamma')$ un homomorphisme d'anneaux à
puissances divisées. Soient $I \subset J \subset B$ et
$I' \subset J' \subset B'$ des idéaux. Supposons
\begin{enumerate}
\item $B/I \to B'/I'$ soit plat, et
\item $J' = JB' + I'$.
\end{enumerate}
Alors l'application canonique
$$
D_{B, \gamma}(J) \otimes_B B' \longrightarrow D_{B', \gamma'}(J')
$$
est un isomorphisme.
\end{lemma}

\begin{proof}
Posons $D = D_{B, \gamma}(J)$. Choisissons des éléments $f_t \in J$ engendrant $J/I$.
Posons $\mathcal{R} = \{(r_0, r_t) \in I \oplus \bigoplus\nolimits_{t \in T} B
\mid \sum r_t f_t = r_0 \text{ dans }B\}$ comme dans la démonstration du
lemme \ref{lemma-describe-divided-power-envelope}. Ce lemme montre que
$$
D = B\langle x_t \rangle/ K
$$
où $K$ est engendré par les éléments $x_t - f_t$ et
$\delta_n(\sum r_t x_t - r_0)$ pour $(r_0, r_t) \in \mathcal{R}$.
Nous voyons ainsi que
\begin{equation}
\label{equation-base-change}
D \otimes_B B' = B'\langle x_t \rangle/K'
\end{equation}
où $K'$ est engendré par les images dans $B'\langle x_t \rangle$
des générateurs de $K$ énumérés ci-dessus. Soit $f'_t \in B'$ l'image
de $f_t$. Par l'hypothèse (1), les éléments $f'_t \in J'$
engendrent $J'/I'$, et $x_t - f'_t \in K'$. Posons
$$
\mathcal{R}' =
\{(r'_0, r'_t) \in I' \oplus \bigoplus\nolimits_{t \in T} B'
\mid \sum r'_t f'_t = r'_0 \text{ dans }B'\}
$$
Pour achever la démonstration, nous devons montrer que
$\delta'_n(\sum r'_t x_t - r'_0) \in K'$
pour $(r'_0, r'_t) \in \mathcal{R}'$, car alors la présentation
(\ref{equation-base-change}) de $D \otimes_B B'$ est identique
à la présentation de $D_{B', \gamma'}(J')$ obtenue dans le
lemme \ref{lemma-describe-divided-power-envelope} à partir des générateurs $f'_t$.
Supposons que $(r'_0, r'_t) \in \mathcal{R}'$. Alors
$\sum r'_t f'_t = 0$ dans $B'/I'$. Puisque $B/I \to B'/I'$ est plat par
l'hypothèse (1), nous pouvons appliquer le critère équationnel de platitude
(Algèbre, lemme \ref{algebra-lemma-flat-eq}) et constater
qu'il existe un entier $m > 0$ ainsi que des éléments
$r_{jt} \in B$ et $c_j \in B'$, $j = 1, \ldots, m$, tels
que
$$
r_{j0} = \sum\nolimits_t r_{jt} f_t  \in I \text{ pour } j = 1, \ldots, m
$$
et
$$
i'_t  = r'_t - \sum\nolimits_j c_j r_{jt} \in I' \text{ pour tout }t
$$
Notons que cela implique aussi
$r'_0 = \sum_t i'_t f_t + \sum_j c_j r_{j0}$.
Nous avons alors
\begin{align*}
\delta'_n(\sum\nolimits_t r'_t x_t  - r'_0)
& =
\delta'_n(
\sum\nolimits_t i'_t x_t +
\sum\nolimits_{t, j} c_j r_{jt} x_t -
\sum\nolimits_t i'_t f_t -
\sum\nolimits_j c_j r_{j0}) \\
& =
\delta'_n(
\sum\nolimits_t i'_t(x_t - f_t) +
\sum\nolimits_j c_j (\sum\nolimits_t r_{jt} x_t  - r_{j0}))
\end{align*}
Puisque $\delta_n(a + b) = \sum_{m = 0, \ldots, n} \delta_m(a) \delta_{n - m}(b)$
et que $\delta_m(\sum i'_t(x_t - f_t))$ appartient à l'idéal
engendré par les $x_t - f_t \in K'$ pour $m > 0$, il suffit de prouver que
$\delta_n(\sum c_j (\sum r_{jt} x_t  - r_{j0}))$ appartient à $K'$.
Pour cela, nous utilisons
$$
\delta_n(\sum\nolimits_j c_j (\sum\nolimits_t r_{jt} x_t  - r_{j0}))
=
\sum c_1^{n_1} \ldots c_m^{n_m}
\delta_{n_1}(\sum r_{1t} x_t  - r_{10}) \ldots
\delta_{n_m}(\sum r_{mt} x_t  - r_{m0})
$$
où la somme porte sur $n_1 + \ldots + n_m = n$. Cela démontre l'assertion voulue.
\end{proof}





\section{Quelques épaississements explicites à puissances divisées}
\label{section-explicit-thickenings}

\noindent
Les constructions de cette section nous aideront à définir la connexion
sur un cristal en modules sur le site cristallin.

\begin{lemma}
\label{lemma-divided-power-first-order-thickening}
Soit $(A, I, \gamma)$ un anneau à puissances divisées. Soit $M$ un $A$-module.
Soit $B = A \oplus M$ comme $A$-algèbre, où $M$ est un idéal de carré nul,
et posons $J = I \oplus M$. Posons
$$
\delta_n(x + z) = \gamma_n(x) + \gamma_{n - 1}(x)z
$$
pour $x \in I$ et $z \in M$.
Alors $\delta$ est une structure de puissances divisées et
$A \to B$ est un homomorphisme d'anneaux à puissances divisées de
$(A, I, \gamma)$ vers $(B, J, \delta)$.
\end{lemma}

\begin{proof}
Nous devons vérifier les conditions (1) -- (5) de la
définition \ref{dpa-definition-divided-powers} d'Algèbre à puissances divisées.
Nous le démontrerons directement dans ce cas ; la démonstration du
lemme suivant donne toutefois une méthode qui évite les calculs.
Les conditions (1) et (3) sont claires. La condition (2) résulte de
\begin{align*}
\delta_n(x + z)\delta_m(x + z)
& =
(\gamma_n(x) + \gamma_{n - 1}(x)z)(\gamma_m(x) + \gamma_{m - 1}(x)z) \\
& = \gamma_n(x)\gamma_m(x) + \gamma_n(x)\gamma_{m - 1}(x)z +
\gamma_{n - 1}(x)\gamma_m(x)z \\
& =
\frac{(n + m)!}{n!m!} \gamma_{n + m}(x) +
\left(\frac{(n + m - 1)!}{n!(m - 1)!} +
\frac{(n + m - 1)!}{(n - 1)!m!}\right)
\gamma_{n + m - 1}(x) z \\
& =
\frac{(n + m)!}{n!m!}\delta_{n + m}(x + z)
\end{align*}
La condition (5) résulte de
\begin{align*}
\delta_n(\delta_m(x + z))
& =
\delta_n(\gamma_m(x) + \gamma_{m - 1}(x)z) \\
& =
\gamma_n(\gamma_m(x)) + \gamma_{n - 1}(\gamma_m(x))\gamma_{m - 1}(x)z \\
& =
\frac{(nm)!}{n! (m!)^n} \gamma_{nm}(x) +
\frac{((n - 1)m)!}{(n - 1)! (m!)^{n - 1}}
\gamma_{(n - 1)m}(x) \gamma_{m - 1}(x) z \\
& = \frac{(nm)!}{n! (m!)^n}(\gamma_{nm}(x) + \gamma_{nm - 1}(x) z)
\end{align*}
par un calcul élémentaire d'arithmétique. Pour prouver (4), nous devons voir que
$$
\delta_n(x + x' + z + z')
=
\gamma_n(x + x') + \gamma_{n - 1}(x + x')(z + z')
$$
est égal à
$$
\sum\nolimits_{i = 0}^n
(\gamma_i(x) + \gamma_{i - 1}(x)z)
(\gamma_{n - i}(x') + \gamma_{n - i - 1}(x')z')
$$
Cela résulte aisément du regroupement des coefficients de
$1$, $z$ et $z'$, et de l'emploi de la condition (4) pour $\gamma$.
\end{proof}

\begin{lemma}
\label{lemma-divided-power-second-order-thickening}
Soit $(A, I, \gamma)$ un anneau à puissances divisées. Soient $M$, $N$ des $A$-modules.
Soit $q : M \times M \to N$ une application $A$-bilinéaire.
Munissons $B = A \oplus M \oplus N$ de la structure de $A$-algèbre dont le produit est
$$
(x, z, w)\cdot (x', z', w') = (xx', xz' + x'z, xw' + x'w + q(z, z') + q(z', z))
$$
et posons $J = I \oplus M \oplus N$. Posons
$$
\delta_n(x, z, w) = (\gamma_n(x), \gamma_{n - 1}(x)z,
\gamma_{n - 1}(x)w + \gamma_{n - 2}(x)q(z, z))
$$
pour $(x, z, w) \in J$.
Alors $\delta$ est une structure de puissances divisées et
$A \to B$ est un homomorphisme d'anneaux à puissances divisées de
$(A, I, \gamma)$ vers $(B, J, \delta)$.
\end{lemma}

\begin{proof}
Supposons que nous voulions démontrer la propriété (4) de la
définition \ref{dpa-definition-divided-powers} d'Algèbre à puissances divisées.
Choisissons $(x, z, w)$ et $(x', z', w')$ dans $J$.
Choisissons une application
$$
A_0 = \mathbf{Z}\langle s, s'\rangle \longrightarrow A,\quad
s \longmapsto x,
s' \longmapsto x'
$$
ce qui est possible par la propriété universelle des algèbres
polynomiales à puissances divisées. Posons $M_0 = A_0 \oplus A_0$ et
$N_0 = A_0 \oplus A_0 \oplus M_0 \otimes_{A_0} M_0$.
Soit $q_0 : M_0 \times M_0 \to N_0$ l'application évidente.
Définissons $M_0 \to M$ comme l'application $A_0$-linéaire qui envoie
les vecteurs de base de $M_0$ sur $z$ et $z'$. Définissons $N_0 \to N$
comme l'application $A_0$-linéaire qui envoie les deux premiers vecteurs de base
de $N_0$ sur $w$ et $w'$ et qui utilise
$M_0 \otimes_{A_0} M_0 \to M \otimes_A M \xrightarrow{q} N$
sur le dernier facteur direct. Il suffit alors de démontrer
l'identité (4) dans la situation $(A_0, M_0, N_0, q_0)$.
On procède de même pour les autres identités. Nous sommes ainsi ramenés au cas
d'un anneau $A$ sans $\mathbf{Z}$-torsion et de modules sans torsion sur $A$.
Dans ce cas, il suffit de montrer que
$$
n! \delta_n(x, z, w) = (x, z, w)^n
$$
dans l'anneau $A$ ; voir Algèbre à puissances divisées,
lemme \ref{dpa-lemma-silly}. Pour le voir, notons que
$$
(x, z, w)^2 = (x^2, 2xz, 2xw + 2q(z, z))
$$
et, par récurrence,
$$
(x, z, w)^n = (x^n, nx^{n - 1}z, nx^{n - 1}w + n(n - 1)x^{n - 2}q(z, z))
$$
D'autre part,
$$
n! \delta_n(x, z, w) = (n!\gamma_n(x), n!\gamma_{n - 1}(x)z,
n!\gamma_{n - 1}(x)w + n!\gamma_{n - 2}(x) q(z, z))
$$
ce qui coïncide avec l'expression précédente. La démonstration est achevée.
\end{proof}




\section{Compatibilité}
\label{section-compatibility}

\noindent
La lecture de cette section n'est pas indispensable ; elle explique comment notre étude
s'articule avec celle de \cite{Berthelot}.
Considérons la notion technique suivante.

\begin{definition}
\label{definition-compatible}
Soient $(A, I, \gamma)$ et $(B, J, \delta)$ des anneaux à puissances divisées.
Soit $A \to B$ un homomorphisme d'anneaux. Nous disons que
{\it $\delta$ est compatible avec $\gamma$}
s'il existe une structure de puissances divisées $\bar\gamma$ sur
$J + IB$ telle que
$$
(A, I, \gamma) \to (B, J + IB, \bar \gamma)\quad\text{et}\quad
(B, J, \delta) \to (B, J + IB, \bar \gamma)
$$
soient des homomorphismes d'anneaux à puissances divisées.
\end{definition}

\noindent
Soit $p$ un nombre premier. Soit $(A, I, \gamma)$ un anneau à puissances divisées.
Soit $A \to C$ un homomorphisme d'anneaux tel que $p$ soit nilpotent dans $C$.
Supposons que $\gamma$ se prolonge à $IC$ (voir le
lemme \ref{dpa-lemma-gamma-extends} d'Algèbre à puissances divisées).
Dans cette situation, le (gros) site cristallin affine de
$\Spec(C)$ sur $\Spec(A)$,
tel qu'il est défini dans \cite{Berthelot},
est l'opposé de la catégorie des systèmes
$$
(B, J, \delta, A \to B, C \to B/J)
$$
où
\begin{enumerate}
\item $(B, J, \delta)$ est un anneau à puissances divisées tel que $p$ soit nilpotent dans $B$,
\item $\delta$ est compatible avec $\gamma$, et
\item le diagramme
$$
\xymatrix{
B \ar[r] & B/J \\
A \ar[u] \ar[r] & C \ar[u]
}
$$
est commutatif.
\end{enumerate}
Les conditions
« $\gamma$ se prolonge à $C$ et $\delta$ est compatible avec $\gamma$ »
sont utilisées dans \cite{Berthelot} afin d'assurer que
la cohomologie cristalline de $\Spec(C)$ est la même que celle de
$\Spec(C/IC)$. Nous éviterons cette difficulté
en travaillant exclusivement avec des $C$ tels que $IC = 0$\footnote{Il y aura bien sûr
un prix à payer.}. Dans ce cas,
pour un système $(B, J, \delta, A \to B, C \to B/J)$ comme ci-dessus,
la commutativité du diagramme affiché implique $IB \subset J$, et
la compatibilité équivaut à la condition que
$(A, I, \gamma) \to (B, J, \delta)$ soit un homomorphisme d'anneaux à
puissances divisées.




\section{Site cristallin affine}
\label{section-affine-site}

\noindent
Dans cette section, nous étudions la variante algébrique du site cristallin.
La situation de base dans laquelle nous mènerons cette étude sera la
suivante.

\begin{situation}
\label{situation-affine}
Ici, $p$ est un nombre premier, $(A, I, \gamma)$ est un anneau à puissances
divisées tel que $A$ soit une $\mathbf{Z}_{(p)}$-algèbre, et $A \to C$ est un
homomorphisme d'anneaux tel que $IC = 0$ et que $p$ soit nilpotent dans $C$.
\end{situation}

\noindent
En général, le nombre premier $p$ appartiendra à l'idéal
à puissances divisées $I$.

\begin{definition}
\label{definition-affine-thickening}
Dans la situation \ref{situation-affine} :
\begin{enumerate}
\item Un {\it épaississement à puissances divisées} de $C$ sur $(A, I, \gamma)$
est un homomorphisme d'algèbres à puissances divisées $(A, I, \gamma) \to (B, J, \delta)$
tel que $p$ soit nilpotent dans $B$, muni d'un homomorphisme d'anneaux $C \to B/J$ tel que
$$
\xymatrix{
B \ar[r] & B/J \\
& C \ar[u] \\
A \ar[uu] \ar[r] & A/I \ar[u]
}
$$
soit commutatif.
\item Un {\it homomorphisme d'épaississements à puissances divisées}
$$
(B, J, \delta, C \to B/J) \longrightarrow (B', J', \delta', C \to B'/J')
$$
est un homomorphisme $\varphi : B \to B'$ de $A$-algèbres à puissances divisées tel
que $C \to B/J \to B'/J'$ soit l'application donnée $C \to B'/J'$.
\item Nous notons $\text{CRIS}(C/A, I, \gamma)$, ou simplement $\text{CRIS}(C/A)$,
la catégorie des épaississements à puissances divisées de $C$ sur $(A, I, \gamma)$.
\item Nous notons $\text{Cris}(C/A, I, \gamma)$, ou simplement $\text{Cris}(C/A)$,
la sous-catégorie pleine formée des $(B, J, \delta, C \to B/J)$ tels que
$C \to B/J$ soit un isomorphisme. Nous notons souvent un tel objet
$(B \to C, \delta)$, en sous-entendant $J = \Ker(B \to C)$.
\end{enumerate}
\end{definition}

\noindent
Notons que, pour un épaississement à puissances divisées $(B, J, \delta)$ comme ci-dessus,
l'idéal $J$ est localement nilpotent ; voir le
lemme \ref{dpa-lemma-nil} d'Algèbre à puissances divisées.
Il existe un foncteur canonique
\begin{equation}
\label{equation-forget-affine}
\text{CRIS}(C/A) \longrightarrow C\text{-algèbres},\quad
(B, J, \delta) \longmapsto B/J
\end{equation}
Cette catégorie n'admet en général ni égalisateurs ni produits fibrés.
Elle n'admet pas non plus, en général, d'objet initial ($=$ colimite vide).

\begin{lemma}
\label{lemma-affine-thickenings-colimits}
Dans la situation \ref{situation-affine}.
\begin{enumerate}
\item $\text{CRIS}(C/A)$ admet les produits finis (mais non les produits infinis),
\item $\text{CRIS}(C/A)$ admet toutes les colimites finies non vides et
(\ref{equation-forget-affine}) commute à celles-ci, et
\item $\text{Cris}(C/A)$ admet toutes les colimites finies non vides et
$\text{Cris}(C/A) \to \text{CRIS}(C/A)$ commute à celles-ci.
\end{enumerate}
\end{lemma}

\begin{proof}
Le produit vide, c'est-à-dire l'objet final de la catégorie des épaississements à
puissances divisées de $C$ sur $(A, I, \gamma)$, est l'anneau nul considéré
comme une $A$-algèbre, muni de l'idéal nul et de l'unique structure de puissances divisées
sur cet idéal, puis de l'unique homomorphisme de $C$ vers
l'anneau nul. Si $(B_t, J_t, \delta_t)_{t \in T}$ est une famille d'objets de
$\text{CRIS}(C/A)$, nous pouvons former le produit
$(\prod_t B_t, \prod_t J_t, \prod_t \delta_t)$ comme dans le
lemme \ref{dpa-lemma-limits} d'Algèbre à puissances divisées.
L'application $C \to \prod B_t/\prod J_t = \prod B_t/J_t$ est évidente.
Toutefois, $p$ n'est nécessairement nilpotent dans $\prod_t B_t$
que si $T$ est fini.

\medskip\noindent
Étant donnés deux objets $(B, J, \gamma)$ et $(B', J', \gamma')$ de
$\text{CRIS}(C/A)$, nous pouvons former un diagramme cocartésien
$$
\xymatrix{
(B, J, \delta) \ar[r] & (B'', J'', \delta'') \\
(A, I, \gamma) \ar[r] \ar[u] & (B', J', \delta') \ar[u]
}
$$
dans la catégorie des anneaux à puissances divisées. Nous avons alors
$$
B''/J'' = B/J \otimes_{A/I} B'/J' \longleftarrow C \otimes_{A/I} C
$$
voir Algèbre à puissances divisées, remarque \ref{dpa-remark-pushout}.
Notons $J'' \subset K \subset B''$
l'idéal tel que
$$
\xymatrix{
B''/J'' \ar[r] & B''/K \\
C \otimes_{A/I} C \ar[r] \ar[u] & C \ar[u]
}
$$
soit cocartésien, c'est-à-dire $B''/K \cong B/J \otimes_C B'/J'$.
Soit $D_{B''}(K) = (D, \bar K, \bar \delta)$
l'enveloppe à puissances divisées de $K$ dans $B''$ relativement à
$(B'', J'', \delta'')$. On vérifie alors aisément que
$(D, \bar K, \bar \delta)$ est un coproduit de $(B, J, \delta)$ et
$(B', J', \delta')$ dans $\text{CRIS}(C/A)$.

\medskip\noindent
Passons maintenant aux coégalisateurs. Soient
$\alpha, \beta : (B, J, \delta) \to (B', J', \delta')$ des morphismes de
$\text{CRIS}(C/A)$. Considérons $B'' = B'/ (\alpha(b) - \beta(b))$. Soit
$J'' \subset B''$ l'image de $J'$. Soit
$D_{B''}(J'') = (D, \bar J, \bar\delta)$ l'enveloppe à puissances divisées de
$J''$ dans $B''$ relativement à $(B', J', \delta')$. On vérifie alors aisément
que $(D, \bar J, \bar \delta)$ est le coégalisateur de $(B, J, \delta)$ et
$(B', J', \delta')$ dans $\text{CRIS}(C/A)$.

\medskip\noindent
D'après Catégories, lemme \ref{categories-lemma-almost-finite-colimits-exist},
$\text{CRIS}(C/A)$ admet toutes les colimites finies non vides. Les constructions
ci-dessus montrent que (\ref{equation-forget-affine}) commute à celles-ci.
Cela implique formellement (3), puisque $\text{Cris}(C/A)$ est la catégorie-fibre
de (\ref{equation-forget-affine}) au-dessus de $C$.
\end{proof}

\begin{remark}
\label{remark-completed-affine-site}
Dans la situation \ref{situation-affine}, nous notons
$\text{Cris}^\wedge(C/A)$ la catégorie dont les objets sont
les couples $(B \to C, \delta)$ tels que
\begin{enumerate}
\item $B$ est une $A$-algèbre $p$-adiquement complète,
\item $B \to C$ est une surjection de $A$-algèbres,
\item $\delta$ est une structure de puissances divisées sur $\Ker(B \to C)$,
\item $A \to B$ est un homomorphisme d'anneaux à puissances divisées.
\end{enumerate}
Les morphismes sont définis comme dans la définition \ref{definition-affine-thickening}.
Alors $\text{Cris}(C/A) \subset \text{Cris}^\wedge(C/A)$ est la sous-catégorie
pleine formée des $B$ tels que $p$ soit nilpotent dans $B$.
Réciproquement, tout objet $(B \to C, \delta)$ de $\text{Cris}^\wedge(C/A)$
est égal à la limite
$$
(B \to C, \delta) = \lim_e (B/p^eB \to C, \delta)
$$
où, pour $e \gg 0$, l'objet $(B/p^eB \to C, \delta)$ appartient
à $\text{Cris}(C/A)$ ; voir le
lemme \ref{dpa-lemma-extend-to-completion} d'Algèbre à puissances divisées.
En particulier, $\text{Cris}^\wedge(C/A)$ est une sous-catégorie pleine
de la catégorie des pro-objets de $\text{Cris}(C/A)$ ; voir
Catégories, remarque \ref{categories-remark-pro-category}.
\end{remark}

\begin{lemma}
\label{lemma-list-properties}
Dans la situation \ref{situation-affine}.
Soit $P \to C$ une surjection de $A$-algèbres de noyau $J$.
Écrivons $D_{P, \gamma}(J) = (D, \bar J, \bar\gamma)$.
Soit $(D^\wedge, J^\wedge, \bar\gamma^\wedge)$ le
complété $p$-adique de $D$ ; voir le
lemme \ref{dpa-lemma-extend-to-completion} d'Algèbre à puissances divisées.
Pour tout $e \geq 1$, posons $P_e = P/p^eP$ et désignons par $J_e \subset P_e$
l'image de $J$, puis écrivons
$D_{P_e, \gamma}(J_e) = (D_e, \bar J_e, \bar\gamma)$.
Alors, pour tout $e$ assez grand, nous avons :
\begin{enumerate}
\item $p^eD \subset \bar J$ et $p^eD^\wedge \subset \bar J^\wedge$
sont stables par les puissances divisées,
\item $D^\wedge/p^eD^\wedge = D/p^eD = D_e$ comme anneaux à puissances divisées,
\item $(D_e, \bar J_e, \bar\gamma)$ est un objet de $\text{Cris}(C/A)$,
\item $(D^\wedge, \bar J^\wedge, \bar\gamma^\wedge)$ est égal à
$\lim_e (D_e, \bar J_e, \bar\gamma)$, et
\item $(D^\wedge, \bar J^\wedge, \bar\gamma^\wedge)$ est un objet de
$\text{Cris}^\wedge(C/A)$.
\end{enumerate}
\end{lemma}

\begin{proof}
L'assertion (1) résulte du
lemme \ref{dpa-lemma-extend-to-completion} d'Algèbre à puissances divisées.
Une propriété générale de la complétion $p$-adique donne
$D/p^eD = D^\wedge/p^eD^\wedge$. Puisque $D/p^eD$ est un anneau à puissances divisées
et que $P \to D/p^eD$ se factorise par $P_e$, la propriété universelle de
$D_e$ fournit une application $D_e \to D/p^eD$. Réciproquement, la propriété universelle
de $D$ fournit une application $D \to D_e$ qui se factorise par $D/p^eD$. Nous omettons
la vérification que ces applications sont inverses l'une de l'autre. Cela prouve (2).
Si $e$ est assez grand, alors $p^eC = 0$, d'où (3).
L'assertion (4) résulte du
lemme \ref{dpa-lemma-extend-to-completion} d'Algèbre à puissances divisées.
L'assertion (5) est claire d'après les définitions.
\end{proof}

\begin{lemma}
\label{lemma-set-generators}
Dans la situation \ref{situation-affine}.
Soit $P$ une algèbre polynomiale sur $A$ et soit
$P \to C$ une surjection de $A$-algèbres de noyau $J$.
Avec $(D_e, \bar J_e, \bar\gamma)$ comme dans le lemme \ref{lemma-list-properties},
pour tout objet $(B, J_B, \delta)$ de $\text{CRIS}(C/A)$, il
existe un $e$ et un morphisme $D_e \to B$ dans $\text{CRIS}(C/A)$.
\end{lemma}

\begin{proof}
Nous pouvons trouver un homomorphisme de $A$-algèbres $P \to B$
relevant l'application $C \to B/J_B$. D'après notre définition de
$\text{CRIS}(C/A)$, nous avons $p^eB = 0$ pour
un certain $e$ ; ainsi $P \to B$ se factorise sous la forme $P \to P_e \to B$.
La propriété universelle de l'enveloppe à puissances divisées montre alors
que $P_e \to B$ se factorise par $D_e$.
\end{proof}

\begin{lemma}
\label{lemma-generator-completion}
Dans la situation \ref{situation-affine}.
Soit $P$ une algèbre polynomiale sur $A$ et soit
$P \to C$ une surjection de $A$-algèbres de noyau $J$.
Soit $(D, \bar J, \bar\gamma)$ le complété $p$-adique de
$D_{P, \gamma}(J)$. Pour tout objet $(B \to C, \delta)$ de
$\text{Cris}^\wedge(C/A)$, il
existe un morphisme $D \to B$ dans $\text{Cris}^\wedge(C/A)$.
\end{lemma}

\begin{proof}
Nous pouvons trouver un homomorphisme de $A$-algèbres $P \to B$ compatible
avec les applications vers $C$. D'après notre définition de
$\text{Cris}^\wedge(C/A)$, l'application $P \to B$ se factorise sous la forme
$P \to D_{P, \gamma}(J) \to B$. Puisque $B$ est $p$-adiquement complet,
nous pouvons factoriser cette application par $D$.
\end{proof}



\section{Module des différentielles}
\label{section-differentials}

\noindent
Dans cette section, nous développons une théorie des modules de différentielles
pour les anneaux à puissances divisées.

\begin{definition}
\label{definition-derivation}
Soit $A$ un anneau. Soit $(B, J, \delta)$ un anneau à puissances divisées.
Soit $A \to B$ un homomorphisme d'anneaux. Soit $M$ un $B$-module.
Une {\it $A$-dérivation à puissances divisées} à valeurs dans $M$ est une application
$\theta : B \to M$ additive, qui annule les éléments
de $A$, vérifie la règle de Leibniz
$\theta(bb') = b\theta(b') + b'\theta(b)$ et satisfait
$$
\theta(\delta_n(x)) = \delta_{n - 1}(x)\theta(x)
$$
pour tout $n \geq 1$ et tout $x \in J$.
\end{definition}

\noindent
Dans la situation de la définition, comme pour les dérivations
usuelles, il existe une {\it $A$-dérivation universelle à puissances divisées}
$$
\text{d}_{B/A, \delta} : B \to \Omega_{B/A, \delta}
$$
telle que toute $A$-dérivation à puissances divisées $\theta : B \to M$ s'écrive
$\theta = \xi \circ d_{B/A, \delta}$ pour une unique application $B$-linéaire
$\xi : \Omega_{B/A, \delta} \to M$. Si $(A, I, \gamma) \to (B, J, \delta)$
est un homomorphisme d'anneaux à puissances divisées, nous pouvons oublier les
puissances divisées sur $A$ et considérer les dérivations à puissances divisées de
$B$ sur $A$. Voici quelques propriétés fondamentales du module universel
des différentielles (à puissances divisées).

\begin{lemma}
\label{lemma-omega}
Soit $A$ un anneau. Soit $(B, J, \delta)$ un anneau à puissances divisées et
$A \to B$ un homomorphisme d'anneaux. 
\begin{enumerate}
\item Considérons $B[x]$ muni de l'idéal à puissances divisées $(JB[x], \delta')$,
où $\delta'$ est le prolongement de $\delta$ à $B[x]$. Alors
$$
\Omega_{B[x]/A, \delta'} =
\Omega_{B/A, \delta} \otimes_B B[x] \oplus B[x]\text{d}x.
$$
\item Considérons $B\langle x \rangle$ muni de l'idéal à puissances divisées
$(JB\langle x \rangle + B\langle x \rangle_{+}, \delta')$. Alors
$$
\Omega_{B\langle x\rangle/A, \delta'} =
\Omega_{B/A, \delta} \otimes_B B\langle x \rangle \oplus
B\langle x\rangle \text{d}x.
$$
\item Soit $K \subset J$ un idéal stable par $\delta_n$ pour
tout $n > 0$. Posons $B' = B/K$ et notons $\delta'$ la structure de
puissances divisées induite sur $J/K$. Alors $\Omega_{B'/A, \delta'}$ est le quotient
de $\Omega_{B/A, \delta} \otimes_B B'$ par le sous-$B'$-module engendré
par les $\text{d}k$, pour $k \in K$.
\end{enumerate}
\end{lemma}

\begin{proof}
Ces assertions se déduisent directement de la construction de $\Omega_{B/A, \delta}$
comme le $B$-module libre sur les éléments $\text{d}b$, modulo les relations
\begin{enumerate}
\item $\text{d}(b + b') = \text{d}b + \text{d}b'$, $b, b' \in B$,
\item $\text{d}a = 0$, $a \in A$,
\item $\text{d}(bb') = b \text{d}b' + b' \text{d}b$, $b, b' \in B$,
\item $\text{d}\delta_n(f) = \delta_{n - 1}(f)\text{d}f$, $f \in J$, $n > 1$.
\end{enumerate}
Notons que la dernière relation explique pourquoi nous obtenons « la même » réponse pour
l'algèbre polynomiale à puissances divisées et l'algèbre polynomiale usuelle :
dans le premier cas, $x$ appartient à l'idéal à puissances divisées, donc
$\text{d}x^{[n]} = x^{[n - 1]}\text{d}x$.
\end{proof}

\noindent
Soit $(A, I, \gamma)$ un anneau à puissances divisées. Dans ce cadre, la
bonne notion de puissances de $I$ est donnée par les puissances divisées
$$
I^{[n]} = \text{idéal engendré par }
\gamma_{e_1}(x_1) \ldots \gamma_{e_t}(x_t)
\text{ avec }\sum e_j \geq n\text{ et }x_j \in I.
$$
Bien entendu, $I^n \subset I^{[n]}$. Notons que $I^{[1]} = I$.
Nous posons parfois aussi $I^{[0]} = A$.

\begin{lemma}
\label{lemma-diagonal-and-differentials}
Soit $(A, I, \gamma) \to (B, J, \delta)$ un homomorphisme
d'anneaux à puissances divisées. Soit $(B(1), J(1), \delta(1))$ le coproduit
de $(B, J, \delta)$ avec lui-même sur $(A, I, \gamma)$, c'est-à-dire
tel que
$$
\xymatrix{
(B, J, \delta) \ar[r] & (B(1), J(1), \delta(1)) \\
(A, I, \gamma) \ar[r] \ar[u] & (B, J, \delta) \ar[u]
}
$$
est cocartésien. Notons $K = \Ker(B(1) \to B)$.
Alors $K \cap J(1) \subset J(1)$ est stable par la structure de
puissances divisées et
$$
\Omega_{B/A, \delta} = K/ \left(K^2 + (K \cap J(1))^{[2]}\right)
$$
canoniquement.
\end{lemma}

\begin{proof}
Le fait que $K \cap J(1) \subset J(1)$ soit stable par la structure de
puissances divisées résulte de ce que $B(1) \to B$ est un homomorphisme
d'anneaux à puissances divisées.

\medskip\noindent
Rappelons que $K/K^2$ possède une structure canonique de $B$-module.
Notons $s_0, s_1 : B \to B(1)$ les deux coprojections et considérons
l'application $\text{d} : B \to K/(K^2 +(K \cap J(1))^{[2]})$ donnée par
$b \mapsto s_1(b) - s_0(b)$. Il est clair que $\text{d}$ est additive,
s'annule sur $A$ et satisfait à la règle de Leibniz.
Nous affirmons que $\text{d}$ est une $A$-dérivation à puissances divisées.
Soit $x \in J$. Posons $y = s_1(x)$ et $z = s_0(x)$.
Nous écrirons $\delta$ au lieu de $\delta(1)$ pour
la structure de puissances divisées sur $J(1)$.
Il faut montrer que $\delta_n(y) - \delta_n(z) = \delta_{n - 1}(y)(y - z)$
modulo $K^2 +(K \cap J(1))^{[2]}$ pour $n \geq 1$.
L'égalité vaut pour $n = 1$. Supposons $n > 1$.
Notons que $\delta_i(y - z)$ appartient à $(K \cap J(1))^{[2]}$ si $i > 1$.
En calculant modulo $K^2 + (K \cap J(1))^{[2]}$, nous obtenons
$$
\delta_n(z) = \delta_n(z - y + y) =
\sum\nolimits_{i = 0}^n \delta_i(z - y)\delta_{n - i}(y) =
\delta_{n - 1}(y) \delta_1(z - y) + \delta_n(y)
$$
Cela démontre l'égalité voulue.

\medskip\noindent
Soit $M$ un $B$-module. Soit $\theta : B \to M$ une
$A$-dérivation à puissances divisées.
Posons $D = B \oplus M$, où $M$ est un idéal de carré nul. Définissons une
structure de puissances divisées sur $J \oplus M \subset D$ en posant
$\delta_n(x + m) = \delta_n(x) + \delta_{n - 1}(x)m$ pour $n > 1$ ; voir le
lemme \ref{lemma-divided-power-first-order-thickening}.
Il existe deux homomorphismes d'algèbres à puissances divisées $B \to D$ : le premier
est l'inclusion et le second est l'application $b \mapsto b + \theta(b)$.
On obtient donc un homomorphisme canonique $B(1) \to D$ d'algèbres à
puissances divisées sur $(A, I, \gamma)$. Il induit une application $K \to M$
qui s'annule sur $K^2$ (car $M$ est un idéal de carré nul) ainsi que sur
$(K \cap J(1))^{[2]}$, puisque $M^{[2]} = 0$. Par construction, le composé
$B \to K/K^2 + (K \cap J(1))^{[2]} \to M$ est égal à $\theta$.
Il s'ensuit que $\text{d}$
est une $A$-dérivation universelle à puissances divisées, ce qui achève la démonstration.
\end{proof}

\begin{remark}
\label{remark-filtration-differentials}
Soit $A \to B$ un homomorphisme d'anneaux et soit $(J, \delta)$ une structure de
puissances divisées sur $B$. Le module universel $\Omega_{B/A, \delta}$
est muni d'une structure supplémentaire : le $B$-sous-module
$N$ de $\Omega_{B/A, \delta}$ engendré par $\text{d}_{B/A, \delta}(J)$.
Au moyen de l'isomorphisme donné dans le
lemme \ref{lemma-diagonal-and-differentials},
celui-ci correspond à l'image de
$K \cap J(1)$ dans $\Omega_{B/A, \delta}$. Considérons l'$A$-algèbre
$D = B \oplus \Omega^1_{B/A, \delta}$, munie de l'idéal $\bar J = J \oplus N$
et des puissances divisées $\bar \delta$ définies comme dans la démonstration du lemme.
Alors $(D, \bar J, \bar \delta)$ est un anneau à puissances divisées
et les deux applications $B \to D$ données par $b \mapsto b$ et
$b \mapsto b + \text{d}_{B/A, \delta}(b)$
sont des homomorphismes d'anneaux à puissances divisées sur $A$. De plus, $N$
est le plus petit sous-module de $\Omega_{B/A, \delta}$ pour lequel il en soit ainsi.
\end{remark}

\begin{lemma}
\label{lemma-diagonal-and-differentials-affine-site}
Dans la situation \ref{situation-affine}.
Soit $(B, J, \delta)$ un objet de $\text{CRIS}(C/A)$.
Soit $(B(1), J(1), \delta(1))$ le coproduit de $(B, J, \delta)$
avec lui-même dans $\text{CRIS}(C/A)$. Notons
$K = \Ker(B(1) \to B)$. Alors $K \cap J(1) \subset J(1)$
est stable par la structure de puissances divisées et
$$
\Omega_{B/A, \delta} = K/ \left(K^2 + (K \cap J(1))^{[2]}\right)
$$
canoniquement.
\end{lemma}

\begin{proof}
La démonstration est mot pour mot celle du
lemme \ref{lemma-diagonal-and-differentials}.
Le seul point à vérifier est que l'anneau à
puissances divisées $D = B \oplus M$ est un objet de $\text{CRIS}(C/A)$
et que les deux applications $B \to D$ sont des morphismes de $\text{CRIS}(C/A)$.
Comme $D/(J \oplus M) = B/J$, l'application $C \to B/J$ permet de considérer
$D$ comme un objet de $\text{CRIS}(C/A)$,
et l'assertion concernant les morphismes résulte immédiatement de la construction.
\end{proof}

\begin{lemma}
\label{lemma-module-differentials-divided-power-envelope}
Soit $(A, I, \gamma)$ un anneau à puissances divisées. Soit $A \to B$ un homomorphisme
d'anneaux et soit $IB \subset J \subset B$ un idéal. Soit
$D_{B, \gamma}(J) = (D, \bar J, \bar \gamma)$ l'enveloppe à puissances divisées.
Alors on a
$$
\Omega_{D/A, \bar\gamma} = \Omega_{B/A} \otimes_B D
$$
\end{lemma}

\begin{proof}[Première démonstration]
Soit $M$ un $D$-module. Nous affirmons que la donnée d'une $A$-dérivation
$\vartheta : B \to M$ équivaut à celle d'une
$A$-dérivation à puissances divisées $\theta : D \to M$. Le lemme de Yoneda
montre que cette assertion implique l'énoncé.

\medskip\noindent
Considérons l'épaississement de carré nul $D \oplus M$ de $D$.
Il existe une structure de puissances divisées $\delta$ sur $\bar J \oplus M$
si l'on annule sur $M$ les opérations supérieures de puissances divisées.
Autrement dit, on pose
$\delta_n(x + m) = \bar\gamma_n(x) + \bar\gamma_{n - 1}(x)m$ pour
tout $x \in \bar J$ et tout $m \in M$ ; voir le
lemme \ref{lemma-divided-power-first-order-thickening}.
Considérons l'homomorphisme de $A$-algèbres $B \to D \oplus M$ dont la première
composante est l'application $B \to D$ et dont la seconde composante
est $\vartheta$. La propriété universelle fournit un
homomorphisme correspondant $D \to D \oplus M$ d'algèbres à puissances divisées,
dont la seconde composante est la
$A$-dérivation à puissances divisées $\theta$ correspondant à $\vartheta$.
\end{proof}

\begin{proof}[Deuxième démonstration]
Démontrons d'abord l'assertion lorsque $B$ est plat sur $A$. Dans ce cas, $\gamma$
se prolonge en une structure de puissances divisées $\gamma'$ sur $IB$ ; voir
Algèbre à puissances divisées, lemme \ref{dpa-lemma-gamma-extends}.
Ainsi $D = D_{B, \gamma'}(J)$ est un quotient de
l'anneau à puissances divisées $(D', J', \delta)$, où $D' =  B\langle x_t \rangle$
et $J' = IB\langle x_t \rangle + B\langle x_t \rangle_{+}$,
par les éléments $x_t - f_t$ et $\delta_n(\sum r_t x_t - r_0)$ ; voir le
lemme \ref{lemma-describe-divided-power-envelope} pour les notations
et les explications. Notons $\text{d} : D' \to \Omega_{D'/A, \delta}$
la dérivation universelle. On a
$$
\Omega_{D'/A, \delta} =
\Omega_{B/A} \otimes_B D' \oplus \bigoplus D' \text{d}x_t,
$$
voir le lemme \ref{lemma-omega}. On en déduit que $\Omega_{D/A, \bar\gamma}$
est le quotient de $\Omega_{D'/A, \delta} \otimes_{D'} D$ par le sous-module
engendré par les images sous $\text{d}$ des générateurs du
noyau de $D' \to D$ énumérés ci-dessus ; voir le lemme \ref{lemma-omega}.
Comme $\text{d}(x_t - f_t) = - \text{d}f_t + \text{d}x_t$,
on a $\text{d}x_t = \text{d}f_t$ dans le quotient.
En particulier, $\Omega_{B/A} \otimes_B D \to \Omega_{D/A, \gamma}$
est surjectif, de noyau engendré par les images sous $\text{d}$
des éléments $\delta_n(\sum r_t x_t - r_0)$.
Or, étant donnée une relation $\sum r_tf_t - r_0 = 0$ dans $B$, avec
$r_t \in B$ et $r_0 \in IB$, on a
\begin{align*}
\text{d}\delta_n(\sum r_t x_t - r_0)
& =
\delta_{n - 1}(\sum r_t x_t - r_0)\text{d}(\sum r_t x_t - r_0)
\\
& =
\delta_{n - 1}(\sum r_t x_t - r_0)
\left(
\sum r_t\text{d}(x_t - f_t) + \sum (x_t - f_t)\text{d}r_t
\right)
\end{align*}
car $\sum r_tf_t - r_0 = 0$ dans $B$. Cette expression est donc déjà nulle dans
$\Omega_{B/A} \otimes_A D$, ce qui achève la démonstration lorsque $B$ est plat sur $A$.

\medskip\noindent
Dans le cas général, écrivons $B$ comme quotient d'un anneau de polynômes
$P \to B$ et soit $J' \subset P$ l'image réciproque de $J$. Alors
$D = D'/K'$, avec les notations du
lemme \ref{lemma-divided-power-envelop-quotient}.
D'après le cas traité au premier paragraphe de la démonstration, on a
$\Omega_{D'/A, \bar\gamma'} = \Omega_{P/A} \otimes_P D'$. Alors
$\Omega_{D/A, \bar \gamma}$ est le quotient de $\Omega_{P/A} \otimes_P D$
par le sous-module engendré par les $\text{d}\bar\gamma_n'(k)$, où $k$
parcourt le noyau de $P \to B$ ; voir le
lemme \ref{lemma-omega} et la description de $K'$ donnée au
lemme \ref{lemma-divided-power-envelop-quotient}. Comme
$\text{d}\bar\gamma_n'(k) = \bar\gamma'_{n - 1}(k)\text{d}k$, on voit
qu'il suffit encore de quotienter par le sous-module engendré par les
$\text{d}k$ avec $k \in \Ker(P \to B)$ ; puisque $\Omega_{B/A}$
est le quotient de $\Omega_{P/A} \otimes_A B$ par ces éléments
(Algèbre, lemme \ref{algebra-lemma-differential-seq}), ce qui achève la démonstration.
\end{proof}

\begin{remark}
\label{remark-divided-powers-de-rham-complex}
Soit $A \to B$ un homomorphisme d'anneaux et soit $(J, \delta)$ une structure de
puissances divisées sur $B$. Posons
$\Omega_{B/A, \delta}^i = \wedge^i_B \Omega_{B/A, \delta}$
où $\Omega_{B/A, \delta}$ est le but de la dérivation universelle à puissances divisées
sur $A$, $\text{d} = \text{d}_{B/A} : B \to \Omega_{B/A, \delta}$.
Notons que $\Omega_{B/A, \delta}$ est le quotient de $\Omega_{B/A}$ par le
$B$-sous-module engendré par les éléments
$\text{d}\delta_n(x) - \delta_{n - 1}(x)\text{d}x$ pour $x \in J$.
Nous affirmons que le lemme \ref{algebra-lemma-de-rham-complex} d'Algèbre s'applique.
Pour le voir, il suffit de vérifier que les éléments
$\text{d}\delta_n(x) - \delta_{n - 1}(x)\text{d}x$
de $\Omega_B$ ont une image nulle dans $\Omega^2_{B/A, \delta}$.
On observe que
$$
\text{d}(\delta_{n - 1}(x)) \wedge \text{d}x
= \delta_{n - 2}(x) \text{d}x \wedge \text{d}x = 0
$$
dans $\Omega^2_{B/A, \delta}$, ce qui donne l'annulation voulue. On obtient donc un
{\it complexe de de Rham à puissances divisées}
$$
\Omega^0_{B/A, \delta} \to \Omega^1_{B/A, \delta} \to
\Omega^2_{B/A, \delta} \to \ldots
$$
qui jouera un rôle important dans la suite.
\end{remark}

\begin{remark}
\label{remark-connection}
Soit $A \to B$ un homomorphisme d'anneaux. Soit $\Omega_{B/A} \to \Omega$
un quotient satisfaisant aux hypothèses du
lemme \ref{algebra-lemma-de-rham-complex} d'Algèbre.
Soit $M$ un $B$-module. Une {\it connexion} est une application additive
$$
\nabla : M \longrightarrow M \otimes_B \Omega
$$
telle que $\nabla(bm) = b \nabla(m) + m \otimes \text{d}b$
pour $b \in B$ et $m \in M$. Dans cette situation, on peut définir des applications
$$
\nabla : M \otimes_B \Omega^i \longrightarrow M \otimes_B \Omega^{i + 1}
$$
par la règle $\nabla(m \otimes \omega) = \nabla(m) \wedge \omega +
m \otimes \text{d}\omega$. Celle-ci est bien définie car, si $b \in B$, alors
\begin{align*}
\nabla(bm \otimes \omega) - \nabla(m \otimes b\omega)
& =
\nabla(bm) \wedge \omega + bm \otimes \text{d}\omega
- \nabla(m) \wedge b\omega - m \otimes \text{d}(b\omega) \\
& =
b\nabla(m) \wedge \omega + m \otimes \text{d}b \wedge \omega
+ bm \otimes \text{d}\omega \\
& \ \ \ \ \ \ - b\nabla(m) \wedge \omega - bm \otimes \text{d}(\omega)
- m \otimes \text{d}b \wedge \omega = 0
\end{align*}
Comme il est d'usage, on dit que la connexion est {\it intégrable} si et
seulement si le composé
$$
M \xrightarrow{\nabla} M \otimes_B \Omega^1
\xrightarrow{\nabla} M \otimes_B \Omega^2
$$
est nul. Dans ce cas, on obtient un complexe
$$
M \xrightarrow{\nabla} M \otimes_B \Omega^1
\xrightarrow{\nabla} M \otimes_B \Omega^2
\xrightarrow{\nabla} M \otimes_B \Omega^3
\xrightarrow{\nabla} M \otimes_B \Omega^4 \to \ldots
$$
appelé complexe de de Rham de la connexion.
\end{remark}

\begin{remark}
\label{remark-base-change-connection}
Considérons un diagramme commutatif d'anneaux
$$
\xymatrix{
B \ar[r]_\varphi & B' \\
A \ar[u] \ar[r] & A' \ar[u]
}
$$
Soient $\Omega_{B/A} \to \Omega$ et $\Omega_{B'/A'} \to \Omega'$
des quotients satisfaisant aux hypothèses du
lemme \ref{algebra-lemma-de-rham-complex} d'Algèbre.
Supposons donnée une application $\varphi : \Omega \to \Omega'$ qui
s'insère dans un diagramme commutatif
$$
\xymatrix{
\Omega_{B/A} \ar[r] \ar[d] &
\Omega_{B'/A'} \ar[d] \\
\Omega \ar[r]^{\varphi} &
\Omega'
}
$$
où la flèche horizontale supérieure est l'application canonique
$\Omega_{B/A} \to \Omega_{B'/A'}$ induite par $\varphi : B \to B'$.
Dans cette situation, pour toute paire $(M, \nabla)$ où $M$ est un $B$-module
et $\nabla : M \to M \otimes_B \Omega$ une connexion,
on obtient par {\it changement de base} la paire $(M \otimes_B B', \nabla')$, où
$$
\nabla' :
M \otimes_B B'
\longrightarrow
(M \otimes_B B') \otimes_{B'} \Omega' = M \otimes_B \Omega'
$$
est définie par la règle
$$
\nabla'(m \otimes b') =
\sum m_i \otimes b'\text{d}\varphi(b_i) + m \otimes \text{d}b' 
$$
si $\nabla(m) = \sum m_i \otimes \text{d}b_i$. Si $\nabla$ est intégrable,
alors $\nabla'$ l'est aussi et il existe dans ce cas une application canonique de
complexes de de Rham (remarque \ref{remark-connection})
\begin{equation}
\label{equation-base-change-map-complexes}
M \otimes_B \Omega^\bullet
\longrightarrow
(M \otimes_B B') \otimes_{B'} (\Omega')^\bullet =
M \otimes_B (\Omega')^\bullet
\end{equation}
qui envoie $m \otimes \eta$ sur $m \otimes \varphi(\eta)$.
\end{remark}

\begin{lemma}
\label{lemma-differentials-completion}
Soit $A \to B$ un homomorphisme d'anneaux et soit $(J, \delta)$ une structure de
puissances divisées sur $B$. Soit $p$ un nombre premier. Supposons que $A$ soit une
$\mathbf{Z}_{(p)}$-algèbre et que $p$ soit nilpotent dans $B/J$. Alors
on a
$$
\lim_e \Omega_{B_e/A, \bar\delta} =
\lim_e \Omega_{B/A, \delta}/p^e\Omega_{B/A, \delta} =
\lim_e \Omega_{B^\wedge/A, \delta^\wedge}/p^e \Omega_{B^\wedge/A, \delta^\wedge}
$$
Pour les notations et les explications, voir la démonstration.
\end{lemma}

\begin{proof}
D'après Algèbre à puissances divisées, lemme \ref{dpa-lemma-extend-to-completion},
la structure $\delta$ se prolonge
à $B_e = B/p^eB$ pour tout $e$ assez grand. La première limite est donc
bien définie. Le lemme fournit aussi une structure de puissances divisées $\delta^\wedge$
sur le complété $B^\wedge = \lim_e B_e$ ; la dernière limite est donc
elle aussi bien définie. D'après le lemme \ref{lemma-omega}
et le fait que $\text{d}p^e = 0$ (toujours),
on voit que la surjection
$\Omega_{B/A, \delta} \to \Omega_{B_e/A, \bar\delta}$ a pour noyau
$p^e\Omega_{B/A, \delta}$. Il en va de même du noyau de
$\Omega_{B^\wedge/A, \delta^\wedge} \to \Omega_{B_e/A, \bar\delta}$.
Le lemme en résulte.
\end{proof}



\section{Schémas à puissances divisées}
\label{section-divided-power-schemes}

\noindent
Quelques remarques sur la globalisation des notions précédentes.

\begin{definition}
\label{definition-divided-power-structure}
Soit $\mathcal{C}$ un site. Soit $\mathcal{O}$ un faisceau d'anneaux
sur $\mathcal{C}$. Soit $\mathcal{I} \subset \mathcal{O}$ un
faisceau d'idéaux. Une {\it structure de puissances divisées $\gamma$} sur $\mathcal{I}$
est une suite d'applications $\gamma_n : \mathcal{I} \to \mathcal{I}$, $n \geq 1$,
telle que, pour tout objet $U$ de $\mathcal{C}$, le triplet
$$
(\mathcal{O}(U), \mathcal{I}(U), \gamma)
$$
soit un anneau à puissances divisées.
\end{definition}

\noindent
Bien entendu, cela s'applique en particulier aux faisceaux d'anneaux sur
les espaces topologiques. Il est toutefois utile de se placer dans ce cadre plus général,
car le faisceau structural du site cristallin vit sur un... site !
Dans ce chapitre, un triplet $(\mathcal{C}, \mathcal{I}, \gamma)$ comme dans la
définition précédente est parfois appelé {\it topos à puissances divisées}.
Étant donnés un second triplet $(\mathcal{C}', \mathcal{I}', \gamma')$ et
un morphisme de topos annelés
$(f, f^\sharp) : (\Sh(\mathcal{C}), \mathcal{O}) \to
(\Sh(\mathcal{C}'), \mathcal{O}')$
on dit que $(f, f^\sharp)$ induit un {\it morphisme de topos à
puissances divisées} si $f^\sharp(f^{-1}\mathcal{I}') \subset \mathcal{I}$
et si les diagrammes
$$
\xymatrix{
f^{-1}\mathcal{I}' \ar[d]_{f^{-1}\gamma'_n} \ar[r]_{f^\sharp} &
\mathcal{I} \ar[d]^{\gamma_n} \\
f^{-1}\mathcal{I}' \ar[r]^{f^\sharp} & \mathcal{I}
}
$$
sont commutatifs pour tout $n \geq 1$. Si $f$ provient d'un morphisme de
sites induit par un foncteur $u : \mathcal{C}' \to \mathcal{C}$, cela signifie
simplement que
$$
(\mathcal{O}'(U'), \mathcal{I}'(U'), \gamma')
\longrightarrow
(\mathcal{O}(u(U')), \mathcal{I}(u(U')), \gamma)
$$
est un homomorphisme d'anneaux à puissances divisées pour tout $U' \in \Ob(\mathcal{C}')$.

\medskip\noindent
Dans le cas des schémas, on exige que l'idéal à puissances divisées soit
{\bf quasi-cohérent}. À cela près, la définition est exactement
la même que dans le cas des topos. La voici.

\begin{definition}
\label{definition-divided-power-scheme}
Un {\it schéma à puissances divisées} est un triplet $(S, \mathcal{I}, \gamma)$
où $S$ est un schéma, $\mathcal{I}$ un faisceau quasi-cohérent
d'idéaux et $\gamma$ une structure de puissances divisées sur $\mathcal{I}$.
Un {\it morphisme de schémas à puissances divisées}
$(S, \mathcal{I}, \gamma) \to (S', \mathcal{I}', \gamma')$ est
un morphisme de schémas $f : S \to S'$ tel que
$f^{-1}\mathcal{I}'\mathcal{O}_S \subset \mathcal{I}$ et que
$$
(\mathcal{O}_{S'}(U'), \mathcal{I}'(U'), \gamma')
\longrightarrow
(\mathcal{O}_S(f^{-1}U'), \mathcal{I}(f^{-1}U'), \gamma)
$$
soit un homomorphisme d'anneaux à puissances divisées pour tout ouvert $U' \subset S'$.
\end{definition}

\noindent
Rappelons qu'il existe une correspondance bijective entre les faisceaux
quasi-cohérents d'idéaux et les immersions fermées ; voir
Morphismes, section \ref{morphisms-section-closed-immersions}.
Ainsi, tout schéma à puissances divisées $(T, \mathcal{J}, \gamma)$
détermine une immersion fermée canonique $U \to T$ définie par $\mathcal{J}$.
Réciproquement, une immersion fermée $U \to T$ et une structure de puissances
divisées $\gamma$ sur le faisceau d'idéaux $\mathcal{J}$ associé
à $U \to T$ déterminent un schéma à puissances divisées $(T, \mathcal{J}, \gamma)$.
Dans nombre de situations, on ne considère de tels triplets
$(U, T, \gamma)$ que lorsque le morphisme $U \to T$ est un épaississement ; voir
Compléments sur les morphismes, définition \ref{more-morphisms-definition-thickening}.

\begin{definition}
\label{definition-divided-power-thickening}
Un triplet $(U, T, \gamma)$ comme ci-dessus est appelé {\it épaississement à puissances divisées}
si $U \to T$ est un épaississement.
\end{definition}

\noindent
Les produits fibrés de schémas à puissances divisées existent lorsque l'un des
trois est un épaississement à puissances divisées. Voici un énoncé précis.

\begin{lemma}
\label{lemma-fibre-product}
Soient $f : (T, \mathcal{J}, \delta) \to (S, \mathcal{I}, \gamma)$ et
$f' : (T', \mathcal{J}', \delta') \to (S, \mathcal{I}, \gamma)$
des morphismes de schémas à puissances divisées. Il existe un schéma à puissances divisées
$(T'', \mathcal{J}'', \delta'')$ et un diagramme cartésien
$$
\xymatrix{
T \ar[d]_f & T'' \ar[d] \ar[l] \\
S & T' \ar[l]_{f'}
}
$$
dans la catégorie des schémas à puissances divisées. Le morphisme
$T'' \to T \times_S T'$ est une immersion fermée et le morphisme
$T''_0 \to T_0 \times_{S_0} T'_0$ est un isomorphisme.
\end{lemma}

\begin{proof}
Esquisse. Notons que, par l'intermédiaire de $\Spec(-)$, les deux dernières assertions
sont compatibles avec ce que nous avons vu des coproduits d'algèbres à puissances divisées dans
Algèbre à puissances divisées, remarque \ref{dpa-remark-pushout}.
On construit donc $T''$ comme sous-schéma fermé de $T \times_S T'$
de la façon suivante : pour tous ouverts affines $U \subset S$, $V \subset T$ et
$V' \subset T'$ tels que $f(V), f'(V') \subset U$, on considère le sous-schéma
fermé de $V \times_U V'$ déterminé par la construction de la
remarque \ref{dpa-remark-pushout} d'Algèbre à puissances divisées. Comme les
schémas $V \times_U V'$ sont les membres d'un recouvrement ouvert
de $T \times_S T'$, on procède ainsi : (1) on montre que ces
sous-schémas fermés se recollent, (2) que les
structures de puissances divisées ainsi obtenues se recollent, et (3) que le résultat du recollement
est le produit fibré dans la catégorie des schémas à puissances divisées.
Pour établir (1), il suffit de montrer
que la construction du coproduit dans la catégorie des
anneaux à puissances divisées commute à la localisation (avec une formulation
appropriée) ; cela résulte du
lemme \ref{dpa-lemma-gamma-extends} d'Algèbre à puissances divisées, qui implique
que les puissances divisées se prolongent aux localisations (par platitude).
\end{proof}

\noindent
Faisons l'observation suivante. Supposons que $(U, T, \gamma)$
soit un schéma à puissances divisées. Supposons que $T$ soit un schéma sur $\mathbf{Z}_{(p)}$
et que $p$ soit localement nilpotent sur $U$. Alors
\begin{enumerate}
\item $p$ est localement nilpotent sur $T \Leftrightarrow U \to T$
est un épaississement (voir Algèbre à puissances divisées, lemme \ref{dpa-lemma-nil}), et
\item $p^e\mathcal{O}_T$ est, localement sur $T$, stable par $\gamma$
pour $e \gg 0$ (voir
Algèbre à puissances divisées, lemme \ref{dpa-lemma-extend-to-completion}).
\end{enumerate}
Cela suggère que l'on disposera de bons résultats sur les épaississements à
puissances divisées sous les hypothèses suivantes.

\begin{situation}
\label{situation-global}
Ici, $p$ est un nombre premier et $(S, \mathcal{I}, \gamma)$ est un schéma à
puissances divisées sur $\mathbf{Z}_{(p)}$. On pose $S_0 = V(\mathcal{I}) \subset S$.
Enfin, $X \to S_0$ est un morphisme de schémas tel que $p$ soit
localement nilpotent sur $X$.
\end{situation}

\noindent
C'est dans cette situation que nous définirons le gros et le petit
site cristallin.









\section{Le gros site cristallin}
\label{section-big-site}

\noindent
Définissons d'abord le gros site. Étant donné un schéma à puissances divisées
$(S, \mathcal{I}, \gamma)$, on dit que $(T, \mathcal{J}, \delta)$ est
un schéma à puissances divisées sur $(S, \mathcal{I}, \gamma)$ si
$T$ est muni d'un morphisme $T \to S$ de schémas à puissances
divisées. De même, on dit qu'un épaississement à puissances divisées $(U, T, \delta)$
est un épaississement à puissances divisées sur $(S, \mathcal{I}, \gamma)$
si $T$ est muni d'un morphisme $T \to S$ de schémas à puissances
divisées.

\begin{definition}
\label{definition-divided-power-thickening-X}
Plaçons-nous dans la situation \ref{situation-global}.
\begin{enumerate}
\item Un {\it épaississement à puissances divisées de $X$ relativement à
$(S, \mathcal{I}, \gamma)$} est la donnée d'un épaississement à puissances divisées
$(U, T, \delta)$ sur $(S, \mathcal{I}, \gamma)$
et d'un $S$-morphisme $U \to X$.
\item Un {\it morphisme d'épaississements à puissances divisées de $X$
relativement à $(S, \mathcal{I}, \gamma)$} se définit de la manière
évidente.
\end{enumerate}
La catégorie des épaississements à puissances divisées de $X$ relativement à
$(S, \mathcal{I}, \gamma)$ est notée $\text{CRIS}(X/S, \mathcal{I}, \gamma)$
ou simplement $\text{CRIS}(X/S)$.
\end{definition}

\noindent
Pour tout $(U, T, \delta)$ dans $\text{CRIS}(X/S)$,
$p$ est localement nilpotent sur $T$ ; voir la discussion précédant la
situation \ref{situation-global}.
On peut représenter commodément toutes les données associées à $(U, T, \delta)$
par le diagramme commutatif
$$
\xymatrix{
T \ar[dd] & U \ar[l] \ar[d] \\
& X \ar[d] \\
S & S_0 \ar[l]
}
$$
où $S_0 = V(\mathcal{I}) \subset S$. Les morphismes de $\text{CRIS}(X/S)$
se représentent de même par de grands diagrammes commutatifs. En particulier,
il existe un foncteur d'oubli canonique
\begin{equation}
\label{equation-forget}
\text{CRIS}(X/S) \longrightarrow \Sch/X,\quad
(U, T, \delta) \longmapsto U
\end{equation}
ainsi que son inverse à droite (et adjoint à gauche)
\begin{equation}
\label{equation-endow-trivial}
\Sch/X \longrightarrow \text{CRIS}(X/S),\quad
U \longmapsto (U, U, \emptyset)
\end{equation}
qui est parfois utile.

\begin{lemma}
\label{lemma-divided-power-thickening-fibre-products}
Dans la situation \ref{situation-global}.
La catégorie $\text{CRIS}(X/S)$ admet toutes les limites finies non vides,
en particulier les produits de deux objets et les produits fibrés.
Le foncteur (\ref{equation-forget}) commute aux limites.
\end{lemma}

\begin{proof}
Omis. Indication : voir le lemme \ref{lemma-affine-thickenings-colimits}
pour le cas affine. Voir également
Algèbre à puissances divisées, remarque \ref{dpa-remark-pushout}.
\end{proof}

\begin{lemma}
\label{lemma-divided-power-thickening-base-change-flat}
Dans la situation \ref{situation-global}. Soit
$$
\xymatrix{
(U_3, T_3, \delta_3) \ar[d] \ar[r] & (U_2, T_2, \delta_2) \ar[d] \\
(U_1, T_1, \delta_1) \ar[r] & (U, T, \delta)
}
$$
un carré cartésien dans la catégorie des épaississements à puissances divisées de
$X$ relativement à $(S, \mathcal{I}, \gamma)$. Si $T_2 \to T$ est
plat et $U_2 = T_2 \times_T U$, alors $T_3 = T_1 \times_T T_2$ (comme schémas).
\end{lemma}

\begin{proof}
Cela résulte de ce qu'une structure de puissances divisées se prolonge de manière unique
le long d'un homomorphisme d'anneaux plat. Voir
Algèbre à puissances divisées, lemme \ref{dpa-lemma-gamma-extends}.
\end{proof}

\noindent
Le lemme précédent signifie que le changement de base d'un morphisme plat
d'épaississements à puissances divisées est encore un morphisme plat et qu'il
s'agit en fait du changement de base « usuel » du morphisme. Il en résulte
que la définition suivante est bien définie.

\begin{definition}
\label{definition-big-crystalline-site}
Dans la situation \ref{situation-global}.
\begin{enumerate}
\item Une famille de morphismes $\{(U_i, T_i, \delta_i) \to (U, T, \delta)\}$
d'épaississements à puissances divisées de $X/S$ est un
{\it recouvrement de Zariski, étale, lisse, syntomique ou fppf}
si et seulement si
\begin{enumerate}
\item $U_i = U \times_T T_i$ pour tout $i$, et
\item $\{T_i \to T\}$ est un recouvrement de Zariski, étale, lisse, syntomique ou fppf.
\end{enumerate}
\item Le {\it gros site cristallin} de $X$ sur $(S, \mathcal{I}, \gamma)$
est la catégorie $\text{CRIS}(X/S)$ munie de la topologie de Zariski.
\item Le topos des faisceaux sur $\text{CRIS}(X/S)$ est noté
$(X/S)_{\text{CRIS}}$ ou parfois
$(X/S, \mathcal{I}, \gamma)_{\text{CRIS}}$\footnote{Cela entre en conflit avec
notre convention consistant à noter le topos associé à un site $\mathcal{C}$
par $\Sh(\mathcal{C})$.}.
\end{enumerate}
\end{definition}

\noindent
On dispose de quelques fonctorialités évidentes pour ces topos.

\begin{remark}[Fonctorialité]
\label{remark-functoriality-big-cris}
Soit $p$ un nombre premier.
Soit $(S, \mathcal{I}, \gamma) \to (S', \mathcal{I}', \gamma')$ un
morphisme de schémas à puissances divisées sur $\mathbf{Z}_{(p)}$.
Posons $S_0 = V(\mathcal{I})$ et $S'_0 = V(\mathcal{I}')$.
Soit
$$
\xymatrix{
X \ar[r]_f \ar[d] & Y \ar[d] \\
S_0 \ar[r] & S'_0
}
$$
un diagramme commutatif de morphismes de schémas et supposons que $p$ soit
localement nilpotent sur $X$ et $Y$. On obtient alors un foncteur continu et
cocontinu
$$
\text{CRIS}(X/S) \longrightarrow \text{CRIS}(Y/S')
$$
en faisant correspondre à $(U, T, \delta)$ le triplet $(U, T, \delta)$
où $U \to X \to Y$ est le $S'$-morphisme de $U$ vers $Y$.
On obtient donc un morphisme de topos
$$
f_{\text{CRIS}} : (X/S)_{\text{CRIS}} \longrightarrow (Y/S')_{\text{CRIS}}
$$
voir Sites, section \ref{sites-section-cocontinuous-morphism-topoi}.
\end{remark}

\begin{remark}[Comparaison avec le site de Zariski]
\label{remark-compare-big-zariski}
Dans la situation \ref{situation-global}.
Le foncteur (\ref{equation-forget}) est cocontinu (détails omis) et
commute aux produits et aux produits fibrés
(lemme \ref{lemma-divided-power-thickening-fibre-products}).
On obtient donc un morphisme de topos
$$
U_{X/S} : (X/S)_{\text{CRIS}} \longrightarrow \Sh((\Sch/X)_{Zar})
$$
du gros topos cristallin de $X/S$ vers le gros topos de Zariski de $X$.
Voir Sites, section \ref{sites-section-cocontinuous-morphism-topoi}.
\end{remark}

\begin{remark}[Morphisme structural]
\label{remark-big-structure-morphism}
Dans la situation \ref{situation-global}.
Considérons le sous-schéma fermé $S_0 = V(\mathcal{I}) \subset S$.
Si l'on suppose que $p$ est localement nilpotent sur $S_0$ (ce qui est toujours
le cas en pratique), on obtient une situation comme dans la
définition \ref{definition-divided-power-thickening-X}, avec $S_0$ au lieu
de $X$. On obtient donc un site $\text{CRIS}(S_0/S)$. Si $f : X \to S_0$ est
le morphisme structural de $X$ sur $S$, on obtient un diagramme commutatif
de morphismes de topos annelés
$$
\xymatrix{
(X/S)_{\text{CRIS}}
\ar[r]_{f_{\text{CRIS}}} \ar[d]_{U_{X/S}} &
(S_0/S)_{\text{CRIS}} \ar[d]^{U_{S_0/S}} \\
\Sh((\Sch/X)_{Zar}) \ar[r]^{f_{big}} & \Sh((\Sch/S_0)_{Zar}) \ar[rd] \\
& & \Sh((\Sch/S)_{Zar})
}
$$
d'après la remarque \ref{remark-functoriality-big-cris}. On considère le composé
$(X/S)_{\text{CRIS}} \to \Sh((\Sch/S)_{Zar})$ comme le morphisme structural du
gros site cristallin. Même si $p$ n'est pas localement nilpotent sur $S_0$,
le morphisme structural
$$
(X/S)_{\text{CRIS}} \longrightarrow \Sh((\Sch/S)_{Zar})
$$
est défini, puisque l'on peut emprunter la voie inférieure du diagramme ci-dessus. C'est donc
le morphisme de topos correspondant au foncteur cocontinu
$\text{CRIS}(X/S) \to (\Sch/S)_{Zar}$ donné par la règle
$(U, T, \delta)/S \mapsto U/S$ ; voir
Sites, section \ref{sites-section-cocontinuous-morphism-topoi}.
\end{remark}

\begin{remark}[Compatibilités]
\label{remark-compatibilities-big-cris}
Les morphismes définis ci-dessus satisfont à de nombreuses compatibilités. Par exemple,
dans la situation de la remarque \ref{remark-functoriality-big-cris},
on obtient un diagramme commutatif de topos annelés
$$
\xymatrix{
(X/S)_{\text{CRIS}} \ar[d] \ar[r] & (Y/S')_{\text{CRIS}} \ar[d] \\
\Sh((\Sch/S)_{Zar}) \ar[r] & \Sh((\Sch/S')_{Zar})
}
$$
où les flèches verticales sont les morphismes structuraux.
\end{remark}




\section{Le site cristallin}
\label{section-site}

\noindent
Comme (\ref{equation-forget}) commute aux produits et aux produits
fibrés, les $(U, T, \delta)$ tels que
$U \to X$ soit une immersion ouverte forment une
sous-catégorie pleine stable par produits fibrés (et, plus généralement,
par limites finies non vides). La
définition suivante est donc bien définie.

\begin{definition}
\label{definition-crystalline-site}
Dans la situation \ref{situation-global}.
\begin{enumerate}
\item Le (petit) {\it site cristallin} de $X$ sur
$(S, \mathcal{I}, \gamma)$, noté $\text{Cris}(X/S, \mathcal{I}, \gamma)$
ou simplement $\text{Cris}(X/S)$, est la sous-catégorie pleine de $\text{CRIS}(X/S)$
formée des $(U, T, \delta)$ de $\text{CRIS}(X/S)$ tels que
$U \to X$ soit une immersion ouverte. Elle est munie de la topologie de Zariski.
\item Le topos des faisceaux sur $\text{Cris}(X/S)$ est noté
$(X/S)_{\text{cris}}$ ou parfois
$(X/S, \mathcal{I}, \gamma)_{\text{cris}}$\footnote{Cela entre en conflit avec
notre convention consistant à noter le topos associé à un site $\mathcal{C}$
par $\Sh(\mathcal{C})$.}.
\end{enumerate}
\end{definition}

\noindent
Pour tout $(U, T, \delta)$ dans $\text{Cris}(X/S)$, le morphisme $U \to X$
définit un objet du petit site de Zariski $X_{Zar}$ de $X$. D'où
On obtient donc un foncteur d'oubli canonique
\begin{equation}
\label{equation-forget-small}
\text{Cris}(X/S) \longrightarrow X_{Zar},\quad
(U, T, \delta) \longmapsto U
\end{equation}
et un adjoint à gauche
\begin{equation}
\label{equation-endow-trivial-small}
X_{Zar} \longrightarrow \text{Cris}(X/S),\quad
U \longmapsto (U, U, \emptyset)
\end{equation}
qui est parfois utile.

\medskip\noindent
On peut comparer le petit et le gros site cristallin, de même
que l'on compare le petit et le gros site de Zariski d'un schéma ; voir
Topologies, lemme \ref{topologies-lemma-at-the-bottom}.

\begin{lemma}
\label{lemma-compare-big-small}
Sous les hypothèses de la définition \ref{definition-divided-power-thickening-X}.
Le foncteur d'inclusion
$$
\text{Cris}(X/S) \to \text{CRIS}(X/S)
$$
commute aux limites finies non vides et est pleinement fidèle, continu
et cocontinu. Il existe des morphismes de topos
$$
(X/S)_{\text{cris}} \xrightarrow{i} (X/S)_{\text{CRIS}}
\xrightarrow{\pi} (X/S)_{\text{cris}}
$$
dont le composé est l'identité et dont le premier est induit
par le foncteur d'inclusion. De plus, $\pi_* = i^{-1}$.
\end{lemma}

\begin{proof}
Pour la première assertion, voir le
lemme \ref{lemma-divided-power-thickening-fibre-products}.
Cela fournit un morphisme de topos
$i : (X/S)_{\text{cris}} \to (X/S)_{\text{CRIS}}$ et un adjoint à gauche
$i_!$ tel que $i^{-1}i_! = i^{-1}i_* = \text{id}$ ; voir
Sites, lemmes \ref{sites-lemma-when-shriek},
\ref{sites-lemma-preserve-equalizers} et
\ref{sites-lemma-back-and-forth}.
Nous affirmons que $i_!$ est exact. Si tel est le cas, on peut définir
$\pi$ par les formules $\pi^{-1} = i_!$ et $\pi_* = i^{-1}$,
et tout est clair. Pour démontrer l'assertion, rappelons que l'on sait déjà
que $i_!$ est exact à droite et préserve les produits fibrés (voir les références
indiquées). Il suffit donc de montrer que $i_! * = *$, où $*$ désigne
l'objet final de la catégorie des faisceaux d'ensembles.
Pour cela, il suffit de produire un ensemble d'objets
$(U_i, T_i, \delta_i)$, $i \in I$, de $\text{Cris}(X/S)$ tel que
$$
\coprod\nolimits_{i \in I} h_{(U_i, T_i, \delta_i)} \to *
$$
soit surjectif dans $(X/S)_{\text{CRIS}}$ (détails omis ; indication : utiliser le fait que
$\text{Cris}(X/S)$ admet des produits et que le foncteur
$\text{Cris}(X/S) \to \text{CRIS}(X/S)$ commute aux produits).
Dans le cas affine, cela
résulte du lemme \ref{lemma-set-generators}. Nous omettons la démonstration
dans le cas général.
\end{proof}

\begin{remark}[Fonctorialité]
\label{remark-functoriality-cris}
Soit $p$ un nombre premier.
Soit $(S, \mathcal{I}, \gamma) \to (S', \mathcal{I}', \gamma')$
un morphisme de schémas à puissances divisées sur $\mathbf{Z}_{(p)}$.
Soit
$$
\xymatrix{
X \ar[r]_f \ar[d] & Y \ar[d] \\
S_0 \ar[r] & S'_0
}
$$
un diagramme commutatif de morphismes de schémas et supposons que $p$ soit
localement nilpotent sur $X$ et $Y$. Par analogie avec
Topologies, lemme \ref{topologies-lemma-morphism-big-small}, on définit
$$
f_{\text{cris}} : (X/S)_{\text{cris}} \longrightarrow (Y/S')_{\text{cris}}
$$
par la formule $f_{\text{cris}} = \pi_Y \circ f_{\text{CRIS}} \circ i_X$,
où $i_X$ et $\pi_Y$ sont ceux du lemme \ref{lemma-compare-big-small}
pour $X$ et $Y$, et où $f_{\text{CRIS}}$ est celui de la
remarque \ref{remark-functoriality-big-cris}.
\end{remark}

\begin{remark}[Comparaison avec le site de Zariski]
\label{remark-compare-zariski}
Dans la situation \ref{situation-global}.
Le foncteur (\ref{equation-forget-small}) est continu, cocontinu et
commute aux produits et aux produits fibrés.
On obtient donc un morphisme de topos
$$
u_{X/S} : (X/S)_{\text{cris}} \longrightarrow \Sh(X_{Zar})
$$
reliant le petit topos cristallin de $X/S$ au
petit topos de Zariski de $X$.
Voir Sites, section \ref{sites-section-cocontinuous-morphism-topoi}.
\end{remark}

\begin{lemma}
\label{lemma-localize}
Dans la situation \ref{situation-global}.
Soient $X' \subset X$ et $S' \subset S$ des sous-schémas ouverts tels que
l'image de $X'$ soit contenue dans $S'$. Il existe alors un foncteur pleinement fidèle
$\text{Cris}(X'/S') \to \text{Cris}(X/S)$
qui donne naissance à un morphisme de topos s'insérant dans le diagramme
commutatif
$$
\xymatrix{
(X'/S')_{\text{cris}} \ar[r] \ar[d]_{u_{X'/S'}} &
(X/S)_{\text{cris}} \ar[d]^{u_{X/S}} \\
\Sh(X'_{Zar}) \ar[r] & \Sh(X_{Zar})
}
$$
De plus, ce diagramme est un exemple de localisation de morphismes de
topos au sens de Sites, lemme \ref{sites-lemma-localize-morphism-topoi}.
\end{lemma}

\begin{proof}
Le foncteur pleinement fidèle consiste à considérer les
objets de $\text{Cris}(X'/S')$ comme des épaississements à puissances
divisées $(U, T, \delta)$ de $X$ pour lesquels $U \to X$
se factorise par $X' \subset X$ (alors $T \to S$ se
factorise automatiquement par $S'$). Ce foncteur est manifestement cocontinu,
d'où le morphisme de topos indiqué.
Soit $h_{X'} \in \Sh(X_{Zar})$ le faisceau représentable associé
à $X'$, considéré comme objet de $X_{Zar}$. Il est clair que
$\Sh(X'_{Zar})$ est le topos localisé $\Sh(X_{Zar})/h_{X'}$.
D'autre part, la catégorie $\text{Cris}(X/S)/u_{X/S}^{-1}h_{X'}$
(voir Sites, lemme \ref{sites-lemma-localize-topos-site})
s'identifie canoniquement à $\text{Cris}(X'/S')$ par le foncteur précédent.
Cela achève la démonstration.
\end{proof}

\begin{remark}[Morphisme structural]
\label{remark-structure-morphism}
Dans la situation \ref{situation-global}.
Considérons le sous-schéma fermé $S_0 = V(\mathcal{I}) \subset S$.
Si l'on suppose que $p$ est localement nilpotent sur $S_0$ (ce qui est toujours
le cas en pratique), on obtient une situation comme dans la
définition \ref{definition-divided-power-thickening-X}, avec $S_0$ au lieu
de $X$. On obtient donc un site $\text{Cris}(S_0/S)$. Si $f : X \to S_0$
est le morphisme structural de $X$ sur $S$, on obtient un
diagramme commutatif de topos
$$
\xymatrix{
(X/S)_{\text{cris}}
\ar[r]_{f_{\text{cris}}} \ar[d]_{u_{X/S}} &
(S_0/S)_{\text{cris}} \ar[d]^{u_{S_0/S}} \\
\Sh(X_{Zar}) \ar[r]^{f_{small}} & \Sh(S_{0, Zar}) \ar[rd] \\
& & \Sh(S_{Zar})
}
$$
voir la remarque \ref{remark-functoriality-cris}. On considère le composé
$(X/S)_{\text{cris}} \to \Sh(S_{Zar})$ comme le morphisme structural du
site cristallin. Même si $p$ n'est pas localement nilpotent sur $S_0$,
le morphisme structural
$$
\tau_{X/S} : (X/S)_{\text{cris}} \longrightarrow \Sh(S_{Zar})
$$
est défini, car on peut suivre le chemin inférieur du diagramme ci-dessus.
\end{remark}

\begin{remark}[Compatibilités]
\label{remark-compatibilities}
Les morphismes définis ci-dessus satisfont à de nombreuses compatibilités. Par exemple,
dans la situation de la remarque \ref{remark-functoriality-cris},
on obtient un diagramme commutatif de topos annelés
$$
\xymatrix{
(X/S)_{\text{cris}} \ar[d] \ar[r] & (Y/S')_{\text{cris}} \ar[d] \\
\Sh((\Sch/S)_{Zar}) \ar[r] & \Sh((\Sch/S')_{Zar})
}
$$
où les flèches verticales sont les morphismes structuraux.
\end{remark}



\section{Faisceaux sur le site cristallin}
\label{section-sheaves}

\noindent
On reprend les notations et les hypothèses de la situation \ref{situation-global}.
Afin de traiter simultanément les petits et gros sites cristallins de $X/S$
dans cette section, posons
$$
\mathcal{C} = \text{CRIS}(X/S)
\quad\text{ou}\quad
\mathcal{C} = \text{Cris}(X/S).
$$
Un faisceau $\mathcal{F}$ sur $\mathcal{C}$ détermine une
{\it restriction} $\mathcal{F}_T$ pour tout objet $(U, T, \delta)$
de $\mathcal{C}$. Plus précisément, $\mathcal{F}_T$ est le faisceau de Zariski sur
le schéma $T$ défini par la règle
$$
\mathcal{F}_T(W) = \mathcal{F}(U \cap W, W, \delta|_W)
$$
pour tout ouvert $W \subset T$. De plus, si $f : T \to T'$ est un morphisme
entre les objets
$(U, T, \delta)$ et $(U', T', \delta')$ de $\mathcal{C}$, il existe
un morphisme canonique de {\it comparaison}
\begin{equation}
\label{equation-comparison}
c_f : f^{-1}\mathcal{F}_{T'} \longrightarrow \mathcal{F}_T.
\end{equation}
En effet, si $W' \subset T'$ est ouvert, alors $f$ induit un morphisme
$$
f|_{f^{-1}W'} :
(U \cap f^{-1}(W'), f^{-1}W', \delta|_{f^{-1}W'})
\longrightarrow
(U' \cap W', W', \delta|_{W'})
$$
de $\mathcal{C}$ ; on peut donc utiliser l'application de restriction
$(f|_{f^{-1}W'})^*$ de $\mathcal{F}$ pour définir une application
$\mathcal{F}_{T'}(W') \to \mathcal{F}_T(f^{-1}W')$.
Ces applications sont manifestement compatibles avec les restrictions ultérieures et
définissent donc un $f$-morphisme de $\mathcal{F}_{T'}$ vers $\mathcal{F}_T$ (voir
Faisceaux, section \ref{sheaves-section-presheaves-functorial},
et en particulier
Faisceaux, définition \ref{sheaves-definition-f-map}).
On obtient ainsi un morphisme $c_f$ comme dans (\ref{equation-comparison}).
Notons que, si $f$ est une immersion ouverte, alors $c_f$ est un
isomorphisme, car $\mathcal{F}_T$ n'est dans ce cas que
la restriction de $\mathcal{F}_{T'}$ à $T$.

\medskip\noindent
Réciproquement, si l'on se donne des faisceaux de Zariski $\mathcal{F}_T$ pour tout objet
$(U, T, \delta)$ de $\mathcal{C}$ et des morphismes de comparaison
$c_f$ comme ci-dessus qui (a) sont des isomorphismes pour les immersions ouvertes et (b)
satisfont à une condition de cocycle convenable, on obtient un faisceau sur
$\mathcal{C}$. La démonstration est exactement celle de
Topologies, lemme \ref{topologies-lemma-characterize-sheaf-big}.

\medskip\noindent
Le {\it faisceau structural} sur $\mathcal{C}$ est le faisceau
$\mathcal{O}_{X/S}$ défini par la règle
$$
\mathcal{O}_{X/S} :
(U, T, \delta)
\longmapsto
\Gamma(T, \mathcal{O}_T)
$$
C'est un faisceau par définition des recouvrements dans $\mathcal{C}$.
Supposons que $\mathcal{F}$ soit un faisceau de $\mathcal{O}_{X/S}$-modules.
Dans ce cas, les applications de comparaison (\ref{equation-comparison})
définissent un morphisme de comparaison
\begin{equation}
\label{equation-comparison-modules}
c_f : f^*\mathcal{F}_{T'} \longrightarrow \mathcal{F}_T
\end{equation}
de $\mathcal{O}_T$-modules.

\medskip\noindent
Un autre type d'exemple s'obtient en partant d'un faisceau
$\mathcal{G}$ sur $(\Sch/X)_{Zar}$ ou sur $X_{Zar}$ (suivant que
$\mathcal{C} = \text{CRIS}(X/S)$ ou $\mathcal{C} = \text{Cris}(X/S)$).
Alors $\underline{\mathcal{G}}$, défini par la règle
$$
\underline{\mathcal{G}} :
(U, T, \delta)
\longmapsto
\mathcal{G}(U)
$$
est un faisceau sur $\mathcal{C}$. En particulier, si l'on prend
$\mathcal{G} = \mathbf{G}_a = \mathcal{O}_X$, on obtient
$$
\underline{\mathbf{G}_a} :
(U, T, \delta)
\longmapsto
\Gamma(U, \mathcal{O}_U)
$$
Il existe un morphisme surjectif de faisceaux
$\mathcal{O}_{X/S} \to \underline{\mathbf{G}_a}$ défini par les
morphismes canoniques $\Gamma(T, \mathcal{O}_T) \to \Gamma(U, \mathcal{O}_U)$
pour les objets $(U, T, \delta)$. Le noyau de ce morphisme est noté
$\mathcal{J}_{X/S}$, d'où une suite exacte courte
$$
0 \to
\mathcal{J}_{X/S} \to
\mathcal{O}_{X/S} \to
\underline{\mathbf{G}_a} \to 0
$$
Notons que $\mathcal{J}_{X/S}$ est muni d'une structure canonique de
puissances divisées. En effet, pour chaque objet $(U, T, \delta)$,
la troisième composante $\delta$ {\it est} une structure de puissances divisées sur le
noyau de $\mathcal{O}_T \to \mathcal{O}_U$. Ainsi, le (gros)
topos cristallin est un topos à puissances divisées.





\section{Cristaux en modules}
\label{section-crystals}

\noindent
Un cristal est en fait un objet très général. Toutefois, sa
définition peut être un peu délicate à analyser ; nous la donnons donc d'abord
dans le cas des modules sur les sites cristallins.

\begin{definition}
\label{definition-modules}
Dans la situation \ref{situation-global}.
Soit $\mathcal{C} = \text{CRIS}(X/S)$ ou $\mathcal{C} = \text{Cris}(X/S)$.
Soit $\mathcal{F}$ un faisceau de $\mathcal{O}_{X/S}$-modules sur $\mathcal{C}$.
\begin{enumerate}
\item On dit que $\mathcal{F}$ est {\it localement quasi-cohérent} si, pour tout
objet $(U, T, \delta)$ de $\mathcal{C}$, la restriction $\mathcal{F}_T$
est un $\mathcal{O}_T$-module quasi-cohérent.
\item On dit que $\mathcal{F}$ est {\it quasi-cohérent} s'il est quasi-cohérent
au sens de
Modules sur les sites, définition \ref{sites-modules-definition-site-local}.
\item On dit que $\mathcal{F}$ est un {\it cristal en $\mathcal{O}_{X/S}$-modules}
si tous les morphismes de comparaison (\ref{equation-comparison-modules}) sont des
isomorphismes.
\end{enumerate}
\end{definition}

\noindent
Ces notions sont reliées de la manière suivante.

\begin{lemma}
\label{lemma-crystal-quasi-coherent-modules}
Avec les notations $X/S, \mathcal{I}, \gamma, \mathcal{C}, \mathcal{F}$
de la définition \ref{definition-modules}, les conditions suivantes sont équivalentes :
\begin{enumerate}
\item $\mathcal{F}$ est quasi-cohérent, et
\item $\mathcal{F}$ est localement quasi-cohérent et est un cristal en
$\mathcal{O}_{X/S}$-modules.
\end{enumerate}
\end{lemma}

\begin{proof}
Supposons (1). Soit $f : (U', T', \delta') \to (U, T, \delta)$ un morphisme de
$\mathcal{C}$. Il faut montrer (a) que $\mathcal{F}_T$ est un
$\mathcal{O}_T$-module quasi-cohérent et (b) que $c_f : f^*\mathcal{F}_T \to \mathcal{F}_{T'}$
est un isomorphisme. L'hypothèse signifie qu'il existe un recouvrement
$\{(U_i, T_i, \delta_i) \to (U, T, \delta)\}$ tel que, pour chaque $i$,
la restriction de $\mathcal{F}$ à $\mathcal{C}/(U_i, T_i, \delta_i)$
admette une présentation globale. Comme il suffit de prouver (a) et (b)
localement pour la topologie de Zariski, on peut remplacer $f : (U', T', \delta') \to (U, T, \delta)$
par son changement de base à $(U_i, T_i, \delta_i)$ et supposer que la restriction de $\mathcal{F}$
à $\mathcal{C}/(U, T, \delta)$ admet une
présentation globale
$$
\bigoplus\nolimits_{j \in J}
\mathcal{O}_{X/S}|_{\mathcal{C}/(U, T, \delta)} \longrightarrow
\bigoplus\nolimits_{i \in I}
\mathcal{O}_{X/S}|_{\mathcal{C}/(U, T, \delta)} \longrightarrow
\mathcal{F}|_{\mathcal{C}/(U, T, \delta)}
\longrightarrow 0
$$
Il est clair que cela donne une présentation
$$
\bigoplus\nolimits_{j \in J} \mathcal{O}_T \longrightarrow
\bigoplus\nolimits_{i \in I} \mathcal{O}_T \longrightarrow
\mathcal{F}_T
\longrightarrow 0
$$
et donc (a) est vérifiée. De plus, cette présentation se restreint à $T'$
en une présentation analogue de $\mathcal{F}_{T'}$, ce qui établit (b).

\medskip\noindent
Supposons (2). Soit $(U, T, \delta)$ un objet de $\mathcal{C}$.
Il faut trouver un recouvrement de $(U, T, \delta)$ tel que $\mathcal{F}$ admette une
présentation globale après restriction aux localisations de $\mathcal{C}$
associées aux membres du recouvrement. On peut donc supposer $T$ affine.
Dans ce cas, on peut choisir une présentation
$$
\bigoplus\nolimits_{j \in J} \mathcal{O}_T \longrightarrow
\bigoplus\nolimits_{i \in I} \mathcal{O}_T \longrightarrow
\mathcal{F}_T
\longrightarrow 0
$$
puisque $\mathcal{F}_T$ est supposé être un $\mathcal{O}_T$-module quasi-cohérent.
La propriété de cristal de $\mathcal{F}$ montre alors que l'image réciproque de cette présentation donne
une présentation de $\mathcal{F}_{T'}$ pour tout morphisme
$f : (U', T', \delta') \to (U, T, \delta)$ de $\mathcal{C}$.
On obtient ainsi la présentation voulue de $\mathcal{F}|_{\mathcal{C}/(U, T, \delta)}$.
\end{proof}

\begin{definition}
\label{definition-crystal-quasi-coherent-modules}
Si $\mathcal{F}$ satisfait aux conditions équivalentes du
lemme \ref{lemma-crystal-quasi-coherent-modules}, on dit
que $\mathcal{F}$ est un
{\it cristal en modules quasi-cohérents}.
On dit que $\mathcal{F}$ est un {\it cristal en modules localement libres de type fini}
si, de plus, $\mathcal{F}$ est localement libre de type fini.
\end{definition}

\noindent
Bien entendu, comme le montre le lemme \ref{lemma-crystal-quasi-coherent-modules}, cette
terminologie est quelque peu lourde, puisqu'un module quasi-cohérent est toujours un cristal.
Mais c'est la terminologie usuelle dans la littérature.

\begin{remark}
\label{remark-crystal}
Pour formuler la notion générale de cristal, on utilise le langage
des champs et des morphismes fortement cartésiens ; voir
Champs, définition \ref{stacks-definition-stack} et
Catégories, définition \ref{categories-definition-cartesian-over-C}.
Dans la situation \ref{situation-global}, soit
$p : \mathcal{C} \to \text{Cris}(X/S)$ un champ.
Un {\it cristal en objets de $\mathcal{C}$ sur $X$ relativement à $S$}
est une {\it section cartésienne} $\sigma : \text{Cris}(X/S) \to \mathcal{C}$,
c'est-à-dire un foncteur $\sigma$ tel que $p \circ \sigma = \text{id}$
et tel que $\sigma(f)$ soit fortement cartésien pour tout
morphisme $f$ de $\text{Cris}(X/S)$. De même pour le gros site cristallin.
\end{remark}





\section{Faisceau des différentielles}
\label{section-differentials-sheaf}

\noindent
Dans cette section, nous nous en tiendrons au (petit) site cristallin,
qui paraît plus naturel. On globalise
la définition \ref{definition-derivation} comme suit.

\begin{definition}
\label{definition-global-derivation}
Dans la situation \ref{situation-global}, soit
$\mathcal{F}$ un faisceau de $\mathcal{O}_{X/S}$-modules sur
$\text{Cris}(X/S)$. Une
{\it $S$-dérivation $D : \mathcal{O}_{X/S} \to \mathcal{F}$}
est un morphisme de faisceaux tel que, pour tout objet $(U, T, \delta)$ de
$\text{Cris}(X/S)$, l'application
$$
D : \Gamma(T, \mathcal{O}_T) \longrightarrow \Gamma(T, \mathcal{F})
$$
soit une $\Gamma(V, \mathcal{O}_V)$-dérivation à puissances divisées, où $V \subset S$
est un ouvert quelconque tel que $T \to S$ se factorise par $V$.
\end{definition}

\noindent
Cela signifie que $D$ est additive, satisfait à la règle de Leibniz, s'annule sur les
fonctions provenant de $S$ et vérifie $D(f^{[n]}) = f^{[n - 1]}D(f)$
pour toute section locale $f$ de l'idéal à puissances divisées $\mathcal{J}_{X/S}$.
C'est un cas particulier d'une notion très générale que nous allons maintenant décrire.

\medskip\noindent
Comparons la discussion qui suit avec
Modules sur les sites, section \ref{sites-modules-section-differentials}. Soit
$\mathcal{C}$ un site, soit $\mathcal{A} \to \mathcal{B}$ un
morphisme de faisceaux d'anneaux sur $\mathcal{C}$, soit $\mathcal{J} \subset \mathcal{B}$
un faisceau d'idéaux, soit $\delta$ une structure de puissances divisées sur
$\mathcal{J}$ et soit $\mathcal{F}$ un faisceau de $\mathcal{B}$-modules.
Il existe alors une notion de {\it $\mathcal{A}$-dérivation à puissances divisées}
$D : \mathcal{B} \to \mathcal{F}$. Cela signifie que $D$ est $\mathcal{A}$-linéaire,
satisfait à la règle de Leibniz et vérifie
$D(\delta_n(x)) = \delta_{n - 1}(x)D(x)$ pour les sections locales $x$ de
$\mathcal{J}$. Dans cette situation, il existe une
{\it $\mathcal{A}$-dérivation universelle à puissances divisées}
$$
\text{d}_{\mathcal{B}/\mathcal{A}, \delta} :
\mathcal{B}
\longrightarrow
\Omega_{\mathcal{B}/\mathcal{A}, \delta}
$$
De plus, $\text{d}_{\mathcal{B}/\mathcal{A}, \delta}$ est le composé
$$
\mathcal{B}
\longrightarrow
\Omega_{\mathcal{B}/\mathcal{A}}
\longrightarrow
\Omega_{\mathcal{B}/\mathcal{A}, \delta}
$$
où le premier morphisme est la dérivation universelle construite dans la démonstration
de Modules sur les sites, lemme \ref{sites-modules-lemma-universal-module},
et le second est le morphisme quotient par le sous-module engendré par
les sections locales
$\text{d}_{\mathcal{B}/\mathcal{A}}(\delta_n(x)) -
\delta_{n - 1}(x)\text{d}_{\mathcal{B}/\mathcal{A}}(x)$.

\medskip\noindent
Passons à la version relative. Supposons que
$(f, f^\sharp) : (\Sh(\mathcal{C}), \mathcal{O}) \to
(\Sh(\mathcal{C}'), \mathcal{O}')$ soit un morphisme de topos annelés,
que $\mathcal{J} \subset \mathcal{O}$ soit un faisceau d'idéaux, que $\delta$ soit une
structure de puissances divisées sur $\mathcal{J}$ et que $\mathcal{F}$ soit un faisceau
de $\mathcal{O}$-modules. Dans cette situation, on dit que
$D : \mathcal{O} \to \mathcal{F}$ est une $\mathcal{O}'$-dérivation à puissances divisées
si $D$ est une $f^{-1}\mathcal{O}'$-dérivation à puissances divisées au sens défini ci-dessus.
En outre, on écrit
$$
\Omega_{\mathcal{O}/\mathcal{O}', \delta} =
\Omega_{\mathcal{O}/f^{-1}\mathcal{O}', \delta}
$$
qui reçoit la
$\mathcal{O}'$-dérivation universelle à puissances divisées.

\medskip\noindent
En appliquant cette construction au morphisme structural
$$
(X/S)_{\text{Cris}} \longrightarrow \Sh(S_{Zar})
$$
(voir la remarque \ref{remark-structure-morphism}), on retrouve la notion
définie dans la définition \ref{definition-global-derivation}.
En particulier, il existe une dérivation universelle à puissances divisées
$$
d_{X/S} : \mathcal{O}_{X/S} \to \Omega_{X/S}
$$
Notons que nous omettons dans la notation l'indication de la compatibilité du
module des différentielles avec les puissances divisées (il paraît en effet
peu probable que l'on considère jamais le module usuel des différentielles du
faisceau structural sur le site cristallin).

\begin{lemma}
\label{lemma-module-differentials-divided-power-scheme}
Soit $(T, \mathcal{J}, \delta)$ un schéma à puissances divisées.
Soit $T \to S$ un morphisme de schémas.
Le quotient $\Omega_{T/S} \to \Omega_{T/S, \delta}$
décrit ci-dessus est un $\mathcal{O}_T$-module quasi-cohérent.
Pour tout ouvert affine $W \subset T$ dont l'image est contenue dans un
ouvert affine $V \subset S$, on a
$$
\Gamma(W, \Omega_{T/S, \delta}) =
\Omega_{\Gamma(W, \mathcal{O}_W)/\Gamma(V, \mathcal{O}_V), \delta}
$$
où le second membre est
celui construit dans la section \ref{section-differentials}.
\end{lemma}

\begin{proof}
Omis.
\end{proof}

\begin{lemma}
\label{lemma-module-of-differentials}
Dans la situation \ref{situation-global}.
Pour $(U, T, \delta)$ dans $\text{Cris}(X/S)$, la restriction
$(\Omega_{X/S})_T$ à $T$ est $\Omega_{T/S, \delta}$ et la restriction
$\text{d}_{X/S}|_T$ est égale à $\text{d}_{T/S, \delta}$.
\end{lemma}

\begin{proof}
Omis.
\end{proof}

\begin{lemma}
\label{lemma-module-of-differentials-on-affine}
Dans la situation \ref{situation-global}.
Pour tout objet affine $(U, T, \delta)$ de $\text{Cris}(X/S)$
s'envoyant dans un ouvert affine $V \subset S$, on a
$$
\Gamma((U, T, \delta), \Omega_{X/S}) =
\Omega_{\Gamma(T, \mathcal{O}_T)/\Gamma(V, \mathcal{O}_V), \delta}
$$
où le second membre est
celui construit dans la section \ref{section-differentials}.
\end{lemma}

\begin{proof}
Cela résulte des lemmes \ref{lemma-module-differentials-divided-power-scheme} et
\ref{lemma-module-of-differentials}.
\end{proof}

\begin{lemma}
\label{lemma-describe-omega-small}
Dans la situation \ref{situation-global}.
Soit $(U, T, \delta)$ un objet de $\text{Cris}(X/S)$.
Soit
$$
(U(1), T(1), \delta(1)) = (U, T, \delta) \times (U, T, \delta)
$$
dans $\text{Cris}(X/S)$. Soit $\mathcal{K} \subset \mathcal{O}_{T(1)}$
le faisceau quasi-cohérent d'idéaux correspondant à l'immersion
fermée $\Delta : T \to T(1)$. Alors
$\mathcal{K} \subset \mathcal{J}_{T(1)}$ est stable par la
structure de puissances divisées sur $\mathcal{J}_{T(1)}$ et l'on a
$$
(\Omega_{X/S})_T = \mathcal{K}/\mathcal{K}^{[2]}
$$
\end{lemma}

\begin{proof}
Notons que $U = U(1)$, puisque $U \to X$ est une immersion ouverte et que
le foncteur (\ref{equation-forget-small}) commute aux produits. On en déduit que
$\mathcal{K} \subset \mathcal{J}_{T(1)}$. Cela étant, le lemme s'obtient
en travaillant localement sur des ouverts affines de $T$ et en utilisant les
lemmes \ref{lemma-module-of-differentials-on-affine} et
\ref{lemma-diagonal-and-differentials-affine-site}.
\end{proof}

\noindent
Le faisceau $\Omega_{X/S}$ n'est pas un cristal en
$\mathcal{O}_{X/S}$-modules quasi-cohérents. Il satisfait néanmoins à deux propriétés
étroitement apparentées (comparer avec le
lemme \ref{lemma-crystal-quasi-coherent-modules}).

\begin{lemma}
\label{lemma-omega-locally-quasi-coherent}
Dans la situation \ref{situation-global}.
Le faisceau des différentielles $\Omega_{X/S}$ possède les deux
propriétés suivantes :
\begin{enumerate}
\item $\Omega_{X/S}$ est localement quasi-cohérent, et
\item pour tout morphisme $(U, T, \delta) \to (U', T', \delta')$
de $\text{Cris}(X/S)$ tel que $f : T \to T'$ soit une immersion fermée,
le morphisme $c_f : f^*(\Omega_{X/S})_{T'} \to (\Omega_{X/S})_T$ est surjectif.
\end{enumerate}
\end{lemma}

\begin{proof}
L'assertion (1) résulte de la combinaison des
lemmes \ref{lemma-module-differentials-divided-power-scheme} et
\ref{lemma-module-of-differentials}.
L'assertion (2) résulte du fait que
$(\Omega_{X/S})_T = \Omega_{T/S, \delta}$
est un quotient de $\Omega_{T/S}$ et que $f^*\Omega_{T'/S} \to \Omega_{T/S}$
est surjectif.
\end{proof}







\section{Deux épaississements universels}
\label{section-universal-thickenings}

\noindent
Les constructions de cette section nous aideront à définir une connexion sur
un cristal en modules sur le site cristallin. En quelque sorte, les constructions
présentées ici sont les versions « faisceautisées et universelles » de celles de la
section \ref{section-explicit-thickenings}.

\begin{remark}
\label{remark-first-order-thickening}
Dans la situation \ref{situation-global}.
Soit $(U, T, \delta)$ un objet de $\text{Cris}(X/S)$.
Notons $\Omega_{T/S, \delta} = (\Omega_{X/S})_T$ ; voir le
lemme \ref{lemma-module-of-differentials}.
Décrivons explicitement un épaississement du premier ordre $T'$ de
$T$. Plus précisément, posons
$$
\mathcal{O}_{T'} = \mathcal{O}_T \oplus \Omega_{T/S, \delta}
$$
muni de la structure d'algèbre pour laquelle $\Omega_{T/S, \delta}$ est un
idéal de carré nul. Soit $\mathcal{J} \subset \mathcal{O}_T$
le faisceau d'idéaux de l'immersion fermée $U \to T$. Posons
$\mathcal{J}' = \mathcal{J} \oplus \Omega_{T/S, \delta}$.
Définissons une structure de puissances divisées sur $\mathcal{J}'$ en posant
$$
\delta_n'(f, \omega) = (\delta_n(f), \delta_{n - 1}(f)\omega),
$$
voir le lemme \ref{lemma-divided-power-first-order-thickening}.
Il existe deux homomorphismes d'anneaux
$$
p_0, p_1 : \mathcal{O}_T \to \mathcal{O}_{T'}
$$
Le premier est donné par $f \mapsto (f, 0)$ et le second par
$f \mapsto (f, \text{d}_{T/S, \delta}f)$. Notons que tous deux sont compatibles
avec les structures de puissances divisées sur $\mathcal{J}$ et $\mathcal{J}'$,
de même que le morphisme quotient $\mathcal{O}_{T'} \to \mathcal{O}_T$.
On obtient ainsi un objet $(U, T', \delta')$ de $\text{Cris}(X/S)$
et un diagramme commutatif
$$
\xymatrix{
& T \ar[ld]_{\text{id}} \ar[d]^i \ar[rd]^{\text{id}} \\
T & T' \ar[l]_{p_0} \ar[r]^{p_1} & T
}
$$
dans $\text{Cris}(X/S)$ tel que $i$ soit un épaississement du premier ordre dont le faisceau
d'idéaux s'identifie à $\Omega_{T/S, \delta}$ et tel que
$p_1 - p_0 : \mathcal{O}_T \to \mathcal{O}_{T'}$
s'identifie à la dérivation universelle $\text{d}_{T/S, \delta}$
composée avec l'inclusion $\Omega_{T/S, \delta} \to \mathcal{O}_{T'}$.
\end{remark}

\begin{remark}
\label{remark-second-order-thickening}
Dans la situation \ref{situation-global}.
Soit $(U, T, \delta)$ un objet de $\text{Cris}(X/S)$.
Notons $\Omega_{T/S, \delta} = (\Omega_{X/S})_T$ ; voir le
lemme \ref{lemma-module-of-differentials}.
Notons également $\Omega^2_{T/S, \delta}$ sa seconde puissance
extérieure. Décrivons explicitement un épaississement du second ordre $T''$ de $T$.
Plus précisément, posons
$$
\mathcal{O}_{T''} =
\mathcal{O}_T \oplus \Omega_{T/S, \delta} \oplus \Omega_{T/S, \delta}
\oplus \Omega^2_{T/S, \delta}
$$
muni de la structure d'algèbre définie de la manière suivante :
$$
(f, \omega_1, \omega_2, \eta) \cdot
(f', \omega_1', \omega_2', \eta') =
(ff', f\omega_1' + f'\omega_1, f\omega_2' + f'\omega_2,
f\eta' + f'\eta + \omega_1 \wedge \omega_2' + \omega_1' \wedge \omega_2).
$$
Soit $\mathcal{J} \subset \mathcal{O}_T$
le faisceau d'idéaux de l'immersion fermée $U \to T$. Soit
$\mathcal{J}''$ l'image réciproque de $\mathcal{J}$ par la
projection $\mathcal{O}_{T''} \to \mathcal{O}_T$.
Définissons une structure de puissances divisées sur $\mathcal{J}''$ en posant
$$
\delta_n''(f, \omega_1, \omega_2, \eta) =
(\delta_n(f), \delta_{n - 1}(f)\omega_1, \delta_{n - 1}(f)\omega_2,
\delta_{n - 1}(f)\eta + \delta_{n - 2}(f)\omega_1 \wedge \omega_2)
$$
voir le lemme \ref{lemma-divided-power-second-order-thickening}.
Il existe trois homomorphismes d'anneaux
$q_0, q_1, q_2 : \mathcal{O}_T \to \mathcal{O}_{T''}$
donnés par
\begin{align*}
q_0(f) & = (f, 0, 0, 0), \\
q_1(f) & = (f, \text{d}f, 0, 0), \\
q_2(f) & = (f, \text{d}f, \text{d}f, 0)
\end{align*}
où $\text{d} = \text{d}_{T/S, \delta}$.
Notons que tous trois sont compatibles avec les structures de puissances divisées
sur $\mathcal{J}$ et $\mathcal{J}''$. Il existe trois homomorphismes d'anneaux
$q_{01}, q_{12}, q_{02} : \mathcal{O}_{T'} \to \mathcal{O}_{T''}$
où $\mathcal{O}_{T'}$ est défini comme dans la remarque \ref{remark-first-order-thickening}.
Plus précisément, posons
\begin{align*}
q_{01}(f, \omega) & = (f, \omega, 0, 0), \\
q_{12}(f, \omega) & =
(f, \text{d}f, \omega, \text{d}\omega), \\
q_{02}(f, \omega) & = (f, \omega, \omega, 0)
\end{align*}
Ces homomorphismes sont eux aussi compatibles avec les structures de puissances
divisées données. Vérifions ces assertions pour $q_{12}$. Notons
que $q_{12}$ est un homomorphisme d'anneaux, puisque
\begin{align*}
q_{12}(f, \omega)q_{12}(g, \eta) & =
(f, \text{d}f, \omega, \text{d}\omega)(g, \text{d}g, \eta, \text{d}\eta) \\
& =
(fg, f\text{d}g + g \text{d}f, f\eta + g\omega,
f\text{d}\eta + g\text{d}\omega + \text{d}f \wedge \eta +
\text{d}g \wedge \omega) \\
& = q_{12}(fg, f\eta + g\omega) = q_{12}((f, \omega)(g, \eta))
\end{align*}
Notons que $q_{12}$ est compatible avec les puissances divisées, car
\begin{align*}
\delta_n''(q_{12}(f, \omega)) & =
\delta_n''((f, \text{d}f, \omega, \text{d}\omega)) \\
& =
(\delta_n(f), \delta_{n - 1}(f)\text{d}f, \delta_{n - 1}(f)\omega,
\delta_{n - 1}(f)\text{d}\omega + \delta_{n - 2}(f)\text{d}(f) \wedge \omega)
\\
& = q_{12}((\delta_n(f), \delta_{n - 1}(f)\omega)) =
q_{12}(\delta'_n(f, \omega))
\end{align*}
Les vérifications pour $q_{01}$ et $q_{02}$ sont plus faciles.
Notons que $q_0 = q_{01} \circ p_0$, $q_1 = q_{01} \circ p_1$,
$q_1 = q_{12} \circ p_0$, $q_2 = q_{12} \circ p_1$,
$q_0 = q_{02} \circ p_0$ et $q_2 = q_{02} \circ p_1$.
Ainsi $(U, T'', \delta'')$ est un objet de $\text{Cris}(X/S)$
et l'on obtient des morphismes
$$
\xymatrix{
T''
\ar@<2ex>[r]
\ar@<0ex>[r]
\ar@<-2ex>[r]
&
T'
\ar@<1ex>[r]
\ar@<-1ex>[r]
&
T
}
$$
de $\text{Cris}(X/S)$ satisfaisant aux relations décrites ci-dessus.
Dans les applications, les notations $q_i : T'' \to T$ et
$q_{ij} : T'' \to T'$ désigneront les morphismes associés aux
homomorphismes d'anneaux décrits ci-dessus.
\end{remark}






\section{Le complexe de de Rham}
\label{section-de-Rham}

\noindent
Dans la situation \ref{situation-global}.
Sur le (petit) site cristallin, on définit
$\Omega^i_{X/S} = \wedge^i_{\mathcal{O}_{X/S}} \Omega_{X/S}$
pour $i \geq 0$. La $S$-dérivation universelle $\text{d}_{X/S}$ définit
le {\it complexe de de Rham}
$$
\mathcal{O}_{X/S} \to \Omega^1_{X/S} \to \Omega^2_{X/S} \to \ldots
$$
sur $\text{Cris}(X/S)$ ; voir le
lemme \ref{lemma-module-of-differentials-on-affine} et la
remarque \ref{remark-divided-powers-de-rham-complex}.


\section{Connexions}
\label{section-connections}

\noindent
Dans la situation \ref{situation-global}.
Étant donné un $\mathcal{O}_{X/S}$-module $\mathcal{F}$ sur $\text{Cris}(X/S)$,
une {\it connexion} est un morphisme de faisceaux abéliens
$$
\nabla :
\mathcal{F}
\longrightarrow
\mathcal{F} \otimes_{\mathcal{O}_{X/S}} \Omega_{X/S}
$$
tel que $\nabla(f s) = f\nabla(s) + s \otimes \text{d}f$
pour les sections locales $s, f$ de $\mathcal{F}$ et de $\mathcal{O}_{X/S}$, respectivement.
Une connexion détermine des morphismes canoniques
$
\nabla :
\mathcal{F} \otimes_{\mathcal{O}_{X/S}} \Omega^i_{X/S}
\longrightarrow
\mathcal{F} \otimes_{\mathcal{O}_{X/S}} \Omega^{i + 1}_{X/S}
$
définis par la règle $\nabla(s \otimes \omega) =
\nabla(s) \wedge \omega + s \otimes \text{d}\omega$
comme dans la remarque \ref{remark-connection}. On dit que la connexion est
{\it intégrable} si $\nabla \circ \nabla = 0$. Si $\nabla$ est intégrable,
on obtient le {\it complexe de de Rham}
$$
\mathcal{F} \to
\mathcal{F} \otimes_{\mathcal{O}_{X/S}} \Omega^1_{X/S} \to
\mathcal{F} \otimes_{\mathcal{O}_{X/S}} \Omega^2_{X/S} \to \ldots
$$
sur $\text{Cris}(X/S)$. Tout cristal en
$\mathcal{O}_{X/S}$-modules est en fait muni d'une connexion
intégrable canonique.

\begin{lemma}
\label{lemma-automatic-connection}
Dans la situation \ref{situation-global}.
Soit $\mathcal{F}$ un cristal en $\mathcal{O}_{X/S}$-modules
sur $\text{Cris}(X/S)$. Alors $\mathcal{F}$ est muni d'une
connexion intégrable canonique.
\end{lemma}

\begin{proof}
Soit $(U, T, \delta)$ un objet de $\text{Cris}(X/S)$.
Soit $(U, T', \delta')$ l'épaississement infinitésimal de $T$
par $(\Omega_{X/S})_T = \Omega_{T/S, \delta}$
construit dans la remarque \ref{remark-first-order-thickening}.
Il est muni de projections $p_0, p_1 : T' \to T$
et d'une diagonale $i : T \to T'$. Par hypothèse, on obtient des
isomorphismes
$$
p_0^*\mathcal{F}_T \xrightarrow{c_0}
\mathcal{F}_{T'} \xleftarrow{c_1}
p_1^*\mathcal{F}_T
$$
de $\mathcal{O}_{T'}$-modules. L'image réciproque de $c = c_1^{-1} \circ c_0$
sur $T$ par $i$ est l'identité
de $\mathcal{F}_T$. Par conséquent, si $s \in \Gamma(T, \mathcal{F}_T)$,
alors $\nabla(s) = p_1^*s - c(p_0^*s)$ est une section de
$p_1^*\mathcal{F}_T$ dont l'image réciproque par $i$ est nulle. Ainsi,
$\nabla(s)$ est une section de
$$
\mathcal{F}_T
\otimes_{\mathcal{O}_T}
\Omega_{T/S, \delta}
$$
car il s'agit du noyau de $p_1^*\mathcal{F}_T \to \mathcal{F}_T$,
puisque $\mathcal{O}_{T'} = \mathcal{O}_T \oplus \Omega_{T/S, \delta}$
par construction. On vérifie aisément que $\nabla(fs) =
f\nabla(s) + s \otimes \text{d}(f)$ en utilisant la description de
$\text{d}$ dans la remarque \ref{remark-first-order-thickening}.

\medskip\noindent
La famille d'applications
$$
\nabla : \Gamma(T, \mathcal{F}_T) \to
\Gamma(T, \mathcal{F}_T \otimes_{\mathcal{O}_T} \Omega_{T/S, \delta})
$$
ainsi obtenue est fonctorielle en $T$, car la construction de $T'$
est fonctorielle en $T$. On obtient donc une connexion.

\medskip\noindent
Pour montrer que la connexion est intégrable, considérons
l'objet $(U, T'', \delta'')$ construit dans la
remarque \ref{remark-second-order-thickening}.
Comme $\mathcal{F}$ est un faisceau, on voit que
$$
\xymatrix{
q_0^*\mathcal{F}_T \ar[rr]_{q_{01}^*c} \ar[rd]_{q_{02}^*c} & &
q_1^*\mathcal{F}_T \ar[ld]^{q_{12}^*c} \\
& q_2^*\mathcal{F}_T
}
$$
est un diagramme commutatif de $\mathcal{O}_{T''}$-modules. Pour
$s \in \Gamma(T, \mathcal{F}_T)$, on a
$c(p_0^*s) = p_1^*s - \nabla(s)$. Écrivons
$\nabla(s) = \sum p_1^*s_i \cdot \omega_i$, où $s_i$ est une section locale
de $\mathcal{F}_T$ et $\omega_i$ une section locale de $\Omega_{T/S, \delta}$.
On considère $\omega_i$ comme une section locale du faisceau structural
$\mathcal{O}_{T'}$ et l'on écrit donc un produit plutôt qu'un produit
tensoriel. D'une part,
\begin{align*}
q_{12}^*c \circ q_{01}^*c(q_0^*s) & = 
q_{12}^*c(q_1^*s - \sum q_1^*s_i \cdot q_{01}^*\omega_i) \\
& =
q_2^*s - \sum q_2^*s_i \cdot q_{12}^*\omega_i -
\sum q_2^*s_i \cdot q_{01}^*\omega_i +
\sum q_{12}^*\nabla(s_i) \cdot q_{01}^*\omega_i
\end{align*}
et, d'autre part,
$$
q_{02}^*c(q_0^*s) = q_2^*s - \sum q_2^*s_i \cdot q_{02}^*\omega_i.
$$
Les formules de la remarque \ref{remark-second-order-thickening} montrent
que
$q_{01}^*\omega_i + q_{12}^*\omega_i - q_{02}^*\omega_i = \text{d}\omega_i$.
La différence des deux expressions ci-dessus est donc
$$
\sum q_2^*s_i \cdot \text{d}\omega_i -
\sum q_{12}^*\nabla(s_i) \cdot q_{01}^*\omega_i
$$
Notons que
$q_{12}^*\omega \cdot q_{01}^*\omega' = \omega' \wedge \omega =
- \omega \wedge \omega'$ par définition de la multiplication sur
$\mathcal{O}_{T''}$. L'expression ci-dessus est donc $\nabla^2(s)$, considérée
comme une section du sous-faisceau $\mathcal{F}_T \otimes \Omega^2_{T/S, \delta}$ de
$q_2^*\mathcal{F}$. On obtient ainsi la condition d'intégrabilité.
\end{proof}




\section{Algèbre cosimpliciale}
\label{section-cosimplicial}

\noindent
Cette section devrait être déplacée ailleurs. Un
{\it anneau cosimplicial} est un objet cosimplicial
de la catégorie des anneaux. Étant donné un anneau $R$, une
{\it $R$-algèbre cosimpliciale} est un objet cosimplicial de la
catégorie des $R$-algèbres. Un {\it idéal cosimplicial} d'un anneau
cosimplicial $A_*$ est la donnée d'idéaux $I_n \subset A_n$, pour tout $n$, tels
que $A(f)(I_n) \subset I_m$ pour tout $f : [n] \to [m]$ dans $\Delta$.

\medskip\noindent
Soit $A_*$ un anneau cosimplicial. Soit $\mathcal{C}$ la catégorie
des couples $(A, M)$ où $A$ est un anneau et $M$ un module sur $A$.
Un morphisme $(A, M) \to (A', M')$ consiste en un homomorphisme d'anneaux $A \to A'$ et
un homomorphisme $A$-linéaire $M \to M'$, où $M'$ est considéré comme un $A$-module
par l'intermédiaire de $A \to A'$ et de la structure de $A'$-module sur $M'$. Cela étant,
on définit un {\it module cosimplicial $M_*$ sur $A_*$} comme un objet
cosimplicial $(A_*, M_*)$ de $\mathcal{C}$ dont la première composante est $A_*$.
Un {\it homomorphisme $\varphi_* : M_* \to N_*$ de modules cosimpliciaux sur
$A_*$} est un morphisme $(A_*, M_*) \to (A_*, N_*)$ d'objets cosimpliciaux
de $\mathcal{C}$ dont la première composante est $1_{A_*}$.

\medskip\noindent
Une {\it homotopie} entre des homomorphismes $\varphi_*, \psi_* : M_* \to N_*$
de modules cosimpliciaux sur $A_*$ est une homotopie entre les
morphismes associés $(A_*, M_*) \to (A_*, N_*)$ dont la première composante est
l'homotopie triviale (duale de
Objets simpliciaux, exemple \ref{simplicial-example-trivial-homotopy}).
Explicitons cette définition. Une telle homotopie est une homotopie
$$
h : M_* \longrightarrow \Hom(\Delta[1], N_*)
$$
entre $\varphi_*$ et $\psi_*$ comme homomorphismes de groupes abéliens
cosimpliciaux, telle que, pour chaque $n$, l'application
$h_n : M_n \to \prod_{\alpha \in \Delta[1]_n} N_n$ soit $A_n$-linéaire.
Le lemme suivant est, pour les modules cosimpliciaux,
une version du lemme \ref{simplicial-lemma-functorial-homotopy}
d'Objets simpliciaux.

\begin{lemma}
\label{lemma-homotopy-tensor}
Soit $A_*$ un anneau cosimplicial. Soient $\varphi_*, \psi_* : K_* \to M_*$
des homomorphismes de modules cosimpliciaux sur $A_*$.
\begin{enumerate}
\item
\label{item-tensor}
Si $\varphi_*$ et $\psi_*$ sont homotopes, alors
$$
\varphi_* \otimes 1, \psi_* \otimes 1 :
K_* \otimes_{A_*} L_* \longrightarrow M_* \otimes_{A_*} L_*
$$
sont homotopes pour tout $A_*$-module cosimplicial $L_*$.
\item
\label{item-wedge}
Si $\varphi_*$ et $\psi_*$ sont homotopes, alors
$$
\wedge^i(\varphi_*), \wedge^i(\psi_*) :
\wedge^i(K_*) \longrightarrow \wedge^i(M_*)
$$
sont homotopes.
\item
\label{item-base-change}
Si $\varphi_*$ et $\psi_*$ sont homotopes et si $A_* \to B_*$
est un homomorphisme d'anneaux cosimpliciaux, alors
$$
\varphi_* \otimes 1, \psi_* \otimes 1 :
K_* \otimes_{A_*} B_* \longrightarrow M_* \otimes_{A_*} B_*
$$
sont homotopes comme homomorphismes de modules cosimpliciaux sur $B_*$.
\item
\label{item-completion}
Si $I_* \subset A_*$ est un idéal cosimplicial, alors les applications
induites
$$
\varphi^\wedge_*, \psi^\wedge_* :
K_*^\wedge \longrightarrow M_*^\wedge
$$
entre les complétés sont homotopes.
\item Ajouter ici d'autres assertions au besoin, par exemple sur les puissances symétriques.
\end{enumerate}
\end{lemma}

\begin{proof}
Soit $h : K_* \longrightarrow \Hom(\Delta[1], M_*)$ l'homotopie
donnée. En degré $n$, on a
$$
h_n = (h_{n, \alpha}) :
K_n \longrightarrow
\prod\nolimits_{\alpha \in \Delta[1]_n} M_n
$$
voir Objets simpliciaux, section \ref{simplicial-section-homotopy-cosimplicial}.
Pour qu'une famille de $h_{n, \alpha}$ constitue une homotopie,
il faut et il suffit que, pour tout $f : [n] \to [m]$, on
ait
$$
h_{m, \alpha} \circ M_*(f) = N_*(f) \circ h_{n, \alpha \circ f}
$$
voir
Objets simpliciaux, équation (\ref{simplicial-equation-property-homotopy-cosimplicial}).
On doit aussi avoir $\psi_n = h_{n, 0 : [n] \to [1]}$ et
$\varphi_n = h_{n, 1 : [n] \to [1]}$.

\medskip\noindent
Dans chacun des cas du lemme, on peut construire les morphismes correspondants.
Cas (\ref{item-tensor}). On peut utiliser l'homotopie $h \otimes 1$ définie
en degré $n$ en posant
$$
(h \otimes 1)_{n, \alpha} = h_{n, \alpha} \otimes 1_{L_n} :
K_n \otimes_{A_n} L_n
\longrightarrow
M_n \otimes_{A_n} L_n.
$$
Cas (\ref{item-wedge}). On peut utiliser l'homotopie $\wedge^ih$ définie
en degré $n$ en posant
$$
\wedge^i(h)_{n, \alpha} = \wedge^i(h_{n, \alpha}) :
\wedge^i_{A_n}(K_n)
\longrightarrow
\wedge^i_{A_n}(M_n).
$$
Cas (\ref{item-base-change}). On peut utiliser l'homotopie $h \otimes 1$ définie
en degré $n$ en posant
$$
(h \otimes 1)_{n, \alpha} = h_{n, \alpha} \otimes 1 :
K_n \otimes_{A_n} B_n
\longrightarrow
M_n \otimes_{A_n} B_n.
$$
Cas (\ref{item-completion}). On peut utiliser l'homotopie $h^\wedge$ définie
en degré $n$ en posant
$$
(h^\wedge)_{n, \alpha} = h_{n, \alpha}^\wedge :
K_n^\wedge
\longrightarrow
M_n^\wedge.
$$
Cette construction convient, car chaque $h_{n, \alpha}$ est $A_n$-linéaire.
\end{proof}






\section{Cristaux en modules quasi-cohérents}
\label{section-quasi-coherent-crystals}

\noindent
Dans la situation \ref{situation-affine}.
Posons $X = \Spec(C)$ et $S = \Spec(A)$. Nous allons
classifier les cristaux en modules quasi-cohérents sur $\text{Cris}(X/S)$.
Auparavant, fixons quelques notations.

\medskip\noindent
Choisissons un anneau de polynômes $P = A[x_i]$ sur $A$ et une surjection $P \to C$
de $A$-algèbres, de noyau $J = \Ker(P \to C)$. Posons
\begin{equation}
\label{equation-D}
D = \lim_e D_{P, \gamma}(J) / p^eD_{P, \gamma}(J)
\end{equation}
pour désigner l'enveloppe à puissances divisées complétée $p$-adiquement.
Cet anneau est muni d'un idéal à puissances divisées $\bar J$ et d'une structure de
puissances divisées $\bar \gamma$ ; voir le lemme \ref{lemma-list-properties}.
Posons $D_e = D/p^eD$ et notons $\bar J_e$ l'image de $\bar J$ dans $D_e$.
Nous emploierons l'abréviation
\begin{equation}
\label{equation-omega-D}
\Omega_D = \lim_e \Omega_{D_e/A, \bar\gamma} =
\lim_e \Omega_{D/A, \bar\gamma}/p^e\Omega_{D/A, \bar\gamma}
\end{equation}
pour le complété $p$-adique du module des différentielles à puissances divisées ;
voir le lemme \ref{lemma-differentials-completion}.
C'est aussi le complété $p$-adique de
$\Omega_{D_{P, \gamma}(J)/A, \bar\gamma}$
qui est libre, ayant pour base les $\text{d}x_i$ ; voir le
lemme \ref{lemma-module-differentials-divided-power-envelope}.
Ainsi, tout élément de $\Omega_D$ s'écrit de manière unique comme une somme
$\sum f_i\text{d}x_i$ telle que, pour tout $e$, seuls un nombre fini de $f_i$
n'appartiennent pas à $p^eD$. De plus, les applications
$\text{d}_{D_e/A, \bar\gamma} : D_e \to \Omega_{D_e/A, \bar\gamma}$
sont compatibles et définissent une $A$-dérivation à puissances divisées
\begin{equation}
\label{equation-derivation-D}
\text{d} : D \longrightarrow \Omega_D
\end{equation}
après complétion $p$-adique.

\medskip\noindent
Nous aurons également besoin des « produits $\Spec(D(n))$ de $\Spec(D)$ » ; voir la
proposition \ref{proposition-compute-cohomology} et sa démonstration pour une
explication. Formellement, ils sont définis comme suit. Pour $n \geq 0$, soit
$J(n) = \Ker(P \otimes_A \ldots \otimes_A P \to C)$, où
le produit tensoriel comporte $n + 1$ facteurs. On pose
\begin{equation}
\label{equation-Dn}
D(n) = \lim_e
D_{P \otimes_A \ldots \otimes_A P, \gamma}(J(n))/
p^eD_{P \otimes_A \ldots \otimes_A P, \gamma}(J(n))
\end{equation}
qui est le complété $p$-adique de l'enveloppe à puissances divisées.
On note $\bar J(n)$ son idéal à puissances divisées et $\bar \gamma(n)$
sa structure de puissances divisées. On introduit aussi $D(n)_e = D(n)/p^eD(n)$ ainsi
que le module des différentielles complété $p$-adiquement
\begin{equation}
\label{equation-omega-Dn}
\Omega_{D(n)} = \lim_e \Omega_{D(n)_e/A, \bar\gamma} =
\lim_e \Omega_{D(n)/A, \bar\gamma}/p^e\Omega_{D(n)/A, \bar\gamma}
\end{equation}
et la dérivation
\begin{equation}
\label{equation-derivation-Dn}
\text{d} : D(n) \longrightarrow \Omega_{D(n)}
\end{equation}
Bien entendu, $D = D(0)$. Notons que les anneaux $D(0), D(1), D(2), \ldots$
forment un objet cosimplicial de la catégorie des anneaux à puissances divisées.

\begin{lemma}
\label{lemma-structure-Dn}
Soient $D$ et $D(n)$ comme dans (\ref{equation-D}) et (\ref{equation-Dn}).
La coprojection $P \to P \otimes_A \ldots \otimes_A P$,
$f \mapsto f \otimes 1 \otimes \ldots \otimes 1$
induit un isomorphisme
\begin{equation}
\label{equation-structure-Dn}
D(n) = \lim_e D\langle \xi_i(j) \rangle/p^eD\langle \xi_i(j) \rangle
\end{equation}
de $D$-algèbres, où
$$
\xi_i(j) = x_i \otimes 1 \otimes \ldots \otimes 1 -
1 \otimes \ldots \otimes 1 \otimes x_i \otimes 1 \otimes \ldots \otimes 1
$$
pour $j = 1, \ldots, n$, le second $x_i$ étant placé dans le facteur numéro $j + 1$
facteur ; rappelons que la construction de $D(n)$ part du
produit tensoriel de $n + 1$ copies de $P$ sur $A$.
\end{lemma}

\begin{proof}
On a
$$
P \otimes_A \ldots \otimes_A P = P[\xi_i(j)]
$$
et $J(n)$ est engendré par $J$ et les éléments $\xi_i(j)$.
Le lemme résulte donc du
lemme \ref{lemma-divided-power-envelope-add-variables}.
\end{proof}

\begin{lemma}
\label{lemma-property-Dn}
Soient $D$ et $D(n)$ comme dans (\ref{equation-D}) et (\ref{equation-Dn}).
Alors $(D, \bar J, \bar\gamma)$ et $(D(n), \bar J(n), \bar\gamma(n))$
sont des objets de $\text{Cris}^\wedge(C/A)$ (voir la
remarque \ref{remark-completed-affine-site}), et
$$
D(n) = \coprod\nolimits_{j = 0, \ldots, n} D
$$
dans $\text{Cris}^\wedge(C/A)$.
\end{lemma}

\begin{proof}
La première assertion est claire. Pour la seconde, si $(B \to C, \delta)$ est un
objet de $\text{Cris}^\wedge(C/A)$, alors on a
$$
\Mor_{\text{Cris}^\wedge(C/A)}(D, B) = 
\Hom_A((P, J), (B, \Ker(B \to C)))
$$
et de même pour $D(n)$ en remplaçant $(P, J)$ par
$(P \otimes_A \ldots \otimes_A P, J(n))$. La propriété relative aux coproduits en résulte,
puisque $P \otimes_A \ldots \otimes_A P$ est un coproduit.
\end{proof}

\noindent
Dans le lemme ci-dessous, nous considérerons des couples $(M, \nabla)$ satisfaisant aux
conditions suivantes :
\begin{enumerate}
\item
\label{item-complete}
$M$ est $p$-adiquement complet en tant que $D$-module,
\item
\label{item-connection}
$\nabla : M \to M \otimes^\wedge_D \Omega_D$ est une connexion, c'est-à-dire
$\nabla(fm) = m \otimes \text{d}f + f\nabla(m)$,
\item
\label{item-integrable}
$\nabla$ est intégrable
(voir la remarque \ref{remark-connection}), et
\item
\label{item-topologically-quasi-nilpotent}
$\nabla$ est {\it topologiquement quasi-nilpotente} : si l'on écrit
$\nabla(m) = \sum \theta_i(m)\text{d}x_i$ pour des opérateurs
$\theta_i : M \to M$, alors, pour tout $m \in M$, il n'existe qu'un nombre fini
de couples $(i, k)$ tels que $\theta_i^k(m) \not \in pM$.
\end{enumerate}
Dans la littérature, les opérateurs $\theta_i$ sont parfois notés
$\nabla_{\partial/\partial x_i}$.
Dans le lemme suivant, nous construisons un foncteur de la catégorie des cristaux en modules
quasi-cohérents sur $\text{Cris}(X/S)$ vers la catégorie de ces couples. Nous montrerons que
ce foncteur est une équivalence dans la
proposition \ref{proposition-crystals-on-affine}.

\begin{lemma}
\label{lemma-crystals-on-affine}
Dans la situation ci-dessus, il existe un foncteur
$$
\begin{matrix}
\text{cristaux en} \\
\mathcal{O}_{X/S}\text{-modules quasi-cohérents sur }\text{Cris}(X/S)
\end{matrix}
\longrightarrow
\begin{matrix}
\text{couples }(M, \nabla)\text{ vérifiant} \\
\text{(\ref{item-complete}), (\ref{item-connection}),
(\ref{item-integrable}) et (\ref{item-topologically-quasi-nilpotent})}
\end{matrix}
$$
\end{lemma}

\begin{proof}
Soit $\mathcal{F}$ un cristal en modules quasi-cohérents sur $X/S$.
Posons $T_e = \Spec(D_e)$, de sorte que $(X, T_e, \bar\gamma)$ est un objet
de $\text{Cris}(X/S)$ pour $e \gg 0$. On dispose de morphismes
$$
(X, T_e, \bar\gamma) \to (X, T_{e + 1}, \bar\gamma) \to \ldots
$$
qui sont des immersions fermées. Posons
$$
M =
\lim_e \Gamma((X, T_e, \bar\gamma), \mathcal{F}) =
\lim_e \Gamma(T_e, \mathcal{F}_{T_e}) = \lim_e M_e
$$
Notons que, puisque $\mathcal{F}$ est localement quasi-cohérent, on a
$\mathcal{F}_{T_e} = \widetilde{M_e}$. Puisque $\mathcal{F}$ est un
cristal, on a $M_e = M_{e + 1}/p^eM_{e + 1}$. On voit donc que
$M_e = M/p^eM$ et que $M$ est complet $p$-adiquement ; voir
Algèbre, lemme \ref{algebra-lemma-limit-complete}.

\medskip\noindent
Par le lemme \ref{lemma-automatic-connection}, on sait que $\mathcal{F}$
est muni d'une connexion intégrable canonique
$\nabla : \mathcal{F} \to \mathcal{F} \otimes \Omega_{X/S}$.
En évaluant cette connexion sur les objets $T_e$ construits ci-dessus,
on obtient une connexion intégrable canonique
$$
\nabla : M \longrightarrow M \otimes^\wedge_D \Omega_D
$$
Pour voir qu'elle est topologiquement quasi-nilpotente, explicitons ce que cela signifie.

\medskip\noindent
On peut maintenant appliquer le même procédé aux anneaux $D(n)$.
On obtient ainsi un module $p$-adiquement complet sur $D(n)$, noté $M(n)$. En utilisant encore
la propriété de cristal de $\mathcal{F}$, on obtient des isomorphismes
$$
M \otimes^\wedge_{D, p_0} D(1) \rightarrow M(1)
\leftarrow M \otimes^\wedge_{D, p_1} D(1)
$$
voir la démonstration du lemme \ref{lemma-automatic-connection}.
Notons $c$ le composé de gauche à droite. Prenons $m \in M$.
Posons $\xi_i = x_i \otimes 1 - 1 \otimes x_i$.
D'après (\ref{equation-structure-Dn}), on peut écrire de manière unique
$$
c(m \otimes 1) = \sum\nolimits_K \theta_K(m) \otimes \prod \xi_i^{[k_i]}
$$
avec $\theta_K(m) \in M$, où la somme porte sur les multi-indices
$K = (k_i)$ tels que $k_i \geq 0$ et $\sum k_i < \infty$. Posons
$\theta_i = \theta_K$, où $K$ a un $1$ à la $i$-ième place et
des zéros ailleurs. On a
$$
\nabla(m) = \sum \theta_i(m) \text{d}x_i.
$$
Cela se voit en comparant avec la définition de
$\nabla$. En effet, l'équation qui le définit est
$p_1^*m = \nabla(m) - c(p_0^*m)$ dans le lemme \ref{lemma-automatic-connection},
mais le signe convient parce que, dans le projet Stacks, on utilise systématiquement
$\text{d}f = p_1(f) - p_0(f)$ modulo le carré de l'idéal de la diagonale,
et que $\xi_i = x_i \otimes 1 - 1 \otimes x_i$ s'envoie donc sur $-\text{d}x_i$
modulo le carré de l'idéal de la diagonale.

\medskip\noindent
Notons $q_i : D \to D(2)$ et $q_{ij} : D(1) \to D(2)$ les coprojections
correspondant aux indices $i, j$. Comme dans le dernier paragraphe de la démonstration du
lemme \ref{lemma-automatic-connection},
on voit que
$$
q_{02}^*c = q_{12}^*c \circ q_{01}^*c.
$$
Cela signifie que
$$
\sum\nolimits_{K''} \theta_{K''}(m) \otimes \prod {\zeta''_i}^{[k''_i]}
=
\sum\nolimits_{K', K} \theta_{K'}(\theta_K(m))
\otimes \prod {\zeta'_i}^{[k'_i]} \prod \zeta_i^{[k_i]}
$$
dans $M \otimes^\wedge_{D, q_2} D(2)$, où
\begin{align*}
\zeta_i & = x_i \otimes 1 \otimes 1 - 1 \otimes x_i \otimes 1,\\
\zeta'_i & = 1 \otimes x_i \otimes 1 - 1 \otimes 1 \otimes x_i,\\
\zeta''_i & = x_i \otimes 1 \otimes 1 - 1 \otimes 1 \otimes x_i.
\end{align*}
En particulier, $\zeta''_i = \zeta_i + \zeta'_i$, et
$D(2)$ est le complété $p$-adique de l'anneau de polynômes
à puissances divisées en $\zeta_i, \zeta'_i$ sur $q_2(D)$ ; voir le lemme \ref{lemma-structure-Dn}.
En comparant les coefficients dans l'expression ci-dessus, on obtient immédiatement
$\theta_i \circ \theta_j = \theta_j \circ \theta_i$
(ce qui fournit une autre démonstration de l'intégrabilité de $\nabla$) et
$$
\theta_K(m) = (\prod \theta_i^{k_i})(m).
$$
En particulier, puisque la somme exprimant $c(m \otimes 1)$ ci-dessus doit converger
$p$-adiquement, on conclut que, pour tout $i$ et tout $m \in M$, les termes
$\theta_i^k(m)$ non nuls modulo $p$ sont en nombre fini.
\end{proof}

\begin{proposition}
\label{proposition-crystals-on-affine}
Le foncteur
$$
\begin{matrix}
\text{cristaux en} \\
\mathcal{O}_{X/S}\text{-modules quasi-cohérents sur }\text{Cris}(X/S)
\end{matrix}
\longrightarrow
\begin{matrix}
\text{couples }(M, \nabla)\text{ vérifiant} \\
\text{(\ref{item-complete}), (\ref{item-connection}),
(\ref{item-integrable}) et (\ref{item-topologically-quasi-nilpotent})}
\end{matrix}
$$
du lemme \ref{lemma-crystals-on-affine}
est une équivalence de catégories.
\end{proposition}

\begin{proof}
Soit $(M, \nabla)$ donné. Construisons
un cristal en modules quasi-cohérents $\mathcal{F}$.
Écrivons $\nabla(m) = \sum \theta_i(m)\text{d}x_i$.
Alors $\theta_i \circ \theta_j = \theta_j \circ \theta_i$, et l'on
peut poser $\theta_K(m) = (\prod \theta_i^{k_i})(m)$ pour tout multi-indice
$K = (k_i)$ tel que $k_i \geq 0$ et $\sum k_i < \infty$.

\medskip\noindent
Soit $(U, T, \delta)$ un objet quelconque de $\text{Cris}(X/S)$ avec $T$ affine.
Écrivons $T = \Spec(B)$ ; l'idéal de $U \to T$ sera noté $J_B \subset B$.
Par le lemme \ref{lemma-set-generators}, il existe un entier $e$ et un morphisme
$$
f : (U, T, \delta) \longrightarrow (X, T_e, \bar\gamma)
$$
où $T_e = \Spec(D_e)$, comme dans la démonstration du
lemme \ref{lemma-crystals-on-affine}.
Choisissons de tels $e$ et $f$ ; notons encore $f : D \to B$ l'homomorphisme
correspondant de $A$-algèbres à puissances divisées. On définit $\mathcal{F}_T$ comme le
faisceau quasi-cohérent de $\mathcal{O}_T$-modules associé au $B$-module
$$
M \otimes_{D, f} B.
$$
Il faut toutefois montrer que cette définition ne dépend pas du choix de $f$.
Supposons que $g : D \to B$ soit un second morphisme de ce type. Puisque $f$ et $g$
sont des morphismes de $\text{Cris}(X/S)$, on voit que l'image de
$f - g : D \to B$ est contenue dans l'idéal à puissances divisées $J_B$.
Posons $\xi_i = f(x_i) - g(x_i) \in J_B$. Par analogie avec la démonstration
du lemme \ref{lemma-crystals-on-affine}, on définit un isomorphisme
$$
c_{f, g} : M \otimes_{D, f} B \longrightarrow M \otimes_{D, g} B
$$
par la formule
$$
m \otimes 1 \longmapsto 
\sum\nolimits_K \theta_K(m) \otimes \prod \xi_i^{[k_i]}
$$
qui a un sens d'après les remarques ci-dessus et le fait que $\nabla$
est topologiquement quasi-nilpotente (la somme est donc finie !).
Un calcul montre que
$$
c_{g, h} \circ c_{f, g} = c_{f, h}
$$
si l'on se donne un troisième morphisme
$h : (U, T, \delta) \longrightarrow (X, T_e, \bar\gamma)$.
On a aussi $c_{f, f} = 1$.
Ces applications sont donc toutes des isomorphismes, et l'on voit que
le module $\mathcal{F}_T$ ne dépend pas du choix de $f$.

\medskip\noindent
Si $a : (U', T', \delta') \to (U, T, \delta)$ est un morphisme entre objets affines
de $\text{Cris}(X/S)$, alors, en choisissant $f' = f \circ a$, il est clair
qu'il existe un isomorphisme canonique
$a^*\mathcal{F}_T \to \mathcal{F}_{T'}$. On omet de vérifier que cette
application ne dépend pas du choix de $f$. En prenant ces applications comme applications de restriction,
il est clair que l'on obtient un cristal en modules quasi-cohérents
sur la sous-catégorie pleine de $\text{Cris}(X/S)$ formée des objets affines.
On omet de montrer que celui-ci s'étend en un cristal sur tout
$\text{Cris}(X/S)$. On omet aussi de montrer que ce procédé est un foncteur
et qu'il est quasi-inverse du foncteur construit dans le
lemme \ref{lemma-crystals-on-affine}.
\end{proof}

\begin{lemma}
\label{lemma-crystals-on-affine-smooth}
Dans la situation \ref{situation-affine}.
Soient $A \to P' \to C$ des homomorphismes d'anneaux tels que $A \to P'$ soit lisse et $P' \to C$
surjectif de noyau $J'$. Soit $D'$ le complété $p$-adique de
$D_{P', \gamma}(J')$. Alors il existe un choix de données $A \to P \to C$
comme ci-dessus et des homomorphismes de $A$-algèbres à puissances divisées
$$
a : D \longrightarrow D',\quad b : D' \longrightarrow D
$$
compatibles avec les applications $D \to C$ et $D' \to C$ tels que
$a \circ b = \text{id}_{D'}$. Ces applications induisent
une équivalence entre la catégorie des couples $(M, \nabla)$ vérifiant
(\ref{item-complete}), (\ref{item-connection}),
(\ref{item-integrable}) et (\ref{item-topologically-quasi-nilpotent})
sur $D$ et celle des couples $(M', \nabla')$ vérifiant
(\ref{item-complete}), (\ref{item-connection}),
(\ref{item-integrable}) et
(\ref{item-topologically-quasi-nilpotent})\footnote{Cette condition
est délicate à formuler pour $(M', \nabla')$ sur $D'$. Voir la démonstration.} sur $D'$.
En particulier, l'équivalence de catégories de la
proposition \ref{proposition-crystals-on-affine}
vaut aussi pour le foncteur correspondant à valeurs dans la catégorie des couples sur $D'$.
\end{lemma}

\begin{proof}
Choisissons une algèbre polynomiale $P = A[y_1, \ldots, y_m]$ sur $A$ et une
surjection $P \to P'$. Définissons $P \to C$ comme le composé $P \to P' \to C$.
On obtient un homomorphisme surjectif $a : D \to D'$ par fonctorialité des
enveloppes à puissances divisées et de la complétion. Prenons $e$ assez grand pour que
$D_e$ soit un épaississement à puissances
divisées de $C$ sur $A$. Alors $D_e \to C$ est une surjection dont le noyau
est localement nilpotent ; voir Algèbre à puissances divisées, lemme \ref{dpa-lemma-nil}.
En posant $D'_e = D'/p^eD'$,
on voit que le noyau de $D_e \to D'_e$ est localement nilpotent.
Par conséquent, d'après Algèbre, lemme \ref{algebra-lemma-smooth-strong-lift},
on peut trouver un relèvement $\beta_e : P' \to D_e$ de l'application $P' \to D'_e$.
Notons que $D_{e + i + 1} \to D_{e + i} \times_{D'_{e + i}} D'_{e + i + 1}$
est surjectif de noyau de carré nul pour tout $i \geq 0$, puisque
$p^{e + i}D \to p^{e + i}D'$ est surjectif. En appliquant successivement la propriété
usuelle de relèvement (Algèbre, proposition \ref{algebra-proposition-smooth-formally-smooth})
aux diagrammes
$$
\xymatrix{
P' \ar[r] & D_{e + i} \times_{D'_{e + i}} D'_{e + i + 1} \\
A \ar[u] \ar[r] & D_{e + i + 1} \ar[u]
}
$$
on voit qu'il existe un homomorphisme de $A$-algèbres $\beta : P' \to D$ dont
le composé avec $a$ est l'application donnée $P' \to D'$.
Par la propriété universelle de l'enveloppe à puissances divisées, on obtient un
homomorphisme $D_{P', \gamma}(J') \to D$. Comme $D$ est $p$-adiquement complet, on
obtient $b : D' \to D$ tel que $a \circ b = \text{id}_{D'}$.

\medskip\noindent
Considérons les foncteurs de changement de base
$$
F : (M, \nabla) \longmapsto
(M \otimes^\wedge_{D, a} D', \nabla')
\quad\text{et}\quad
G : (M', \nabla') \longmapsto
(M' \otimes^\wedge_{D', b} D, \nabla)
$$
sur les modules à connexion vérifiant (\ref{item-complete}),
(\ref{item-connection}) et (\ref{item-integrable}).
Voir la remarque \ref{remark-base-change-connection}.
Puisque $a \circ b = \text{id}_{D'}$, on voit que
$F \circ G$ est le foncteur identité. Disons que $(M', \nabla')$
possède la propriété (\ref{item-topologically-quasi-nilpotent}) si elle
est satisfaite par $G(M', \nabla')$. Un argument formel montre alors que, pour achever
la démonstration, il suffit de montrer que $G(F(M, \nabla))$ est isomorphe
à $(M, \nabla)$ lorsque $(M, \nabla)$ satisfait les quatre
conditions (\ref{item-complete}), (\ref{item-connection}),
(\ref{item-integrable}) et (\ref{item-topologically-quasi-nilpotent}).
Pour cela, on utilise l'isomorphisme fonctoriel
$$
c_{\text{id}_D, b \circ a} :
M \otimes_{D, \text{id}_D} D
\longrightarrow
M \otimes_{D, b \circ a} D
$$
de la démonstration de la proposition \ref{proposition-crystals-on-affine}
(qui requiert la quasi-nilpotence topologique de $\nabla$,
supposée ici).
Il reste à montrer que ce morphisme est horizontal, c'est-à-dire
compatible avec les connexions, ce que l'on omet.

\medskip\noindent
La dernière assertion de l'énoncé en résulte.
\end{proof}

\begin{remark}
\label{remark-equivalence-more-general}
L'équivalence de la proposition \ref{proposition-crystals-on-affine}
vaut si l'on part d'une surjection $P \to C$ telle que $P/A$ satisfasse à la
propriété forte de relèvement du
lemme \ref{algebra-lemma-smooth-strong-lift} d'Algèbre.
Pour le montrer, on peut raisonner comme dans la démonstration du
lemme \ref{lemma-crystals-on-affine-smooth}.
(On ajoutera ici les détails si le besoin s'en fait sentir.)
Il existe vraisemblablement aussi une démonstration directe de ce résultat, mais l'avantage
d'utiliser des anneaux de polynômes est que les anneaux $D(n)$ sont les complétés $p$-adiques
d'anneaux de polynômes à puissances divisées, ce qui simplifie l'algèbre.
\end{remark}



\section{Remarques générales sur la cohomologie}
\label{section-cohomology-lqc}

\noindent
Dans cette section, on effectue quelques préparatifs pour ramener la cohomologie
des modules sur le site cristallin d'un schéma affine à
une question algébrique.

\begin{lemma}
\label{lemma-vanishing-lqc}
Dans la situation \ref{situation-global}.
Soit $\mathcal{F}$ un $\mathcal{O}_{X/S}$-module localement quasi-cohérent
sur $\text{Cris}(X/S)$. Alors on a
$$
H^p((U, T, \delta), \mathcal{F}) = 0
$$
pour tout $p > 0$ et tout $(U, T, \delta)$ tel que $T$ ou $U$ soit affine.
\end{lemma}

\begin{proof}
Comme $U \to T$ est un épaississement, on voit que $U$ est affine si et seulement si $T$
est affine ; voir Limites, lemme \ref{limits-lemma-affine}.
Cela étant, appliquons le
lemme \ref{sites-cohomology-lemma-cech-vanish-collection} de Cohomologie sur les sites
à la collection $\mathcal{B}$ des objets affines $(U, T, \delta)$ et à la
collection $\text{Cov}$ des recouvrements ouverts affines
$\mathcal{U} = \{(U_i, T_i, \delta_i) \to (U, T, \delta)\}$. Le
complexe de {\v C}ech
${\check C}^*(\mathcal{U}, \mathcal{F})$ associé à un tel recouvrement est simplement
le complexe de {\v C}ech du $\mathcal{O}_T$-module quasi-cohérent
$\mathcal{F}_T$
(on utilise ici l'hypothèse selon laquelle $\mathcal{F}$ est localement quasi-cohérent)
relativement au recouvrement ouvert affine $\{T_i \to T\}$ du
schéma affine $T$. La cohomologie de {\v C}ech est donc nulle par
Cohomologie des schémas, lemmes
\ref{coherent-lemma-cech-cohomology-quasi-coherent} et
\ref{coherent-lemma-quasi-coherent-affine-cohomology-zero}.
Ainsi, les hypothèses du
lemme \ref{sites-cohomology-lemma-cech-vanish-collection} de Cohomologie sur les sites
sont satisfaites, ce qui conclut.
\end{proof}

\begin{lemma}
\label{lemma-compare}
Dans la situation \ref{situation-global}.
Supposons en outre que $X$ et $S$ soient des schémas affines.
Considérons la sous-catégorie pleine $\mathcal{C} \subset \text{Cris}(X/S)$
formée des épaississements à puissances divisées $(X, T, \delta)$
et munie de la topologie chaotique (voir
Sites, exemple \ref{sites-example-indiscrete}).
Pour tout $\mathcal{O}_{X/S}$-module localement quasi-cohérent $\mathcal{F}$,
on a
$$
R\Gamma(\mathcal{C}, \mathcal{F}|_\mathcal{C}) =
R\Gamma(\text{Cris}(X/S), \mathcal{F})
$$
\end{lemma}

\begin{proof}
Notons $\text{AffineCris}(X/S)$ la sous-catégorie pleine de $\text{Cris}(X/S)$
formée des objets $(U, T, \delta)$ tels que $U$ et $T$ soient affines.
On en fait un site en déclarant qu'une famille de morphismes
$\{(U_i, T_i, \delta_i) \to (U, T, \delta)\}_{i \in I}$ de
$\text{AffineCris}(X/S)$ est un recouvrement si et seulement si elle en est un
dans $\text{Cris}(X/S)$. Avec cette définition, le foncteur d'inclusion
$$
\text{AffineCris}(X/S) \longrightarrow \text{Cris}(X/S)
$$
est un foncteur cocontinu spécial au sens de
Sites, définition \ref{sites-definition-special-cocontinuous-functor}.
La démonstration est exactement la même que celle du
lemme \ref{topologies-lemma-affine-big-site-Zariski} de Topologies.
On voit ainsi que le topos des faisceaux sur $\text{Cris}(X/S)$
est le même que celui des faisceaux sur $\text{AffineCris}(X/S)$
par restriction le long du foncteur d'inclusion affiché.
Il faut donc montrer l'énoncé correspondant pour
l'inclusion $\mathcal{C} \subset \text{AffineCris}(X/S)$.

\medskip\noindent
On utilisera sans autre mention le fait que $\mathcal{C}$ et
$\text{AffineCris}(X/S)$ admettent des produits et des produits fibrés
(on omet les détails ; voir le
lemme \ref{lemma-divided-power-thickening-fibre-products}).
Le foncteur d'inclusion $u : \mathcal{C} \to \text{AffineCris}(X/S)$
est pleinement fidèle, continu et commute aux produits et aux produits fibrés.
On affirme qu'il définit un morphisme de sites annelés
$$
f :
(\text{AffineCris}(X/S), \mathcal{O}_{X/S})
\longrightarrow
(\Sh(\mathcal{C}), \mathcal{O}_{X/S}|_\mathcal{C})
$$
Pour le voir, on utilisera Sites, lemme \ref{sites-lemma-directed-morphism}.
Notons que $\mathcal{C}$ admet des produits fibrés et que $u$ y commute ;
les catégories $\mathcal{I}^u_{(U, T, \delta)}$ sont donc des réunions disjointes
de catégories filtrantes (par Sites, lemme \ref{sites-lemma-almost-directed}, et
Catégories, lemme \ref{categories-lemma-split-into-directed}). Il
suffit donc de montrer que $\mathcal{I}^u_{(U, T, \delta)}$ est connexe.
Le fait qu'elle soit non vide résulte du lemme \ref{lemma-set-generators} : puisque $U$ et $T$
sont affines, ce lemme affirme qu'il existe au moins un objet
$(X, T', \delta')$ de $\mathcal{C}$ et un morphisme
$(U, T, \delta) \to (X, T', \delta')$ d'épaississements à puissances divisées.
La connexité résulte du fait que $\mathcal{C}$ admet des produits
et que $u$ y commute (comparer avec la démonstration du
lemme \ref{sites-lemma-directed} de Sites).

\medskip\noindent
Notons que $f_*\mathcal{F} = \mathcal{F}|_\mathcal{C}$. Pour démontrer le lemme, il suffit donc
d'établir que $R^pf_*\mathcal{F} = 0$ pour $p > 0$ ; voir
Cohomologie sur les sites, lemme \ref{sites-cohomology-lemma-apply-Leray}. Par
Cohomologie sur les sites, lemme \ref{sites-cohomology-lemma-higher-direct-images},
il suffit de montrer que
$H^p(\text{AffineCris}(X/S)/(X, T, \delta), \mathcal{F}) = 0$
pour tout $(X, T, \delta)$.
Cela résulte du lemme \ref{lemma-vanishing-lqc}, car le
topos du site $\text{AffineCris}(X/S)/(X, T, \delta)$
est équivalent au topos du site
$\text{Cris}(X/S)/(X, T, \delta)$ utilisé dans ce lemme.
\end{proof}

\begin{lemma}
\label{lemma-complete}
Dans la situation \ref{situation-affine}.
Posons $\mathcal{C} = (\text{Cris}(C/A))^{opp}$ et
$\mathcal{C}^\wedge = (\text{Cris}^\wedge(C/A))^{opp}$
munis de la topologie chaotique ; pour les notations, voir la
remarque \ref{remark-completed-affine-site}.
Il existe un morphisme de topos
$$
g : \Sh(\mathcal{C}) \longrightarrow \Sh(\mathcal{C}^\wedge)
$$
tel que, si $\mathcal{F}$ est un faisceau de groupes abéliens sur
$\mathcal{C}$, alors
$$
R^pg_*\mathcal{F}(B \to C, \delta) =
\left\{
\begin{matrix}
\lim_e \mathcal{F}(B_e \to C, \delta) & \text{si }p = 0 \\
R^1\lim_e \mathcal{F}(B_e \to C, \delta) & \text{si }p = 1 \\
0 & \text{sinon}
\end{matrix}
\right.
$$
où $B_e = B/p^eB$ pour $e \gg 0$.
\end{lemma}

\begin{proof}
Tout foncteur entre catégories définit, dans le même sens, un morphisme de
topos chaotiques ; en effet, un tel foncteur
peut être considéré comme un foncteur cocontinu entre sites ; voir
Sites, section \ref{sites-section-cocontinuous-morphism-topoi}.
On omet la démonstration de la description de $g_*\mathcal{F}$.
Notons que, dans l'énoncé, $(B_e \to C, \delta)$
n'est un objet de $\text{Cris}(C/A)$ que pour $e$ assez grand.
Soit $\mathcal{I}$ un faisceau abélien injectif sur $\mathcal{C}$.
Alors les applications de transition
$$
\mathcal{I}(B_e \to C, \delta) \leftarrow
\mathcal{I}(B_{e + 1} \to C, \delta)
$$
sont surjectives, car les morphismes
$$
(B_e \to C, \delta)
\longrightarrow
(B_{e + 1} \to C, \delta)
$$
sont des monomorphismes dans la catégorie $\mathcal{C}$. Ainsi, pour un faisceau
abélien injectif, les deux membres de la formule affichée dans le lemme coïncident.
En prenant une résolution injective de $\mathcal{F}$, on obtient aisément
le résultat (les faisceaux sont des préfaisceaux, de sorte que l'exactitude se mesure au
niveau des groupes de sections sur les objets).
\end{proof}

\begin{lemma}
\label{lemma-category-with-covering}
Soit $\mathcal{C}$ une catégorie munie de la topologie chaotique.
Soit $X$ un objet de $\mathcal{C}$ tel que tout objet de
$\mathcal{C}$ admette un morphisme vers $X$. Supposons que $\mathcal{C}$
admette les produits de deux objets.
Alors, pour tout faisceau abélien $\mathcal{F}$ sur $\mathcal{C}$,
la cohomologie totale $R\Gamma(\mathcal{C}, \mathcal{F})$ est représentée
par le complexe
$$
\mathcal{F}(X) \to \mathcal{F}(X \times X) \to
\mathcal{F}(X \times X \times X) \to \ldots
$$
associé au groupe abélien cosimplicial $[n] \mapsto \mathcal{F}(X^{n + 1})$.
\end{lemma}

\begin{proof}
Notons que $H^q(X^p, \mathcal{F}) = 0$ pour tout $q > 0$, puisque tout préfaisceau est un
faisceau sur $\mathcal{C}$. L'hypothèse sur $X$ signifie que $h_X \to *$
est surjectif. En utilisant $H^q(X, \mathcal{F}) = H^q(h_X, \mathcal{F})$ et
$H^q(\mathcal{C}, \mathcal{F}) = H^q(*, \mathcal{F})$, on voit que notre
énoncé résulte de
Cohomologie sur les sites,
lemme \ref{sites-cohomology-lemma-cech-to-cohomology-sheaf-sets}.
\end{proof}









\section{Préparatifs cosimpliciaux}
\label{section-cohomology}

\noindent
Dans cette section, on compare la cohomologie cristalline à la cohomologie
de de Rham. On suit \cite{Bhatt}.

\begin{example}
\label{example-cosimplicial-module}
Supposons que $A_*$ soit un anneau cosimplicial quelconque.
Considérons le module cosimplicial $M_*$ défini par la règle
$$
M_n = \bigoplus\nolimits_{i = 0, ..., n} A_n e_i
$$
Pour une application $f : [n] \to [m]$, définissons $M_*(f) : M_n \to M_m$
comme l'unique application $A_*(f)$-linéaire envoyant $e_i$ sur $e_{f(i)}$.
On affirme que l'identité de $M_*$ est homotope à $0$.
En effet, une homotopie est donnée par un morphisme de modules cosimpliciaux
$$
h : M_* \longrightarrow \Hom(\Delta[1], M_*)
$$
voir la section \ref{section-cosimplicial}.
Pour $j \in \{0, \ldots, n + 1\}$, notons $\alpha^n_j : [n] \to [1]$ l'application
définie par $\alpha^n_j(i) = 0 \Leftrightarrow i < j$. Alors
$\Delta[1]_n = \{\alpha^n_0, \ldots, \alpha^n_{n + 1}\}$ et, par conséquent,
$\Hom(\Delta[1], M_*)_n = \prod_{j = 0, \ldots, n + 1} M_n$ ; voir
Objets simpliciaux, sections \ref{simplicial-section-homotopy} et
\ref{simplicial-section-homotopy-cosimplicial}. Plutôt que d'utiliser
cette écriture comme produit, on considère un élément
de $\Hom(\Delta[1], M_*)_n$ comme une fonction $\Delta[1]_n \to M_n$.
Avec cette notation, on définit $h$ en degré $n$ par la règle
$$
h_n(e_i)(\alpha^n_j) =
\left\{
\begin{matrix}
e_{i} & \text{si} & i < j \\
0 & \text{sinon} 
\end{matrix}
\right.
$$
Vérifions d'abord que $h$ est un morphisme de modules cosimpliciaux. Pour
$f : [n] \to [m]$, on va montrer que
\begin{equation}
\label{equation-cosimplicial-morphism}
h_m \circ M_*(f) = \Hom(\Delta[1], M_*)(f) \circ h_n
\end{equation}
Le membre de gauche de (\ref{equation-cosimplicial-morphism}), évalué en
$e_i$, puis en $\alpha^m_j$, vaut
$$
h_m(e_{f(i)})(\alpha^m_j) =
\left\{
\begin{matrix}
e_{f(i)} & \text{si} & f(i) < j \\
0 & \text{sinon}
\end{matrix}
\right.
$$
Notons que $\alpha^m_j \circ f = \alpha^n_{j'}$, où
$0 \leq j' \leq n + 1$ est l'unique indice tel que $f(i) < j$
si et seulement si $i < j'$. Ainsi, le membre de droite de
(\ref{equation-cosimplicial-morphism}), évalué en $e_i$,
puis en $\alpha^m_j$, vaut
$$
M_*(f)(h_n(e_i)(\alpha^m_j \circ f)) =
M_*(f)(h_n(e_i)(\alpha^n_{j'})) =
\left\{
\begin{matrix}
e_{f(i)} & \text{si} & i < j' \\
0 & \text{sinon} 
\end{matrix}
\right.
$$
Notre description de $j'$ montre que les deux résultats sont égaux.
Ainsi, $h$ est un morphisme de modules cosimpliciaux.
Soient $0 : \Delta[0] \to \Delta[1]$ et
$1 : \Delta[0] \to \Delta[1]$ les applications évidentes, et notons
$ev_0, ev_1 : \Hom(\Delta[1], M_*) \to M_*$ les applications
d'évaluation correspondantes. Le lecteur vérifie aisément que
les composés
$$
ev_0 \circ h, ev_1 \circ h : M_* \longrightarrow M_*
$$
sont respectivement $1$ et $0$ ; ainsi, $h$ est l'homotopie cherchée entre
$1$ et $0$.
\end{example}

\begin{lemma}
\label{lemma-vanishing-omega-1}
Avec les notations de (\ref{equation-omega-Dn}), le complexe
$$
\Omega_{D(0)} \to \Omega_{D(1)} \to \Omega_{D(2)} \to \ldots
$$
est homotope à zéro comme $D(*)$-module cosimplicial.
\end{lemma}

\begin{proof}
On va utiliser le principe de
Objets simpliciaux, lemme \ref{simplicial-lemma-functorial-homotopy},
et, plus précisément, le
lemme \ref{lemma-homotopy-tensor},
qui affirme que tout foncteur transforme des applications homotopes entre objets
(co)simpliciaux en applications homotopes.
Le complexe du lemme est égal au complété $p$-adique du
changement de base du module cosimplicial
$$
M_* = \left(
\Omega_{P/A} \to
\Omega_{P \otimes_A P/A} \to
\Omega_{P \otimes_A P \otimes_A P/A} \to \ldots
\right)
$$
par l'homomorphisme d'anneaux cosimpliciaux $P\otimes_A \ldots \otimes_A P \to D(n)$. Cela
résulte du lemme \ref{lemma-module-differentials-divided-power-envelope} ;
voir les commentaires après (\ref{equation-omega-D}). Il
suffit donc de montrer que le module cosimplicial $M_*$ est homotope à zéro
(on utilise le changement de base et la complétion $p$-adique).
On peut même supposer que $A = \mathbf{Z}$ et $P = \mathbf{Z}[\{x_i\}_{i \in I}]$,
car on peut effectuer le changement de base par $\mathbf{Z} \to A$.
Dans ce cas, $P^{\otimes n + 1}$ est l'algèbre polynomiale engendrée par les éléments
$$
x_i(e) = 1 \otimes \ldots \otimes x_i \otimes \ldots \otimes 1
$$
où $x_i$ occupe la $e$-ième place. Les modules du complexe sont libres, ayant
pour base les $\text{d}x_i(e)$. Notons que, si $f : [n] \to [m]$ est une
application, alors
$$
M_*(f)(\text{d}x_i(e)) = \text{d}x_i(f(e))
$$
On voit donc que $M_*$ est une somme directe, indexée par $I$, de copies du module
étudié dans l'exemple \ref{example-cosimplicial-module}, ce qui conclut.
\end{proof}

\begin{lemma}
\label{lemma-vanishing}
Avec les notations de (\ref{equation-Dn}) et (\ref{equation-omega-Dn}),
pour tout module cosimplicial $M_*$ sur $D(*)$ et tout
$i > 0$, le module cosimplicial
$$
M_0 \otimes^\wedge_{D(0)} \Omega^i_{D(0)} \to
M_1 \otimes^\wedge_{D(1)} \Omega^i_{D(1)} \to
M_2 \otimes^\wedge_{D(2)} \Omega^i_{D(2)} \to \ldots
$$
est homotope à zéro, où $\Omega^i_{D(n)}$ est le complété $p$-adique
de la $i$-ième puissance extérieure de $\Omega_{D(n)}$.
\end{lemma}

\begin{proof}
Par le lemme \ref{lemma-vanishing-omega-1}, les endomorphismes $0$ et $1$
de $\Omega_{D(*)}$ sont homotopes.
En appliquant le foncteur $\wedge^i$, on voit qu'il
en va de même pour le module cosimplicial $\wedge^i\Omega_{D(*)}$ ; voir le
lemme \ref{lemma-homotopy-tensor}.
Une autre application du même lemme montre que le complété $p$-adique
$\Omega^i_{D(*)}$ est homotopiquement équivalent à zéro.
En tensorisant par $M_*$, on voit que $M_* \otimes_{D(*)} \Omega^i_{D(*)}$
est homotope à zéro ; voir encore le lemme \ref{lemma-homotopy-tensor}.
Une dernière application du foncteur de complétion $p$-adique achève la démonstration.
\end{proof}



\section{Lemme de Poincar\'e à puissances divisées}
\label{section-poincare}

\noindent
Nous ne considérons que la version la plus simple possible.

\begin{lemma}
\label{lemma-poincare}
Soit $A$ un anneau. Soit $P = A\langle x_i \rangle$ un anneau de polynômes
à puissances divisées sur $A$. Pour tout $A$-module $M$, le complexe
$$
0 \to M \to
M \otimes_A P \to
M \otimes_A \Omega^1_{P/A, \delta} \to
M \otimes_A \Omega^2_{P/A, \delta} \to \ldots
$$
est exact. Soit $D$ le complété $p$-adique de $P$.
Soit $\Omega^i_D$ le complété $p$-adique de la $i$-ième puissance
extérieure de $\Omega_{D/A, \delta}$. Pour tout module complet pour la topologie $p$-adique sur
$A$, noté $M$, le complexe
$$
0 \to M \to
M \otimes^\wedge_A D \to
M \otimes^\wedge_A \Omega^1_D \to
M \otimes^\wedge_A \Omega^2_D \to \ldots
$$
est exact.
\end{lemma}

\begin{proof}
Il suffit de montrer que le complexe
$$
E :
(0 \to A \to P \to \Omega^1_{P/A, \delta} \to
\Omega^2_{P/A, \delta} \to \ldots)
$$
est homotopiquement équivalent à zéro comme complexe de $A$-modules.
Pour tout multi-indice $K = (k_i)$, on peut considérer le sous-complexe $E(K)$
qui, en degré $j$, est formé de
$$
\bigoplus\nolimits_{I = \{i_1, \ldots, i_j\} \subset \text{Supp}(K)}
A
\prod\nolimits_{i \not \in I} x_i^{[k_i]}
\prod\nolimits_{i \in I} x_i^{[k_i - 1]}
\text{d}x_{i_1} \wedge \ldots \wedge \text{d}x_{i_j}
$$
Puisque $E = \bigoplus E(K)$, il suffit de prouver que chacun des
complexes $E(K)$ est homotope à zéro. Si $K = 0$, alors
$E(K) : (A \to A)$ est homotope à zéro. Si $K$ a un support (fini) non vide
$S$, alors le complexe $E(K)$ est isomorphe au complexe
$$
0 \to A \to
\bigoplus\nolimits_{s \in S} A \to
\wedge^2(\bigoplus\nolimits_{s \in S} A) \to
\ldots \to \wedge^{\# S}(\bigoplus\nolimits_{s \in S} A) \to 0
$$
qui est homotope à zéro, par exemple par
Compléments d'algèbre, lemme \ref{more-algebra-lemma-homotopy-koszul-abstract}.
\end{proof}

\noindent
Une autre approche, plus directe, du lemme suivant est
expliquée dans l'exemple \ref{example-integrate}.

\begin{lemma}
\label{lemma-relative-poincare}
Soit $A$ un anneau. Soit $(B, I, \delta)$ un anneau à puissances divisées
tel que $B$ soit une $A$-algèbre.
Soit $P = B\langle x_i \rangle$ un anneau de polynômes à puissances divisées
sur $B$, muni comme d'habitude de l'idéal à puissances divisées $J = IP + B\langle x_i \rangle_{+}$.
Soit $M$ un $B$-module muni d'une connexion intégrable
$\nabla : M \to M \otimes_B \Omega^1_{B/A, \delta}$. Alors le morphisme de
complexes de de Rham
$$
M \otimes_B \Omega^*_{B/A, \delta}
\longrightarrow
M \otimes_P \Omega^*_{P/A, \delta}
$$
est un quasi-isomorphisme. Soient $D$, resp.\ $D'$, les complétés $p$-adiques de
$B$, resp.\ $P$, et soient $\Omega^i_D$, resp.\ $\Omega^i_{D'}$, les complétés $p$-adiques
de $\Omega^i_{B/A, \delta}$,
resp.\ $\Omega^i_{P/A, \delta}$. Soit $M$ un module $p$-adiquement complet
sur $D$, muni d'une connexion intégrable
$\nabla : M \to M \otimes^\wedge_D \Omega^1_D$.
Alors le morphisme de complexes de de Rham
$$
M \otimes^\wedge_D \Omega^*_D
\longrightarrow
M \otimes^\wedge_D \Omega^*_{D'}
$$
est un quasi-isomorphisme.
\end{lemma}

\begin{proof}
Considérons la filtration décroissante $F^*$ de $\Omega^*_{B/A, \delta}$
donnée par les sous-complexes
$F^i(\Omega^*_{B/A, \delta}) = \sigma_{\geq i}\Omega^*_{B/A, \delta}$.
Voir Homologie, section \ref{homology-section-truncations}.
Elle induit une filtration décroissante $F^*$ de $\Omega^*_{P/A, \delta}$
en posant
$$
F^i(\Omega^*_{P/A, \delta}) =
F^i(\Omega^*_{B/A, \delta}) \wedge \Omega^*_{P/A, \delta}.
$$
On a une suite exacte courte scindée
$$
0 \to \Omega^1_{B/A, \delta} \otimes_B P \to
\Omega^1_{P/A, \delta} \to
\Omega^1_{P/B, \delta} \to 0
$$
et le dernier module est libre, ayant pour base les $\text{d}x_i$. Il en résulte que
$F^i(\Omega^*_{P/A, \delta}) \to \Omega^*_{P/A, \delta}$ est, en chaque degré, une
injection scindée et que
$$
\text{gr}^i_F(\Omega^*_{P/A, \delta}) =
\Omega^i_{B/A, \delta} \otimes_B \Omega^*_{P/B, \delta}
$$
comme complexes. On peut donc définir une filtration $F^*$ de
$M \otimes_B \Omega^*_{P/A, \delta}$ en posant
$$
F^i(M \otimes_B \Omega^*_{P/A, \delta}) =
M \otimes_B F^i(\Omega^*_{P/A, \delta})
$$
et l'on a
$$
\text{gr}^i_F(M \otimes_B \Omega^*_{P/A, \delta}) =
M \otimes_B \Omega^i_{B/A, \delta} \otimes_B \Omega^*_{P/B, \delta}
$$
comme complexes.
Par le lemme \ref{lemma-poincare}, chacun de ces complexes est
quasi-isomorphe à $M \otimes_B \Omega^i_{B/A, \delta}$ placé en degré $0$.
On voit donc que la première flèche affichée dans l'énoncé est un morphisme de
complexes filtrés qui induit un quasi-isomorphisme sur les gradués associés. Cela
implique qu'il s'agit d'un quasi-isomorphisme, par exemple par la suite spectrale
associée à un complexe filtré ; voir
Homologie, section \ref{homology-section-filtered-complex}.

\medskip\noindent
La démonstration du second quasi-isomorphisme est exactement la même.
\end{proof}



\section{Cohomologie dans le cas affine}
\label{section-cohomology-affine}

\noindent
Revenons à la situation étudiée dans la
section \ref{section-quasi-coherent-crystals}. On
part de $(A, I, \gamma)$ et de $A/I \to C$, et l'on pose
$X = \Spec(C)$ et $S = \Spec(A)$. On choisit ensuite
un anneau de polynômes $P$ sur $A$ et une surjection $P \to C$ de
noyau $J$. On obtient $D$ et $D(n)$ ; voir
(\ref{equation-D}) et (\ref{equation-Dn}).
Posons $T(n)_e = \Spec(D(n)/p^eD(n))$, de sorte que
$(X, T(n)_e, \delta(n))$ soit un objet de $\text{Cris}(X/S)$.
Soit $\mathcal{F}$ un faisceau de $\mathcal{O}_{X/S}$-modules, et posons
$$
M(n) = \lim_e \Gamma((X, T(n)_e, \delta(n)), \mathcal{F})
$$
pour $n = 0, 1, 2, 3, \ldots$. La famille ainsi obtenue forme un module cosimplicial
sur l'anneau cosimplicial $D(0), D(1), D(2), \ldots$.

\begin{proposition}
\label{proposition-compute-cohomology}
Avec les notations ci-dessus, supposons que
\begin{enumerate}
\item $\mathcal{F}$ soit localement quasi-cohérent, et
\item pour tout morphisme $(U, T, \delta) \to (U', T', \delta')$
de $\text{Cris}(X/S)$ tel que $f : T \to T'$ soit une immersion fermée,
le morphisme $c_f : f^*\mathcal{F}_{T'} \to \mathcal{F}_T$ soit surjectif.
\end{enumerate}
Alors le complexe
$$
M(0) \to M(1) \to M(2) \to \ldots
$$
calcule $R\Gamma(\text{Cris}(X/S), \mathcal{F})$.
\end{proposition}

\begin{proof}
L'hypothèse (1) et le lemme \ref{lemma-compare} montrent que
$R\Gamma(\text{Cris}(X/S), \mathcal{F})$ est isomorphe à
$R\Gamma(\mathcal{C}, \mathcal{F})$. Notons que les catégories
$\mathcal{C}$ utilisées dans les lemmes \ref{lemma-compare} et \ref{lemma-complete}
coïncident. Soit $f : T \to T'$ une immersion fermée comme en (2). La surjectivité
de $c_f : f^*\mathcal{F}_{T'} \to \mathcal{F}_T$ équivaut à celle
de $\mathcal{F}_{T'} \to f_*\mathcal{F}_T$. Ainsi, si
$\mathcal{F}$ satisfait (1) et (2), on obtient une suite exacte courte
$$
0 \to \mathcal{K} \to \mathcal{F}_{T'} \to f_*\mathcal{F}_T \to 0
$$
de $\mathcal{O}_{T'}$-modules quasi-cohérents sur $T'$ ; voir
Schémas, section \ref{schemes-section-quasi-coherent}, et en particulier le
lemme \ref{schemes-lemma-push-forward-quasi-coherent}.
Ainsi, si $T'$ est affine, on conclut que l'application de restriction
$\mathcal{F}(U', T', \delta') \to \mathcal{F}(U, T, \delta)$
est surjective par l'annulation de $H^1(T', \mathcal{K})$ ; voir
Cohomologie des schémas, lemme
\ref{coherent-lemma-quasi-coherent-affine-cohomology-zero}.
Les applications de transition des systèmes inverses du lemme \ref{lemma-complete}
sont donc surjectives. On en déduit que
$R^pg_*(\mathcal{F}|_\mathcal{C}) = 0$ pour tout $p \geq 1$,
où $g$ est comme dans le lemme \ref{lemma-complete}.
L'objet $D$ de la catégorie $\mathcal{C}^\wedge$
satisfait l'hypothèse du lemme \ref{lemma-category-with-covering} par le
lemme \ref{lemma-generator-completion},
avec
$$
D \times \ldots \times D = D(n)
$$
dans $\mathcal{C}$, car $D(n)$ est le coproduit de $n + 1$ copies de
$D$ dans $\text{Cris}^\wedge(C/A)$ ; voir le lemme \ref{lemma-property-Dn}.
Cela conclut.
\end{proof}

\begin{lemma}
\label{lemma-cohomology-is-zero}
On se place sous les hypothèses et avec les notations de la
proposition \ref{proposition-compute-cohomology}.
Alors
$$
H^j(\text{Cris}(X/S), \mathcal{F} \otimes_{\mathcal{O}_{X/S}} \Omega^i_{X/S})
= 0
$$
pour tout $i > 0$ et tout $j \geq 0$.
\end{lemma}

\begin{proof}
Il résulte du lemme \ref{lemma-omega-locally-quasi-coherent} que
$\mathcal{H} = \mathcal{F} \otimes_{\mathcal{O}_{X/S}} \Omega^i_{X/S}$
satisfait aussi les hypothèses (1) et (2) de la
proposition \ref{proposition-compute-cohomology}.
Posons $M(n)_e = \Gamma((X, T(n)_e, \delta(n)), \mathcal{F})$,
de sorte que $M(n) = \lim_e M(n)_e$. Alors
\begin{align*}
\lim_e \Gamma((X, T(n)_e, \delta(n)), \mathcal{H}) & =
\lim_e M(n)_e \otimes_{D(n)_e} \Omega^i_{D(n)}/p^e\Omega^i_{D(n)} \\
& = \lim_e M(n)_e \otimes_{D(n)} \Omega^i_{D(n)}
\end{align*}
Par le
lemme \ref{lemma-vanishing},
les modules cosimpliciaux
$$
M(0)_e \otimes_{D(0)} \Omega^i_{D(0)} \to
M(1)_e \otimes_{D(1)} \Omega^i_{D(1)} \to
M(2)_e \otimes_{D(2)} \Omega^i_{D(2)} \to \ldots
$$
sont homotopes à zéro. Comme les applications de transition
$M(n)_{e + 1} \to M(n)_e$ sont surjectives, on voit que
la limite inverse des complexes associés est acyclique\footnote{En fait,
ils sont même homotopes à zéro, car les homotopies sont compatibles, mais cela n'est pas
nécessaire. Cet argument détourné tient au fait que
la limite $\lim_e M(n)_e \otimes_{D(n)} \Omega^i_{D(n)}$ n'est pas le
complété $p$-adique de $M(n) \otimes_{D(n)} \Omega^i_{D(n)}$, car, sous
les hypothèses du lemme, on ne sait pas si
$M(n)_e = M(n)_{e + 1}/p^eM(n)_{e + 1}$. Si $\mathcal{F}$ est un cristal,
cette égalité est vraie.}.
La cohomologie de $\mathcal{H}$ s'annule donc, d'après la
proposition \ref{proposition-compute-cohomology}.
\end{proof}

\begin{proposition}
\label{proposition-compute-cohomology-crystal}
Sous les hypothèses de la proposition \ref{proposition-compute-cohomology},
supposons maintenant que $\mathcal{F}$ soit un cristal en modules quasi-cohérents.
Soit $(M, \nabla)$ le $D$-module à connexion correspondant ; voir la
proposition \ref{proposition-crystals-on-affine}. Alors le complexe
$$
M \otimes^\wedge_D \Omega^*_D
$$
calcule $R\Gamma(\text{Cris}(X/S), \mathcal{F})$.
\end{proposition}

\begin{proof}
On va le démontrer à l'aide des deux suites spectrales associées au
complexe double $K^{*, *}$ de termes
$$
K^{a, b} = M \otimes_D^\wedge \Omega^a_{D(b)}
$$
Que sait-on jusqu'ici ? Le lemme \ref{lemma-vanishing}
affirme que chaque colonne $K^{a, *}$, pour $a > 0$, est acyclique.
La proposition \ref{proposition-compute-cohomology} affirme que
la première colonne $K^{0, *}$ est quasi-isomorphe à
$R\Gamma(\text{Cris}(X/S), \mathcal{F})$.
La première suite spectrale associée au complexe double
montre donc qu'il existe un quasi-isomorphisme canonique de
$R\Gamma(\text{Cris}(X/S), \mathcal{F})$ vers
$\text{Tot}(K^{*, *})$.

\medskip\noindent
Considérons ensuite les lignes $K^{*, b}$. Par le
lemme \ref{lemma-structure-Dn},
chacun des $b + 1$ homomorphismes $D \to D(b)$ présente $D(b)$ comme le complété $p$-adique
d'une algèbre polynomiale à puissances divisées sur $D$.
Le lemme \ref{lemma-relative-poincare} montre donc que le morphisme
$$
M \otimes^\wedge_D\Omega^*_D
\longrightarrow
M \otimes^\wedge_{D(b)} \Omega^*_{D(b)} = K^{*, b}
$$
est un quasi-isomorphisme. Notons que chacun de ces morphismes induit le {\it même}
morphisme en cohomologie (et même le même morphisme dans la catégorie dérivée), car
l'inverse est donné par le morphisme codiagonal $D(b) \to D$ (correspondant
au morphisme de multiplication $P \otimes_A \ldots \otimes_A P \to P$).
Ainsi, en considérant la page $E_1$ de la seconde suite spectrale,
on obtient
$$
E_1^{a, b} = H^a(M \otimes^\wedge_D\Omega^*_D)
$$
avec les différentielles
$$
E_1^{a, 0} \xrightarrow{0}
E_1^{a, 1} \xrightarrow{1}
E_1^{a, 2} \xrightarrow{0}
E_1^{a, 3} \xrightarrow{1} \ldots
$$
car chacune est la somme alternée des identifications données
$H^a(M \otimes^\wedge_D\Omega^*_D) = E_1^{a, 0} = E_1^{a, 1} = \ldots$.
On voit donc que la page $E_2$ est égale à $H^a(M \otimes^\wedge_D\Omega^*_D)$
sur la première ligne et à zéro ailleurs. Il s'ensuit que l'identification
de $M \otimes^\wedge_D\Omega^*_D$ avec la première ligne induit un
quasi-isomorphisme de $M \otimes^\wedge_D\Omega^*_D$ vers
$\text{Tot}(K^{*, *})$.
\end{proof}

\begin{lemma}
\label{lemma-compute-cohomology-crystal-smooth}
On se place sous les hypothèses de la proposition \ref{proposition-compute-cohomology-crystal}.
Soient $A \to P' \to C$ des homomorphismes d'anneaux tels que $A \to P'$ soit lisse et $P' \to C$
surjectif de noyau $J'$. Soit $D'$ le complété $p$-adique de
$D_{P', \gamma}(J')$. Soit $(M', \nabla')$ le couple sur $D'$
correspondant à $\mathcal{F}$ ; voir le
lemme \ref{lemma-crystals-on-affine-smooth}. Alors le complexe
$$
M' \otimes^\wedge_{D'} \Omega^*_{D'}
$$
calcule $R\Gamma(\text{Cris}(X/S), \mathcal{F})$.
\end{lemma}

\begin{proof}
Choisissons $a : D \to D'$ et $b : D' \to D$ comme dans le
lemme \ref{lemma-crystals-on-affine-smooth}.
Notons que le module obtenu par le changement de base $M = M' \otimes^\wedge_{D', b} D$, muni de sa
connexion $\nabla$, correspond à $\mathcal{F}$. Nous savons donc
que $M \otimes^\wedge_D \Omega_D^*$ calcule la cohomologie cristalline
de $\mathcal{F}$, voir la
proposition \ref{proposition-compute-cohomology-crystal}.
Il suffit donc de montrer que les morphismes de changement de base (induits
par $a$ et $b$)
$$
M' \otimes^\wedge_{D'} \Omega^*_{D'}
\longrightarrow
M \otimes^\wedge_D \Omega^*_D
\quad\text{et}\quad
M \otimes^\wedge_D \Omega^*_D
\longrightarrow
M' \otimes^\wedge_{D'} \Omega^*_{D'}
$$
sont des quasi-isomorphismes. Puisque $a \circ b = \text{id}_{D'}$, on voit
que la composée dans un sens est l'identité du complexe
$M' \otimes^\wedge_{D'} \Omega^*_{D'}$. Il suffit donc de montrer que
le morphisme
$$
M \otimes^\wedge_D \Omega^*_D
\longrightarrow
M \otimes^\wedge_D \Omega^*_D
$$
induit par $b \circ a : D \to D$ est un quasi-isomorphisme. (Notons que nous
avons le même complexe des deux côtés puisque $M = M' \otimes^\wedge_{D', b} D$,
d'où $M \otimes^\wedge_{D, b \circ a} D =
M' \otimes^\wedge_{D', b \circ a \circ b} D =
M' \otimes^\wedge_{D', b} D = M$.) En fait, nous affirmons que, pour tout
homomorphisme de $A$-algèbres à puissances divisées $\rho : D \to D$ compatible
avec l'augmentation vers $C$, le morphisme induit
$M \otimes^\wedge_D \Omega^*_D \to M \otimes^\wedge_{D, \rho} \Omega^*_D$
est un quasi-isomorphisme.

\medskip\noindent
Écrivons $\rho(x_i) = x_i + z_i$. Les éléments $z_i$ appartiennent à
l'idéal à puissances divisées de $D$ car $\rho$ est compatible avec
l'augmentation vers $C$. On peut donc factoriser le morphisme $\rho$
en une composée
$$
D \xrightarrow{\sigma} D\langle \xi_i \rangle^\wedge \xrightarrow{\tau} D
$$
où le premier morphisme est donné par $x_i \mapsto x_i + \xi_i$ et le
second est le morphisme de $D$-algèbres à puissances divisées qui envoie $\xi_i$ sur $z_i$.
(Cela utilise les propriétés universelles des algèbres polynomiales, des algèbres
polynomiales à puissances divisées, des enveloppes à puissances divisées et de la complétion $p$-adique.)
Notons qu'il existe un {\it automorphisme} $\alpha$ de
$D\langle \xi_i \rangle^\wedge$ avec $\alpha(x_i) = x_i - \xi_i$
et $\alpha(\xi_i) = \xi_i$. En appliquant le lemme \ref{lemma-relative-poincare}
à $\alpha \circ \sigma$ (qui envoie $x_i$ sur $x_i$) et en utilisant le fait que
$\alpha$ est un isomorphisme, on conclut que $\sigma$ induit un
quasi-isomorphisme de $M \otimes^\wedge_D \Omega^*_D$ vers
$M \otimes^\wedge_{D, \sigma} \Omega^*_{D\langle \xi_i \rangle^\wedge}$.
D'autre part, le morphisme $\tau$ admet pour inverse à gauche le
morphisme $D \to D\langle \xi_i \rangle^\wedge$, $x_i \mapsto x_i$,
et l'on conclut (en utilisant encore une fois le lemme \ref{lemma-relative-poincare})
que $\tau$ induit un quasi-isomorphisme de
$M \otimes^\wedge_{D, \sigma} \Omega^*_{D\langle \xi_i \rangle^\wedge}$
vers $M \otimes^\wedge_{D, \tau \circ \sigma} \Omega^*_D$. En composant ces
deux quasi-isomorphismes, on obtient que $\rho$ induit un quasi-isomorphisme
$M \otimes^\wedge_D \Omega^*_D \to M \otimes^\wedge_{D, \rho} \Omega^*_D$
ce qu'il fallait démontrer.
\end{proof}










\section{Deux contre-exemples}
\label{section-examples}

\noindent
Avant de passer à quelques résultats positifs de la cohomologie cristalline,
donnons deux exemples qui expliquent pourquoi la cohomologie cristalline
fonctionne assez mal lorsque les schémas considérés ne sont pas
propres sur la base ou sont singuliers. Le premier exemple se trouve
dans \cite{BO}.

\begin{example}
\label{example-torsion}
Soit $A = \mathbf{Z}_p$, avec l'idéal à puissances divisées $(p)$ muni de
son unique structure de puissances divisées $\gamma$. Soit
$C = \mathbf{F}_p[x, y]/(x^2, xy, y^2)$. Choisissons la présentation
$$
C = P/J = \mathbf{Z}_p[x, y]/(x^2, xy, y^2, p)
$$
Soit $D = D_{P, \gamma}(J)^\wedge$, avec pour idéal à puissances divisées
$(\bar J, \bar \gamma)$ comme à la section \ref{section-quasi-coherent-crystals}.
Nous noterons encore $x, y$ les images de $x$ et $y$ dans $D$.
Considérons l'élément
$$
\tau = \bar\gamma_p(x^2)\bar\gamma_p(y^2) - \bar\gamma_p(xy)^2 \in D
$$
On remarque que $p\tau = 0$ car
$$
p! \bar\gamma_p(x^2) \bar\gamma_p(y^2) =
x^{2p} \bar\gamma_p(y^2) = \bar\gamma_p(x^2y^2) =
x^py^p \bar\gamma_p(xy) = p! \bar\gamma_p(xy)^2
$$
dans $D$. On remarque aussi que $\text{d}\tau = 0$ dans $\Omega_D$ car
\begin{align*}
\text{d}(\bar\gamma_p(x^2) \bar\gamma_p(y^2))
& =
\bar\gamma_{p - 1}(x^2)\bar\gamma_p(y^2)\text{d}x^2 +
\bar\gamma_p(x^2)\bar\gamma_{p - 1}(y^2)\text{d}y^2 \\
& =
2 x \bar\gamma_{p - 1}(x^2)\bar\gamma_p(y^2)\text{d}x +
2 y \bar\gamma_p(x^2)\bar\gamma_{p - 1}(y^2)\text{d}y \\
& =
2/(p - 1)!( x^{2p - 1} \bar\gamma_p(y^2)\text{d}x +
y^{2p - 1} \bar\gamma_p(x^2)\text{d}y ) \\
& =
2/(p - 1)!
(x^{p - 1} \bar\gamma_p(xy^2)\text{d}x +
y^{p - 1} \bar\gamma_p(x^2y)\text{d}y) \\
& =
2/(p - 1)!
(x^{p - 1}y^p \bar\gamma_p(xy)\text{d}x +
x^py^{p - 1} \bar\gamma_p(xy)\text{d}y) \\
& =
2 \bar\gamma_{p - 1}(xy) \bar\gamma_p(xy)(y\text{d}x + x \text{d}y) \\
& = 
\text{d}(\bar\gamma_p(xy)^2)
\end{align*}
Enfin, nous affirmons que $\tau \not = 0$ dans $D$. Pour le voir, il suffit de
construire un objet $(B \to \mathbf{F}_p[x, y]/(x^2, xy, y^2), \delta)$
de $\text{Cris}(C/S)$ tel que l'image de $\tau$ dans $B$ ne soit pas nulle.
Pour cela, prenons
$$
B = \mathbf{F}_p[x, y, u, v]/(x^3, x^2y, xy^2, y^3, xu, yu, xv, yv, u^2, v^2)
$$
muni de la surjection évidente vers $C$. Soit $K = \Ker(B \to C)$ et
considérons l'application
$$
\delta_p : K \longrightarrow K,\quad
ax^2 + bxy + cy^2 + du + ev + fuv \longmapsto a^pu + c^pv
$$
On vérifie qu'elle satisfait aux hypothèses (1), (2), (3) du
lemme \ref{dpa-lemma-need-only-gamma-p} d'Algèbre à puissances divisées
et définit donc une structure de puissances divisées. De plus,
l'image de $\tau$ est $uv$, qui est non nul dans $B$.
Posons $X = \Spec(C)$ et $S = \Spec(A)$.
Nous en tirons les conclusions suivantes :
\begin{enumerate}
\item $H^0(\text{Cris}(X/S), \mathcal{O}_{X/S})$ a de la $p$-torsion, et
\item l'endomorphisme d'image réciproque par Frobenius $F^* : H^0(\text{Cris}(X/S), \mathcal{O}_{X/S})
\to H^0(\text{Cris}(X/S), \mathcal{O}_{X/S})$ n'est pas injective.
\end{enumerate}
En effet, $\tau$ définit un élément de torsion non nul de
$H^0(\text{Cris}(X/S), \mathcal{O}_{X/S})$ d'après la
proposition \ref{proposition-compute-cohomology-crystal}.
De même, $F^*(\tau) = \sigma(\tau)$, où $\sigma : D \to D$ est le
morphisme induit par tout relèvement de Frobenius sur $P$. Si l'on choisit $\sigma(x) = x^p$
et $\sigma(y) = y^p$, un calcul simple montre que $F^*(\tau) = 0$.
\end{example}

\noindent
L'exemple suivant montre que, même pour l'espace affine de dimension $r$, la cohomologie
cristalline ne donne pas le résultat attendu.

\begin{example}
\label{example-affine-n-space}
Soit $A = \mathbf{Z}_p$, avec l'idéal à puissances divisées $(p)$ muni de
son unique structure de puissances divisées $\gamma$. Soit
$C = \mathbf{F}_p[x_1, \ldots, x_r]$. Choisissons la présentation
$$
C = P/J = P/pP\quad\text{avec}\quad P = \mathbf{Z}_p[x_1, \ldots, x_r]
$$
Notons que $pP$ est muni de puissances divisées d'après
Algèbre à puissances divisées, lemme \ref{dpa-lemma-gamma-extends}.
Ainsi, en posant $D = P^\wedge$ avec l'idéal à puissances divisées $(p)$, on obtient une
situation comme à la section \ref{section-quasi-coherent-crystals}.
On conclut que $R\Gamma(\text{Cris}(X/S), \mathcal{O}_{X/S})$
est représenté par le complexe
$$
D \to \Omega^1_D \to \Omega^2_D \to \ldots \to \Omega^r_D
$$
voir la proposition \ref{proposition-compute-cohomology-crystal}.
En supposant $r > 0$, on obtient les conclusions suivantes :
\begin{enumerate}
\item La cohomologie cristalline du faisceau structural cristallin
de $X = \mathbf{A}^r_{\mathbf{F}_p}$ sur $S = \Spec(\mathbf{Z}_p)$
est nulle sauf en degrés $0, \ldots, r$.
\item On a $H^0(\text{Cris}(X/S), \mathcal{O}_{X/S}) = \mathbf{Z}_p$.
\item Le groupe de cohomologie $H^r(\text{Cris}(X/S), \mathcal{O}_{X/S})$
est infini et n'est pas un groupe abélien de torsion.
\item Le groupe de cohomologie $H^r(\text{Cris}(X/S), \mathcal{O}_{X/S})$
n'est pas séparé pour la topologie $p$-adique.
\end{enumerate}
Si les deux premiers énoncés sont raisonnables, les assertions (3) et (4) sont
déconcertants ! Ces énoncés se vérifient immédiatement en
explicitant le complexe affiché ci-dessus. Faisons-le simplement
dans le cas $r = 1$. Nous considérons alors le complexe à deux termes
de modules $p$-adiquement complets
$$
\text{d} :
D = \left(
\bigoplus\nolimits_{n \geq 0} \mathbf{Z}_p x^n
\right)^\wedge
\longrightarrow
\Omega^1_D = \left(
\bigoplus\nolimits_{n \geq 1} \mathbf{Z}_p x^{n - 1}\text{d}x
\right)^\wedge
$$
Le morphisme est donné par $\text{diag}(0, 1, 2, 3, 4, \ldots)$, à ceci près que
le premier facteur de la somme directe manque au membre de droite. Il est alors clair
que $\bigoplus_{n > 0} \mathbf{Z}_p/n\mathbf{Z}_p$ est un sous-groupe
du conoyau, lequel est donc infini. En fait, l'élément
$$
\omega = \sum\nolimits_{e > 0} p^e x^{p^{2e} - 1}\text{d}x
$$
n'est manifestement pas un élément de torsion du conoyau. Mais la situation est pire.
En effet, considérons l'élément
$$
\eta = \sum\nolimits_{e > 0} p^e x^{p^e - 1}\text{d}x
$$
Pour tout $t > 0$, l'élément $\eta$ est congru à
$\sum_{e > t} p^e x^{p^e - 1}\text{d}x$ modulo l'image de
$\text{d}$ ; cette somme est divisible par $p^t$. Ainsi, la classe de cohomologie
de $\eta$ dans $H^1(\text{Cris}(X/S), \mathcal{O}_{X/S})$
appartient à $\bigcap p^tH^1(\text{Cris}(X/S), \mathcal{O}_{X/S})$.
Mais $\eta$ n'appartient pas à l'image de $\text{d}$ car il devrait
être l'image de $a + \sum_{e > 0} x^{p^e}$ pour un certain $a \in \mathbf{Z}_p$,
mais cette somme n'appartient pas au module figurant à gauche.
Ainsi, la classe de cohomologie de $\eta$ est non nulle et l'on voit que
l'assertion (4) est vérifiée. (En fait, la classe de cohomologie de $\eta$ n'est pas de torsion.)
En revanche, les groupes de cohomologie
$H^i(\text{Cris}(X/S), \mathcal{O}_{X/S})$
sont des $\mathbf{Z}_p$-modules complets au sens dérivé : ce sont les groupes de cohomologie
d'un complexe de modules $p$-adiquement complets, voir
Compléments d'algèbre, section \ref{more-algebra-section-derived-completion}.
\end{example}








\section{Applications}
\label{section-applications}

\noindent
Dans cette section, nous rassemblons quelques applications des résultats des
sections précédentes.

\begin{proposition}
\label{proposition-compare-with-de-Rham}
Dans la situation \ref{situation-global}.
Soit $\mathcal{F}$ un cristal en modules quasi-cohérents sur
$\text{Cris}(X/S)$. La flèche de troncation de complexes
$$
(\mathcal{F} \to
\mathcal{F} \otimes_{\mathcal{O}_{X/S}} \Omega^1_{X/S} \to
\mathcal{F} \otimes_{\mathcal{O}_{X/S}} \Omega^2_{X/S} \to \ldots)
\longrightarrow \mathcal{F}[0],
$$
bien qu'il ne soit pas un quasi-isomorphisme, le devient après application du foncteur
$Ru_{X/S, *}$. En fait, pour tout $i > 0$, on a 
$$
Ru_{X/S, *}(\mathcal{F} \otimes_{\mathcal{O}_{X/S}} \Omega^i_{X/S}) = 0.
$$
\end{proposition}

\begin{proof}
Le lemme \ref{lemma-automatic-connection} fournit un complexe de de Rham
comme indiqué dans l'énoncé. Abrégeons
$\mathcal{H} = \mathcal{F} \otimes \Omega^i_{X/S}$.
Soit $X' \subset X$ un sous-schéma ouvert affine
dont l'image est contenue dans un sous-schéma ouvert affine $S' \subset S$.
Alors
$$
(Ru_{X/S, *}\mathcal{H})|_{X'_{Zar}} =
Ru_{X'/S', *}(\mathcal{H}|_{\text{Cris}(X'/S')}),
$$
voir le lemme \ref{lemma-localize}. Ainsi, le
lemme \ref{lemma-cohomology-is-zero}
montre que $Ru_{X/S, *}\mathcal{H}$ est un complexe de faisceaux sur
$X_{Zar}$ dont la cohomologie est nulle sur tout ouvert affine.
Comme la topologie de $X$ possède une base formée d'ouverts affines,
cela implique que $Ru_{X/S, *}\mathcal{H}$ est quasi-isomorphe à zéro.
\end{proof}

\begin{remark}
\label{remark-vanishing}
La démonstration de la proposition \ref{proposition-compare-with-de-Rham}
montre que la conclusion
$$
Ru_{X/S, *}(\mathcal{F} \otimes_{\mathcal{O}_{X/S}} \Omega^i_{X/S}) = 0
$$
pour $i > 0$ vaut pour tout $\mathcal{O}_{X/S}$-module
$\mathcal{F}$ qui vérifie les conditions (1) et (2) de la
proposition \ref{proposition-compute-cohomology}.
Cela s'applique aux faisceaux suivants, qui ne sont pas des cristaux :
$\Omega^i_{X/S}$ pour tout $i$, ainsi qu'à tout faisceau de la forme
$\underline{\mathcal{F}}$, où $\mathcal{F}$ est un $\mathcal{O}_X$-module
quasi-cohérent. En particulier, cela s'applique au
faisceau $\underline{\mathcal{O}_X} = \underline{\mathbf{G}_a}$.
Notons toutefois qu'un résultat tel que le lemme \ref{lemma-automatic-connection}
est nécessaire pour produire un complexe de de Rham, ce qui exige que $\mathcal{F}$ soit un cristal.
Par conséquent, à l'heure actuelle, la classe des faisceaux de modules pour lesquels l'énoncé
complet de la proposition \ref{proposition-compare-with-de-Rham} est vérifié
est exactement la catégorie des cristaux en modules quasi-cohérents.
\end{remark}

\noindent
Plaçons-nous dans la situation \ref{situation-global}.
Soit $\mathcal{F}$ un cristal en modules quasi-cohérents sur
$\text{Cris}(X/S)$. Soit $(U, T, \delta)$ un objet de
$\text{Cris}(X/S)$. La proposition \ref{proposition-compare-with-de-Rham}
permet de construire un morphisme canonique
\begin{equation}
\label{equation-restriction}
R\Gamma(\text{Cris}(X/S), \mathcal{F})
\longrightarrow
R\Gamma(T, \mathcal{F}_T \otimes_{\mathcal{O}_T} \Omega^*_{T/S, \delta})
\end{equation}
En effet, on a $R\Gamma(\text{Cris}(X/S), \mathcal{F}) =
R\Gamma(\text{Cris}(X/S), \mathcal{F} \otimes \Omega^*_{X/S})$,
on peut restreindre les classes de cohomologie globale à $T$, et $\Omega_{X/S}$
se restreint à $\Omega_{T/S, \delta}$ par le
lemme \ref{lemma-module-of-differentials}.

















\section{Quelques résultats supplémentaires}
\label{section-missing}

\noindent
Dans cette section, nous mentionnons certains résultats dont nous ne donnons pas encore la démonstration.
Nous les formulons sous forme d'une suite de remarques et ne les transformerons
en véritables lemmes et propositions que lorsque nous ajouterons des
démonstrations détaillées.

\begin{remark}[Images directes supérieures]
\label{remark-compute-direct-image}
Soit $p$ un nombre premier. Soit
$(S, \mathcal{I}, \gamma) \to (S', \mathcal{I}', \gamma')$
un morphisme de schémas à puissances divisées sur $\mathbf{Z}_{(p)}$. Soit
$$
\xymatrix{
X \ar[r]_f \ar[d] & X' \ar[d] \\
S_0 \ar[r] & S'_0
}
$$
un diagramme commutatif de morphismes de schémas, et supposons que $p$ soit
localement nilpotent sur $X$ et $X'$. Soit $\mathcal{F}$ un
$\mathcal{O}_{X/S}$-module sur $\text{Cris}(X/S)$. Alors
$Rf_{\text{cris}, *}\mathcal{F}$ se calcule comme suit.

\medskip\noindent
Étant donné un objet $(U', T', \delta')$ de $\text{Cris}(X'/S')$, posons
$U = X \times_{X'} U' = f^{-1}(U')$ (un sous-schéma ouvert de $X$). Notons
$(T_0, T, \delta)$ le schéma à puissances divisées sur $S$ tel que
$$
\xymatrix{
T \ar[r] \ar[d] & T' \ar[d] \\
S \ar[r] & S'
}
$$
soit cartésien dans la catégorie des schémas à puissances divisées ; voir le
lemme \ref{lemma-fibre-product}. Il existe un
morphisme induit $U \to T_0$, et l'on obtient un morphisme
$(U/T)_{\text{cris}} \to (X/S)_{\text{cris}}$ ; voir la
remarque \ref{remark-functoriality-cris}.
Soit $\mathcal{F}_U$ l'image inverse de $\mathcal{F}$.
Soit $\tau_{U/T} : (U/T)_{\text{cris}} \to T_{Zar}$ le morphisme structural.
Alors on a
\begin{equation}
\label{equation-identify-pushforward}
\left(Rf_{\text{cris}, *}\mathcal{F}\right)_{T'} =
R(T \to T')_*\left(R\tau_{U/T, *} \mathcal{F}_U \right)
\end{equation}
où le membre de gauche est la restriction (voir la
section \ref{section-sheaves}).

\medskip\noindent
Indications : montrer d'abord que $\text{Cris}(U/T)$ est le topos obtenu en localisant $\text{Cris}(X/S)$ (au sens
du lemme \ref{sites-lemma-localize-topos-site} de Sites)
par rapport au faisceau d'ensembles $f_{\text{cris}}^{-1}h_{(U', T', \delta')}$. Réduire ensuite
l'énoncé au cas où $\mathcal{F}$ est un module injectif
et où l'on considère l'image directe de modules, en utilisant le fait que l'image inverse d'un
$\mathcal{O}_{X/S}$-module injectif est un $\mathcal{O}_{U/T}$-module injectif sur
$\text{Cris}(U/T)$. Enfin, vérifier le résultat pour l'image directe ordinaire.
\end{remark}

\begin{remark}[Mayer-Vietoris]
\label{remark-mayer-vietoris}
Dans la situation de la remarque \ref{remark-compute-direct-image},
supposons donné un recouvrement ouvert $X = X' \cup X''$. Notons
$X''' = X' \cap X''$. Soient $f'$, $f''$ et $f'''$ les restrictions de $f$
à $X'$, $X''$ et $X'''$. De plus, soient $\mathcal{F}'$, $\mathcal{F}''$
et $\mathcal{F}'''$ les restrictions de $\mathcal{F}$ aux sites cristallins
de $X'$, $X''$ et $X'''$. Il existe alors un triangle distingué
$$
Rf_{\text{cris}, *}\mathcal{F}
\longrightarrow
Rf'_{\text{cris}, *}\mathcal{F}' \oplus Rf''_{\text{cris}, *}\mathcal{F}''
\longrightarrow
Rf'''_{\text{cris}, *}\mathcal{F}'''
\longrightarrow
Rf_{\text{cris}, *}\mathcal{F}[1]
$$
dans $D(\mathcal{O}_{X'/S'})$.

\medskip\noindent
Indications : c'est une conséquence formelle du fait que les sous-catégories
$\text{Cris}(X'/S)$, $\text{Cris}(X''/S)$, $\text{Cris}(X'''/S)$ correspondent
à des sous-objets ouverts du faisceau final sur $\text{Cris}(X/S)$ et que la
dernière est l'intersection des deux premières.
\end{remark}

\begin{remark}[Complexe de {\v C}ech]
\label{remark-cech-complex}
Soit $p$ un nombre premier. Soit $(A, I, \gamma)$ un anneau à puissances
divisées, $A$ étant une $\mathbf{Z}_{(p)}$-algèbre. Posons $S = \Spec(A)$ et
$S_0 = \Spec(A/I)$. Soit $X$ un schéma séparé\footnote{Cette hypothèse n'est
pas strictement nécessaire : à l'aide d'hyperrecouvrements, la construction de la
remarque s'étend au cas général.} sur
$S_0$, tel que $p$ soit localement nilpotent sur $X$. Soit $\mathcal{F}$ un
cristal en $\mathcal{O}_{X/S}$-modules quasi-cohérents.

\medskip\noindent
Choisissons un recouvrement ouvert affine
$X = \bigcup_{\lambda \in \Lambda} U_\lambda$ de $X$.
Écrivons $U_\lambda = \Spec(C_\lambda)$. Choisissons une algèbre de polynômes
$P_\lambda$ sur $A$ et une surjection $P_\lambda \to C_\lambda$.
Ces choix étant fixés, on peut construire un complexe de {\v C}ech qui
calcule $R\Gamma(\text{Cris}(X/S), \mathcal{F})$.

\medskip\noindent
Pour $n \geq 0$ et $\lambda_0, \ldots, \lambda_n \in \Lambda$,
posons $U_{\lambda_0 \ldots \lambda_n} = U_{\lambda_0} \cap \ldots
\cap U_{\lambda_n}$. Ce schéma est affine par hypothèse. Écrivons
$U_{\lambda_0 \ldots \lambda_n} = \Spec(C_{\lambda_0 \ldots \lambda_n})$.
Posons
$$
P_{\lambda_0 \ldots \lambda_n} =
P_{\lambda_0} \otimes_A \ldots \otimes_A P_{\lambda_n}
$$
qui est muni d'une surjection canonique vers $C_{\lambda_0 \ldots \lambda_n}$.
Notons son noyau $J_{\lambda_0 \ldots \lambda_n}$ et notons
$D_{\lambda_0 \ldots \lambda_n}$
l'enveloppe à puissances divisées, complétée $p$-adiquement, de
$J_{\lambda_0 \ldots \lambda_n}$ dans $P_{\lambda_0 \ldots \lambda_n}$
relativement à $\gamma$. Soit $M_{\lambda_0 \ldots \lambda_n}$ le
$P_{\lambda_0 \ldots \lambda_n}$-module correspondant
à la restriction de $\mathcal{F}$ à
$\text{Cris}(U_{\lambda_0 \ldots \lambda_n}/S)$ d'après
la proposition \ref{proposition-crystals-on-affine}.
Par construction, on obtient un anneau cosimplicial à puissances divisées $D(*)$
dont le terme de degré $n$ est l'anneau
$$
D(n) =
\prod\nolimits_{\lambda_0 \ldots \lambda_n}
D_{\lambda_0 \ldots \lambda_n}
$$
(on utilise la fonctorialité des enveloppes à puissances divisées et la structure
cosimpliciale triviale sur l'anneau $P(*)$ défini de même).
Comme $M_{\lambda_0 \ldots \lambda_n}$ est la « valeur » de $\mathcal{F}$
sur les objets $\Spec(D_{\lambda_0 \ldots \lambda_n})$, on voit que
$M(*)$ défini par la règle
$$
M(n) = \prod\nolimits_{\lambda_0 \ldots \lambda_n}
M_{\lambda_0 \ldots \lambda_n}
$$
définit un module cosimplicial sur $D(*)$. Nous affirmons maintenant que
$$
R\Gamma(\text{Cris}(X/S), \mathcal{F}) = s(M(*))
$$
Ici, $s(-)$ désigne le complexe de cochaînes associé à un module
cosimplicial (voir
Objets simpliciaux, section \ref{simplicial-section-dold-kan-cosimplicial}).

\medskip\noindent
Indications : la démonstration est semblable à celle de la
proposition \ref{proposition-compute-cohomology} (en particulier,
le résultat vaut pour tout module qui vérifie les hypothèses de
cette proposition).
\end{remark}

\begin{remark}[Complexe de {\v C}ech alterné]
\label{remark-alternating-cech-complex}
Soit $p$ un nombre premier. Soit $(A, I, \gamma)$ un anneau à puissances
divisées, $A$ étant une $\mathbf{Z}_{(p)}$-algèbre. Posons $S = \Spec(A)$ et
$S_0 = \Spec(A/I)$. Soit $X$ un schéma séparé quasi-compact
sur $S_0$, tel que $p$ soit localement nilpotent sur $X$. Soit
$\mathcal{F}$ un cristal en $\mathcal{O}_{X/S}$-modules quasi-cohérents.

\medskip\noindent
Choisissons un recouvrement ouvert affine fini
$X = \bigcup_{\lambda \in \Lambda} U_\lambda$ de $X$
et un ordre total sur $\Lambda$.
Écrivons $U_\lambda = \Spec(C_\lambda)$. Choisissons une algèbre de polynômes
$P_\lambda$ sur $A$ et une surjection $P_\lambda \to C_\lambda$.
Ces choix étant fixés, on peut construire un complexe
de {\v C}ech alterné qui calcule $R\Gamma(\text{Cris}(X/S), \mathcal{F})$.

\medskip\noindent
Nous allons utiliser les notations introduites dans la
remarque \ref{remark-cech-complex}.
Notons $\Omega_{\lambda_0 \ldots \lambda_n}$
le module des différentielles, complété $p$-adiquement, de
$D_{\lambda_0 \ldots \lambda_n}$ sur $A$ compatible avec la structure à puissances
divisées. Soit $\nabla$ la connexion intégrable sur
$M_{\lambda_0 \ldots \lambda_n}$ provenant de la
proposition \ref{proposition-crystals-on-affine}.
Considérons le complexe double $M^{\bullet, \bullet}$ dont les
termes sont
$$
M^{n, m} =
\bigoplus\nolimits_{\lambda_0 < \ldots < \lambda_n}
M_{\lambda_0 \ldots \lambda_n}
\otimes^\wedge_{D_{\lambda_0 \ldots \lambda_n}}
\Omega^m_{D_{\lambda_0 \ldots \lambda_n}}.
$$
Pour la différentielle $d_1$ (qui augmente $n$), utilisons la différentielle
de {\v C}ech usuelle et, pour la différentielle $d_2$,
la connexion, c'est-à-dire la différentielle du complexe de de Rham.
Nous affirmons que
$$
R\Gamma(\text{Cris}(X/S), \mathcal{F}) = \text{Tot}(M^{\bullet, \bullet})
$$
Ici, $\text{Tot}(-)$ désigne le complexe total associé à un
complexe double ; voir
Homologie, définition \ref{homology-definition-associated-simple-complex}.

\medskip\noindent
Indications : on a
$$
R\Gamma(\text{Cris}(X/S), \mathcal{F}) = R\Gamma(\text{Cris}(X/S),
\mathcal{F} \otimes_{\mathcal{O}_{X/S}} \Omega_{X/S}^\bullet)
$$
par la proposition \ref{proposition-compare-with-de-Rham}.
Le membre de droite de la formule n'est autre que le complexe de {\v C}ech alterné
du recouvrement $X = \bigcup_{\lambda \in \Lambda} U_\lambda$
(qui induit un recouvrement ouvert du faisceau final de $\text{Cris}(X/S)$)
et du complexe $\mathcal{F} \otimes_{\mathcal{O}_{X/S}} \Omega_{X/S}^\bullet$ ;
voir la proposition \ref{proposition-compute-cohomology-crystal}.
Le résultat découle alors d'un résultat général sur la cohomologie des sites,
selon lequel le complexe de {\v C}ech alterné calcule la cohomologie
dès lors qu'il la calcule correctement sur tous les morceaux (insérer ici une
référence ultérieure).
\end{remark}

\begin{remark}[Quasi-cohérence]
\label{remark-quasi-coherent}
Dans la situation de la remarque \ref{remark-compute-direct-image},
supposons que $S \to S'$ soit quasi-compact et quasi-séparé et
que $X \to S_0$ soit quasi-compact et quasi-séparé. Alors, pour un cristal
en $\mathcal{O}_{X/S}$-modules quasi-cohérents $\mathcal{F}$,
les faisceaux $R^if_{\text{cris}, *}\mathcal{F}$ sont localement quasi-cohérents.

\medskip\noindent
Indications : il faut montrer que les restrictions à $T'$ sont des
$\mathcal{O}_{T'}$-modules quasi-cohérents, où $(U', T', \delta')$ est un objet quelconque de
$\text{Cris}(X'/S')$. Il suffit de le faire lorsque $T'$ est affine.
On utilise la formule (\ref{equation-identify-pushforward}),
le fait que $T \to T'$ est quasi-compact et quasi-séparé (car $T$
est affine au-dessus du changement de base de $T'$ le long de $S \to S'$), ainsi que le
lemme de Cohomologie des schémas
\ref{coherent-lemma-quasi-coherence-higher-direct-images}
pour voir qu'il suffit de montrer que les faisceaux
$R^i\tau_{U/T, *}\mathcal{F}_U$ sont quasi-cohérents.
Notons que $U \to T_0$ est aussi quasi-compact et quasi-séparé ; voir
Schémas, lemmes \ref{schemes-lemma-quasi-compact-permanence} et
\ref{schemes-lemma-compose-after-separated}.

\medskip\noindent
On est ainsi ramené à prouver que $R^i\tau_{X/S, *}\mathcal{F}$
est quasi-cohérent sur $S$ dans le cas où $p$ est localement nilpotent sur $S$. Ici,
$\tau_{X/S}$ est le morphisme structural ; voir la
remarque \ref{remark-structure-morphism}.
On peut travailler localement sur $S$, et donc supposer $S$ affine
(voir le lemme \ref{lemma-localize}). Une récurrence sur le nombre
d'ouverts affines recouvrant $X$ et Mayer--Vietoris
(remarque \ref{remark-mayer-vietoris}) ramènent la question
au cas où $X$ est lui aussi affine (comme dans la démonstration du
lemme de Cohomologie des schémas
\ref{coherent-lemma-quasi-coherence-higher-direct-images}).
Écrivons $X = \Spec(C)$ et $S = \Spec(A)$, de sorte que $(A, I, \gamma)$ et
$A \to C$ soient comme
dans la situation \ref{situation-affine}. Choisissons une algèbre de polynômes
$P$ sur $A$ et une surjection $P \to C$ comme dans la
section \ref{section-quasi-coherent-crystals}.
Soit $(M, \nabla)$ le module correspondant à $\mathcal{F}$ ; voir la
proposition \ref{proposition-crystals-on-affine}.
En appliquant la 
proposition \ref{proposition-compute-cohomology-crystal},
on voit que $R\Gamma(\text{Cris}(X/S), \mathcal{F})$ est représenté par
$M \otimes_D \Omega_D^*$. Notons que la complétion n'est pas nécessaire,
car $p$ est nilpotent dans $A$ ! Il faut montrer que cela est compatible
avec le passage aux ouverts principaux de $S = \Spec(A)$. Supposons que $g \in A$.
On conclut de même que $R\Gamma(\text{Cris}(X_g/S_g), \mathcal{F})$
est calculé par $M_g \otimes_{D_g} \Omega_{D_g}^*$ (on utilise encore le fait que
la complétion $p$-adique n'est pas nécessaire). La conclusion s'ensuit puisque la localisation
est un foncteur exact dans la catégorie des $A$-modules.
\end{remark}

\begin{remark}[Bornitude]
\label{remark-bounded-cohomology}
Dans la situation de la remarque \ref{remark-compute-direct-image},
supposons que $S \to S'$ soit quasi-compact et quasi-séparé et
que $X \to S_0$ soit de type fini et quasi-séparé. Alors il existe
un entier $i_0$ tel que, pour tout cristal
en $\mathcal{O}_{X/S}$-modules quasi-cohérents $\mathcal{F}$,
on ait $R^if_{\text{cris}, *}\mathcal{F} = 0$ pour tout $i > i_0$.

\medskip\noindent
Indications : en raisonnant comme dans la remarque \ref{remark-quasi-coherent} (à l'aide du
lemme de Cohomologie des schémas
\ref{coherent-lemma-quasi-coherence-higher-direct-images})
on se ramène à prouver que $H^i(\text{Cris}(X/S), \mathcal{F}) = 0$ pour $i \gg 0$
dans la situation de la proposition \ref{proposition-compute-cohomology-crystal},
lorsque $C$ est une $A$-algèbre de type fini. C'est clair, car on peut
choisir une algèbre de polynômes en un nombre fini de variables et l'on a $\Omega^i_D = 0$
pour $i \gg 0$.
\end{remark}

\begin{remark}[Bornes particulières]
\label{remark-bounded-cohomology-over-point}
Dans la situation \ref{situation-global}, soit $\mathcal{F}$ un cristal en
$\mathcal{O}_{X/S}$-modules quasi-cohérents. Supposons que $S_0$
ait un unique point et que $X \to S_0$ soit de présentation finie.
\begin{enumerate}
\item Si $\dim X = d$ et si $X/S_0$ est de dimension de plongement $e$, alors
$H^i(\text{Cris}(X/S), \mathcal{F}) = 0$ pour $i > d + e$.
\item Si $X$ est séparé et peut être recouvert par $q$ ouverts affines, et si
$X/S_0$ est de dimension de plongement $e$, alors
$H^i(\text{Cris}(X/S), \mathcal{F}) = 0$ pour $i > q + e$.
\end{enumerate}
Indications : dans le cas (1), on peut utiliser l'égalité
$$
H^i(\text{Cris}(X/S), \mathcal{F}) = H^i(X_{Zar}, Ru_{X/S, *}\mathcal{F})
$$
et le fait que $Ru_{X/S, *}\mathcal{F}$ se calcule localement par un complexe de de Rham
construit à l'aide d'une immersion de $X$ dans un schéma lisse
de dimension $e$ sur $S$
(voir le lemme \ref{lemma-compute-cohomology-crystal-smooth}).
Ces complexes de de Rham sont nuls en tout degré $> e$. Ainsi, (1)
résulte de Cohomologie, proposition
\ref{cohomology-proposition-vanishing-Noetherian}.
Dans le cas (2), on utilise le complexe de {\v C}ech alterné (voir la
remarque \ref{remark-alternating-cech-complex}) pour se ramener au cas où
$X$ est affine. Dans le cas affine, on démontre le résultat à l'aide du complexe de de Rham
associé à une immersion de $X$ dans un schéma lisse de dimension $e$
sur $S$ (la construction d'un tel objet demande un peu de travail).
\end{remark}

\begin{remark}[Morphisme de changement de base]
\label{remark-base-change}
Dans la situation de la remarque \ref{remark-compute-direct-image},
supposons que $S = \Spec(A)$ et $S' = \Spec(A')$ soient affines.
Soit $\mathcal{F}'$ un $\mathcal{O}_{X'/S'}$-module.
Soit $\mathcal{F}$ l'image inverse de $\mathcal{F}'$.
Il existe alors un morphisme canonique de changement de base
$$
L(S' \to S)^*R\tau_{X'/S', *}\mathcal{F}'
\longrightarrow
R\tau_{X/S, *}\mathcal{F}
$$
où $\tau_{X/S}$ et $\tau_{X'/S'}$ sont les morphismes structuraux ; voir la
remarque \ref{remark-structure-morphism}. Après passage aux sections globales,
on obtient un morphisme de changement de base
\begin{equation}
\label{equation-base-change-map}
R\Gamma(\text{Cris}(X'/S'), \mathcal{F}') \otimes^\mathbf{L}_{A'} A
\longrightarrow
R\Gamma(\text{Cris}(X/S), \mathcal{F})
\end{equation}
dans $D(A)$.

\medskip\noindent
Indication : on compose le morphisme très général de changement de base de
Cohomologie sur les sites, remarque \ref{sites-cohomology-remark-base-change},
avec le morphisme canonique
$Lf_{\text{cris}}^*\mathcal{F}' \to
f_{\text{cris}}^*\mathcal{F}' = \mathcal{F}$.
\end{remark}

\begin{remark}[Isomorphisme de changement de base]
\label{remark-base-change-isomorphism}
Le morphisme (\ref{equation-base-change-map}) est un isomorphisme pourvu que
toutes les conditions suivantes soient satisfaites :
\begin{enumerate}
\item $p$ est nilpotent dans $A'$,
\item $\mathcal{F}'$ est un cristal en
$\mathcal{O}_{X'/S'}$-modules quasi-cohérents,
\item $X' \to S'_0$ est un morphisme quasi-compact et quasi-séparé,
\item $X = X' \times_{S'_0} S_0$,
\item $\mathcal{F}'$ est un $\mathcal{O}_{X'/S'}$-module plat,
\item $X' \to S'_0$ est un morphisme d'intersection complète locale (voir
Compléments sur les morphismes, définition \ref{more-morphisms-definition-lci} ; c'est
par exemple le cas si $X' \to S'_0$ est syntomique ou lisse),
\item $X'$ et $S_0$ sont Tor-indépendants sur $S'_0$ (voir
Compléments d'algèbre, définition \ref{more-algebra-definition-tor-independent} ;
c'est par exemple le cas si $S_0 \to S'_0$ ou $X' \to S'_0$ est plat).
\end{enumerate}
Indications : la condition (1) signifie que, dans les arguments ci-dessous, la complétion $p$-adique
n'a aucun effet et peut être ignorée.
À l'aide de la condition (3) et de Mayer--Vietoris (voir la
remarque \ref{remark-mayer-vietoris}), on se ramène au cas
où $X'$ est affine. En fait, par la condition (6), après avoir encore restreint,
on peut supposer que $X' = \Spec(C')$ et que l'on dispose d'une présentation
$C' = A'/I'[x_1, \ldots, x_n]/(\bar f'_1, \ldots, \bar f'_c)$
où $\bar f'_1, \ldots, \bar f'_c$ est une suite Koszul-régulière dans $A'/I'$.
(Cela signifie que, localement pour la topologie lisse, $\bar f'_1, \ldots, \bar f'_c$ forme
une suite régulière ; voir Compléments d'algèbre,
lemme \ref{more-algebra-lemma-Koszul-regular-flat-locally-regular}.)
Choisissons un relèvement de
$\bar f'_i$ en un élément $f'_i \in A'[x_1, \ldots, x_n]$. D'après (4), on voit que
$X = \Spec(C)$, où $C = A/I[x_1, \ldots, x_n]/(\bar f_1, \ldots, \bar f_c)$,
où $f_i \in A[x_1, \ldots, x_n]$ est l'image de $f'_i$.
La propriété (7) montre que $\bar f_1, \ldots, \bar f_c$ est une suite Koszul-régulière
dans $A/I[x_1, \ldots, x_n]$. L'enveloppe à puissances divisées de
$I'A'[x_1, \ldots, x_n] + (f'_1, \ldots, f'_c)$ dans $A'[x_1, \ldots, x_n]$
relativement à $\gamma'$ est
$$
D' = A'[x_1, \ldots, x_n]\langle \xi_1, \ldots, \xi_c \rangle/(\xi_i - f'_i)
$$
voir le lemme \ref{lemma-describe-divided-power-envelope}. On vérifie alors que
$\xi_1 - f'_1, \ldots, \xi_n - f'_n$ est une suite Koszul-régulière dans
l'anneau $A'[x_1, \ldots, x_n]\langle \xi_1, \ldots, \xi_c\rangle$.
De même, l'enveloppe à puissances divisées de
$IA[x_1, \ldots, x_n] + (f_1, \ldots, f_c)$ dans $A[x_1, \ldots, x_n]$
relativement à $\gamma$ est
$$
D = A[x_1, \ldots, x_n]\langle \xi_1, \ldots, \xi_c\rangle/(\xi_i - f_i)
$$
et $\xi_1 - f_1, \ldots, \xi_n - f_n$ est une suite Koszul-régulière dans
l'anneau $A[x_1, \ldots, x_n]\langle \xi_1, \ldots, \xi_c\rangle$.
Il s'ensuit que $D' \otimes_{A'}^\mathbf{L} A = D$. La condition (2)
implique que $\mathcal{F}'$ correspond à un couple $(M', \nabla)$
formé d'un $D'$-module muni d'une connexion ; voir la
proposition \ref{proposition-crystals-on-affine}.
Alors $M = M' \otimes_{D'} D$ correspond à l'image inverse $\mathcal{F}$.
L'hypothèse (5) montre que $M'$ est un $D'$-module plat, donc
$$
M = M' \otimes_{D'} D = M' \otimes_{D'} D' \otimes_{A'}^\mathbf{L} A
= M' \otimes_{A'}^\mathbf{L} A
$$
Comme les modules des différentielles $\Omega_{D'}$ et $\Omega_D$
(tels que définis dans la section \ref{section-quasi-coherent-crystals})
sont des $D'$-modules libres sur les mêmes générateurs, on voit que
$$
M \otimes_D \Omega^\bullet_D =
M' \otimes_{D'} \Omega^\bullet_{D'} \otimes_{D'} D =
M' \otimes_{D'} \Omega^\bullet_{D'} \otimes_{A'}^\mathbf{L} A
$$
ce qui démontre l'assertion par la
proposition \ref{proposition-compute-cohomology-crystal}.
\end{remark}

\begin{remark}[Rlim]
\label{remark-rlim}
Soit $p$ un nombre premier. Soit $(A, I, \gamma)$ un anneau à puissances
divisées, $A$ étant une algèbre sur $\mathbf{Z}_{(p)}$, et supposons $p$ nilpotent
dans $A/I$. Posons $S = \Spec(A)$ et $S_0 = \Spec(A/I)$.
Soit $X$ un schéma sur $S_0$, tel que $p$ soit localement
nilpotent sur $X$. Soit $\mathcal{F}$ un
$\mathcal{O}_{X/S}$-module quelconque. Pour $e \gg 0$, l'idéal $(p^e) \subset I$
est stable par $\gamma$ ; voir
Algèbre à puissances divisées, lemme \ref{dpa-lemma-extend-to-completion}.
Posons $S_e = \Spec(A/p^eA)$ pour $e \gg 0$.
Alors $\text{Cris}(X/S_e)$ est une sous-catégorie pleine de $\text{Cris}(X/S)$,
et notons $\mathcal{F}_e$ la restriction de $\mathcal{F}$ à
$\text{Cris}(X/S_e)$. Alors
$$
R\Gamma(\text{Cris}(X/S), \mathcal{F}) =
R\lim_e R\Gamma(\text{Cris}(X/S_e), \mathcal{F}_e)
$$

\medskip\noindent
Indications : il suffit de le prouver lorsque $\mathcal{F}$ est injectif.
Dans ce cas, les faisceaux $\mathcal{F}_e$ sont eux aussi des
modules injectifs, les applications de transition
$\Gamma(\mathcal{F}_{e + 1}) \to \Gamma(\mathcal{F}_e)$ sont
surjectives, et l'on a
$\Gamma(\mathcal{F}) = \lim_e \Gamma(\mathcal{F}_e)$, car
tout objet de $\text{Cris}(X/S)$ est localement un objet de l'une
des catégories $\text{Cris}(X/S_e)$, par définition de
$\text{Cris}(X/S)$.
\end{remark}

\begin{remark}[Comparaison]
\label{remark-comparison}
Soit $p$ un nombre premier. Soit $(A, I, \gamma)$ un anneau à puissances
divisées, avec $p$ nilpotent dans $A$. Posons $S = \Spec(A)$ et
$S_0 = \Spec(A/I)$. Soit $Y$ un schéma lisse sur $S$, et posons
$X = Y \times_S S_0$. Soit
$\mathcal{F}$ un cristal en $\mathcal{O}_{X/S}$-modules quasi-cohérents.
Alors
\begin{enumerate}
\item $\gamma$ s'étend en une structure à puissances divisées sur l'idéal
de $X$ dans $Y$, de sorte que $(X, Y, \gamma)$ soit un objet de $\text{Cris}(X/S)$,
\item la restriction $\mathcal{F}_Y$ (voir la section \ref{section-sheaves})
est munie d'une connexion intégrable canonique
$\nabla : \mathcal{F}_Y \to
\mathcal{F}_Y \otimes_{\mathcal{O}_Y} \Omega_{Y/S}$, et
\item on a
$$
R\Gamma(\text{Cris}(X/S), \mathcal{F}) =
R\Gamma(Y, \mathcal{F}_Y \otimes_{\mathcal{O}_Y} \Omega^\bullet_{Y/S})
$$
dans $D(A)$.
\end{enumerate}
Indications : voir Algèbre à puissances divisées, lemme \ref{dpa-lemma-gamma-extends}, pour (1).
Voir le lemme \ref{lemma-automatic-connection} pour (2).
Pour la partie (3), notons qu'il existe un morphisme ; voir
(\ref{equation-restriction}). Ce morphisme est un isomorphisme lorsque
$X$ est affine ; voir le
lemme \ref{lemma-compute-cohomology-crystal-smooth}.
Cela montre que $Ru_{X/S, *}\mathcal{F}$ et
$\mathcal{F}_Y \otimes \Omega^\bullet_{Y/S}$ sont quasi-isomorphes
comme complexes sur $Y_{Zar} = X_{Zar}$.
Puisque $R\Gamma(\text{Cris}(X/S), \mathcal{F}) =
R\Gamma(X_{Zar}, Ru_{X/S, *}\mathcal{F})$, le résultat en découle.
\end{remark}

\begin{remark}[Caractère parfait]
\label{remark-perfect}
Soit $p$ un nombre premier. Soit $(A, I, \gamma)$ un anneau à puissances
divisées, avec $p$ nilpotent dans $A$. Posons $S = \Spec(A)$ et
$S_0 = \Spec(A/I)$. Soit $X$ un schéma propre et lisse sur $S_0$.
Soit $\mathcal{F}$ un cristal en
$\mathcal{O}_{X/S}$-modules quasi-cohérents localement libres de type fini.
Alors $R\Gamma(\text{Cris}(X/S), \mathcal{F})$ est un
objet parfait de $D(A)$.

\medskip\noindent
Indications : la remarque \ref{remark-base-change-isomorphism} donne
$$
R\Gamma(\text{Cris}(X/S), \mathcal{F}) \otimes_A^\mathbf{L} A/I
\cong
R\Gamma(\text{Cris}(X/S_0), \mathcal{F}|_{\text{Cris}(X/S_0)})
$$
La remarque \ref{remark-comparison} donne
$$
R\Gamma(\text{Cris}(X/S_0), \mathcal{F}|_{\text{Cris}(X/S_0)}) =
R\Gamma(X, \mathcal{F}_X \otimes \Omega^\bullet_{X/S_0})
$$
En utilisant la filtration bête du complexe de de Rham, on voit que
le dernier complexe affiché est parfait dans $D(A/I)$ dès que les complexes
$$
R\Gamma(X, \mathcal{F}_X \otimes \Omega^q_{X/S_0})
$$
sont des complexes parfaits dans $D(A/I)$ ; voir
Compléments d'algèbre, lemme \ref{more-algebra-lemma-two-out-of-three-perfect}.
C'est vrai par les arguments usuels
de cohomologie cohérente, car $\mathcal{F}_X \otimes \Omega^q_{X/S_0}$
est un faisceau localement libre de type fini et $X \to S_0$ est propre et plat
(insérer ici une référence ultérieure). En appliquant le
lemme \ref{more-algebra-lemma-perfect-modulo-nilpotent-ideal} de Compléments d'algèbre,
on voit que
$$
R\Gamma(\text{Cris}(X/S), \mathcal{F}) \otimes_A^\mathbf{L} A/I^n
$$
est un objet parfait de $D(A/I^n)$ pour tout $n$. Cela ne suffit pas tout à fait,
sauf si $A$ est noethérien. En effet, bien que $I$ soit localement nilpotent
par l'hypothèse de nilpotence de $p$ ; voir
Algèbre à puissances divisées, lemme \ref{dpa-lemma-nil},
on ne peut conclure que $I^n = 0$ pour un certain $n$. Un contre-exemple
est $\mathbf{F}_p\langle x \rangle$. Pour le démontrer en général lorsque
$\mathcal{F} = \mathcal{O}_{X/S}$, l'argument de
\url{https://math.columbia.edu/~dejong/wordpress/?p=2227}
convient. Lorsque les coefficients $\mathcal{F}$ sont non triviaux,
l'argument de \cite{Faltings-very} semble être le suivant. On se ramène au
cas $pA = 0$ par le lemme de Compléments d'algèbre
\ref{more-algebra-lemma-perfect-modulo-nilpotent-ideal}.
Dans ce cas, le morphisme de Frobenius $A \to A$, $a \mapsto a^p$, se factorise
en $A \to A/I \xrightarrow{\varphi} A$ (car $x^p = 0$ pour $x \in I$). Posons
$X^{(1)} = X \otimes_{A/I, \varphi} A$. Le morphisme de Frobenius absolu
de $X$ se factorise par un morphisme $F_X : X \to X^{(1)}$ (une sorte de
Frobenius relatif). Localement sur les ouverts affines, si $X = \Spec(C)$, alors
$X^{(1)} = \Spec( C \otimes_{A/I, \varphi} A)$
et $F_X$ correspond à $C \otimes_{A/I, \varphi} A \to C$,
$c \otimes a \mapsto c^pa$. Cela définit des morphismes de topos annelés
$$
(X/S)_{\text{cris}}
\xrightarrow{(F_X)_{\text{cris}}}
(X^{(1)}/S)_{\text{cris}}
\xrightarrow{u_{X^{(1)}/S}}
\Sh(X^{(1)}_{Zar})
$$
dont le composé est noté $\text{Frob}_X$. On montre alors que
$R\text{Frob}_{X, *}\mathcal{F}$ est représenté par un
complexe parfait de $\mathcal{O}_{X^{(1)}}$-modules (!)
par un calcul local.
\end{remark}

\begin{remark}[Caractère parfait dans le cas complet]
\label{remark-complete-perfect}
Soit $p$ un nombre premier. Soit $(A, I, \gamma)$ un anneau à puissances
divisées, $A$ étant complet $p$-adiquement et $p$ nilpotent dans $A/I$. Posons
$S = \Spec(A)$ et $S_0 = \Spec(A/I)$. Soit $X$ un schéma propre
et lisse sur $S_0$. Soit $\mathcal{F}$ un cristal en
$\mathcal{O}_{X/S}$-modules quasi-cohérents localement libres de type fini.
Alors $R\Gamma(\text{Cris}(X/S), \mathcal{F})$ est un
objet parfait de $D(A)$.

\medskip\noindent
Indications : on sait que $K = R\Gamma(\text{Cris}(X/S), \mathcal{F})$
est la limite dérivée $K = R\lim K_e$ des cohomologies sur $A/p^eA$ ;
voir la remarque \ref{remark-rlim}.
Chaque $K_e$ est un complexe parfait dans $D(A/p^eA)$ par la
remarque \ref{remark-perfect}.
Comme $A$ est complet $p$-adiquement, le résultat
résulte du
lemme \ref{more-algebra-lemma-Rlim-perfect-gives-complete} de Compléments d'algèbre.
\end{remark}

\begin{remark}[Comparaison dans le cas complet]
\label{remark-complete-comparison}
Soit $p$ un nombre premier. Soit $(A, I, \gamma)$ un anneau à puissances
divisées, $A$ étant un anneau noethérien complet $p$-adiquement et $p$ nilpotent
dans $A/I$. Posons $S = \Spec(A)$ et
$S_0 = \Spec(A/I)$. Soit $Y$ un schéma propre et lisse sur $S$, et posons
$X = Y \times_S S_0$. Soit $\mathcal{F}$ un cristal de type fini en
$\mathcal{O}_{X/S}$-modules quasi-cohérents. Alors
\begin{enumerate}
\item il existe un $\mathcal{O}_Y$-module cohérent $\mathcal{F}_Y$
muni d'une connexion intégrable
$$
\nabla :
\mathcal{F}_Y
\longrightarrow
\mathcal{F}_Y \otimes_{\mathcal{O}_Y} \Omega_{Y/S}
$$
tel que $\mathcal{F}_Y/p^e\mathcal{F}_Y$ soit le module muni de sa connexion
sur $A/p^eA$ obtenu dans la remarque \ref{remark-comparison}, et
\item on a
$$
R\Gamma(\text{Cris}(X/S), \mathcal{F}) =
R\Gamma(Y, \mathcal{F}_Y \otimes_{\mathcal{O}_Y} \Omega^\bullet_{Y/S})
$$
dans $D(A)$.
\end{enumerate}
Indications : l'existence de $\mathcal{F}_Y$ résulte du théorème d'existence de Grothendieck
(insérer ici une référence ultérieure). L'isomorphisme des cohomologies résulte
du fait que les deux membres sont calculés comme $R\lim$ des versions modulo $p^e$
(voir la remarque \ref{remark-rlim} pour le membre de gauche ; utiliser le théorème
des fonctions formelles, voir le
théorème \ref{coherent-theorem-formal-functions} de Cohomologie des schémas
pour le membre de droite).
Chacune des versions modulo $p^e$ est isomorphe par la
remarque \ref{remark-comparison}.
\end{remark}




\section{Image inverse le long de morphismes purement inséparables}
\label{section-pull-back-along-pth-root}

\noindent
Par revêtement $\alpha_p$, on entend un morphisme de la forme
$$
X' = \Spec(C[z]/(z^p - c)) \longrightarrow \Spec(C) = X
$$
où $C$ est une $\mathbf{F}_p$-algèbre et $c \in C$. De manière équivalente,
$X'$ est un $\alpha_p$-torseur sur $X$. Un {\it revêtement
$\alpha_p$ itéré}\footnote{Cette notation n'est pas standard.}
est un morphisme de schémas en caractéristique
$p$ qui, localement sur la cible, est composé d'un nombre fini de
revêtements $\alpha_p$. Dans cette section, nous montrons que l'image inverse par
un tel morphisme induit un quasi-isomorphisme en cohomologie cristalline
après inversion du nombre premier $p$. En fait, nous démontrons une version précise
de ce résultat. Commençons par un lemme préliminaire dont la formulation
nécessite quelques notations.

\medskip\noindent
Supposons donné un homomorphisme d'anneaux $B \to B'$ et des quotients $\Omega_B \to \Omega$ et
$\Omega_{B'} \to \Omega'$ vérifiant les hypothèses de la
remarque \ref{remark-base-change-connection}.
Ainsi, (\ref{equation-base-change-map-complexes}) fournit un
morphisme canonique de complexes
$$
c_M^\bullet :
M \otimes_B \Omega^\bullet
\longrightarrow
M \otimes_B (\Omega')^\bullet
$$
pour tout $B$-module $M$ muni d'une connexion intégrable
$\nabla : M \to M \otimes_B \Omega_B$.

\medskip\noindent
Supposons donnés $a \in B$, $z \in B'$ et une application $\theta : B' \to B'$
vérifiant les hypothèses suivantes :
\begin{enumerate}
\item
\label{item-d-a-zero}
$\text{d}(a) = 0$,
\item
\label{item-direct-sum}
$\Omega' = B' \otimes_B \Omega \oplus B'\text{d}z$ ; nous écrivons
$\text{d}(f) = \text{d}_1(f) + \partial_z(f) \text{d}z$
où $\text{d}_1(f) \in B' \otimes \Omega$ et $\partial_z(f) \in B'$
pour tout $f \in B'$,
\item
\label{item-theta-linear}
$\theta : B' \to B'$ est $B$-linéaire,
\item
\label{item-integrate}
$\partial_z \circ \theta = a$,
\item
\label{item-injective}
$B \to B'$ est universellement injectif (et donc $\Omega \to \Omega'$
est injectif),
\item
\label{item-factor}
$af - \theta(\partial_z(f)) \in B$ pour tout $f \in B'$,
\item
\label{item-horizontal}
$(\theta \otimes 1)(\text{d}_1(f)) - \text{d}_1(\theta(f)) \in \Omega$
pour tout $f \in B'$, où
$\theta \otimes 1 : B' \otimes \Omega \to B' \otimes \Omega$
\end{enumerate}
Ces conditions ne sont pas logiquement indépendantes.
Par exemple, l'hypothèse (\ref{item-integrate}) implique
que $\partial_z(af - \theta(\partial_z(f))) = 0$.
Ainsi, si l'image de $B \to B'$ est l'ensemble des
éléments annulés par $\partial_z$, alors (\ref{item-factor})
en découle. Un raisonnement analogue s'applique à la condition (\ref{item-horizontal}).

\begin{lemma}
\label{lemma-find-homotopy}
Dans la situation ci-dessus, il existe un morphisme de complexes
$$
e_M^\bullet :
M \otimes_B (\Omega')^\bullet
\longrightarrow
M \otimes_B \Omega^\bullet
$$
tel que $c_M^\bullet \circ e_M^\bullet$
et $e_M^\bullet \circ c_M^\bullet$ soient homotopes à
la multiplication par $a$.
\end{lemma}

\begin{proof}
Dans cette démonstration, tous les produits tensoriels sont pris sur $B$.
L'hypothèse (\ref{item-direct-sum}) implique que
$$
M \otimes (\Omega')^i =
(B' \otimes M \otimes \Omega^i)
\oplus
(B' \text{d}z \otimes M \otimes \Omega^{i - 1})
$$
pour tout $i \geq 0$. Une famille de générateurs additifs de
$M \otimes (\Omega')^i$ est formée d'éléments de la forme
$f \omega$ et d'éléments de la forme $f \text{d}z \wedge \eta$,
où $f \in B'$, $\omega \in M \otimes \Omega^i$, et
$\eta \in M \otimes \Omega^{i - 1}$.

\medskip\noindent
Pour $f \in B'$, écrivons
$$
\epsilon(f) = af - \theta(\partial_z(f))
\quad\text{et}\quad
\epsilon'(f) = (\theta \otimes 1)(\text{d}_1(f)) - \text{d}_1(\theta(f))
$$
de sorte que $\epsilon(f) \in B$ et $\epsilon'(f) \in \Omega$ en vertu des
hypothèses (\ref{item-factor}) et (\ref{item-horizontal}).
Définissons $e_M^\bullet$ par les formules
$e^i_M(f\omega) = \epsilon(f) \omega$ et
$e^i_M(f \text{d}z \wedge \eta) = \epsilon'(f) \wedge \eta$.
Nous verrons ci-dessous que les applications $e^i_M$ définissent un morphisme de complexes.

\medskip\noindent
Définissons
$$
h^i : M \otimes_B (\Omega')^i \longrightarrow M \otimes_B (\Omega')^{i - 1}
$$
par les formules $h^i(f \omega) = 0$ et
$h^i(f \text{d}z \wedge \eta) = \theta(f) \eta$
pour les éléments ci-dessus. Nous affirmons que
$$
\text{d} \circ h + h \circ \text{d} = a - c_M^\bullet \circ e_M^\bullet
$$
Notons que la multiplication par $a$ est un morphisme de complexes
d'après (\ref{item-d-a-zero}). Ainsi, puisque $c_M^\bullet$ est un morphisme injectif
de complexes en vertu de l'hypothèse (\ref{item-injective}), on conclut que
$e_M^\bullet$ est un morphisme de complexes. Pour démontrer l'affirmation, calculons
\begin{align*}
(\text{d} \circ h + h \circ \text{d})(f \omega)
& =
h\left(\text{d}(f) \wedge \omega + f \nabla(\omega)\right)
\\
& =
\theta(\partial_z(f)) \omega
\\
& =
a f\omega - \epsilon(f)\omega 
\\
& =
a f \omega - c^i_M(e^i_M(f\omega))
\end{align*}
La deuxième égalité vient de ce que $\text{d}z$ n'intervient pas dans $\nabla(\omega)$,
et la troisième résulte de l'hypothèse (6). De même, on a
\begin{align*}
(\text{d} \circ h + h \circ \text{d})(f \text{d}z \wedge \eta)
& =
\text{d}(\theta(f) \eta) +
h\left(\text{d}(f) \wedge \text{d}z \wedge \eta -
f \text{d}z \wedge \nabla(\eta)\right)
\\
& =
\text{d}(\theta(f)) \wedge \eta + \theta(f) \nabla(\eta)
- (\theta \otimes 1)(\text{d}_1(f)) \wedge \eta
- \theta(f) \nabla(\eta)
\\
& =
\text{d}_1(\theta(f)) \wedge \eta +
\partial_z(\theta(f)) \text{d}z \wedge \eta -
(\theta \otimes 1)(\text{d}_1(f)) \wedge \eta
\\
& =
a f \text{d}z \wedge \eta - \epsilon'(f) \wedge \eta \\
& = a f \text{d}z \wedge \eta - c^i_M(e^i_M(f \text{d}z \wedge \eta))
\end{align*}
La deuxième égalité vient de ce que
$\text{d}(f) \wedge \text{d}z \wedge \eta =
- \text{d}z \wedge \text{d}_1(f) \wedge \eta$.
La quatrième égalité résulte de l'hypothèse (\ref{item-integrate}).
D'autre part, il résulte immédiatement des définitions
que $e^i_M(c^i_M(\omega)) = \epsilon(1) \omega = a \omega$.
Cela démontre le lemme.
\end{proof}

\begin{example}
\label{example-integrate}
Un exemple standard de la situation ci-dessus se présente lorsque
$B' = B\langle z \rangle$ est l'algèbre polynomiale à puissances divisées
sur un anneau à puissances divisées $(B, J, \delta)$, munie des puissances divisées
$\delta'$ sur $J' = B'_{+} + JB' \subset B'$. Plus précisément, prenons
$\Omega = \Omega_{B, \delta}$ et $\Omega' = \Omega_{B', \delta'}$.
Dans ce cas, on peut prendre $a = 1$ et
$$
\theta( \sum b_m z^{[m]} ) = \sum b_m z^{[m + 1]}
$$
Notons que
$$
f - \theta(\partial_z(f)) = f(0)
$$
est égal au terme constant. Il s'ensuit que, dans ce cas, le
lemme \ref{lemma-find-homotopy}
redonne le lemme de Poincar\'e cristallin
(lemme \ref{lemma-relative-poincare}).
\end{example}

\begin{lemma}
\label{lemma-computation}
Plaçons-nous dans la situation \ref{situation-affine}. Supposons que $D$ et $\Omega_D$ soient comme dans
(\ref{equation-D}) et (\ref{equation-omega-D}).
Soit $\lambda \in D$. Soit $D'$ le complété $p$-adique de
$$
D[z]\langle \xi \rangle/(\xi - (z^p - \lambda))
$$
et soit $\Omega_{D'}$ le complété $p$-adique du module des
différentielles à puissances divisées de $D'$ sur $A$. Pour tout couple $(M, \nabla)$
sur $D$ vérifiant (\ref{item-complete}), (\ref{item-connection}),
(\ref{item-integrable}) et (\ref{item-topologically-quasi-nilpotent}),
l'application canonique de complexes (\ref{equation-base-change-map-complexes})
$$
c_M^\bullet : M \otimes_D^\wedge \Omega^\bullet_D
\longrightarrow
M \otimes_D^\wedge \Omega^\bullet_{D'}
$$
a la propriété suivante : il existe un morphisme $e_M^\bullet$
dans le sens opposé tel que $c_M^\bullet \circ e_M^\bullet$
et $e_M^\bullet \circ c_M^\bullet$ soient toutes deux homotopes à la multiplication par $p$.
\end{lemma}

\begin{proof}
Nous allons le démontrer à l'aide du lemme \ref{lemma-find-homotopy} avec $a = p$.
Il nous faut donc trouver $\theta : D' \to D'$ et vérifier
(\ref{item-d-a-zero}), (\ref{item-direct-sum}), (\ref{item-theta-linear}),
(\ref{item-integrate}), (\ref{item-injective}), (\ref{item-factor}),
(\ref{item-horizontal}). Commençons par rassembler quelques informations sur les anneaux
$D$ et $D'$ et les modules $\Omega_D$ et $\Omega_{D'}$.

\medskip\noindent
En écrivant
$$
D[z]\langle \xi \rangle/(\xi - (z^p - \lambda))
=
D\langle \xi \rangle[z]/(z^p - \xi - \lambda)
$$
on voit que $D'$ est le complété $p$-adique du $D$-module libre
$$
\bigoplus\nolimits_{i = 0, \ldots, p - 1}
\bigoplus\nolimits_{n \geq 0}
z^i \xi^{[n]} D
$$
où $\xi^{[0]} = 1$.
Il s'ensuit que $D \to D'$ admet une section $D$-linéaire continue ; en particulier,
$D \to D'$ est universellement injectif, c'est-à-dire que (\ref{item-injective}) est vérifiée.
Considérons $D'$ comme une algèbre à puissances divisées
sur $A$, d'idéal à puissances divisées $\overline{J}' = \overline{J}D' + (\xi)$.
Alors $D'$ est aussi le complété $p$-adique de l'enveloppe à puissances divisées
de l'idéal engendré par $z^p - \lambda$ dans $D$, voir le
lemme \ref{lemma-describe-divided-power-envelope}. Ainsi,
$$
\Omega_{D'} = \Omega_D \otimes_D^\wedge D' \oplus D'\text{d}z
$$
d'après le lemme \ref{lemma-module-differentials-divided-power-envelope}.
Cela démontre (\ref{item-direct-sum}). Notons que (\ref{item-d-a-zero})
est évidente.

\medskip\noindent
À ce stade, construisons $\theta$. (Nous avons écrit un script PARI/gp theta.gp
qui vérifie certaines des formules de cette démonstration et se trouve dans le
sous-répertoire scripts du projet Stacks.) Avant cela, calculons
la différentielle des éléments $z^i \xi^{[n]}$. On a
$\text{d}z^i = i z^{i - 1} \text{d}z$. Pour $n \geq 1$, on a
$$
\text{d}\xi^{[n]} =
\xi^{[n - 1]} \text{d}\xi =
- \xi^{[n - 1]}\text{d}\lambda + p z^{p - 1} \xi^{[n - 1]}\text{d}z
$$
car $\xi = z^p - \lambda$. Pour $0 < i < p$ et $n \geq 1$, on a
\begin{align*}
\text{d}(z^i\xi^{[n]})
& =
iz^{i - 1}\xi^{[n]}\text{d}z + z^i\xi^{[n - 1]}\text{d}\xi \\
& =
iz^{i - 1}\xi^{[n]}\text{d}z + z^i\xi^{[n - 1]}\text{d}(z^p - \lambda) \\
& =
- z^i\xi^{[n - 1]}\text{d}\lambda +
(iz^{i - 1}\xi^{[n]} + pz^{i + p - 1}\xi^{[n - 1]})\text{d}z \\
& =
- z^i\xi^{[n - 1]}\text{d}\lambda +
(iz^{i - 1}\xi^{[n]} + pz^{i - 1}(\xi + \lambda)\xi^{[n - 1]})\text{d}z \\
& =
- z^i\xi^{[n - 1]}\text{d}\lambda +
((i + pn)z^{i - 1}\xi^{[n]} + p\lambda z^{i - 1}\xi^{[n - 1]})\text{d}z
\end{align*}
la dernière égalité résultant de $\xi \xi^{[n - 1]} = n\xi^{[n]}$.
On voit donc que
\begin{align*}
\partial_z(z^i) & = i z^{i - 1} \\
\partial_z(\xi^{[n]}) & = p z^{p - 1} \xi^{[n - 1]} \\
\partial_z(z^i\xi^{[n]}) & =
(i + pn) z^{i - 1} \xi^{[n]} + p \lambda z^{i - 1}\xi^{[n - 1]}
\end{align*}
Motivés par ces formules, définissons $\theta$ par les règles
$$
\begin{matrix}
\theta(z^j)
& = & p\frac{z^{j + 1}}{j + 1}
& j = 0, \ldots p - 1, \\
\theta(z^{p - 1}\xi^{[m]})
& = & \xi^{[m + 1]}
& m \geq 1, \\
\theta(z^j \xi^{[m]})
& = &
\frac{p z^{j + 1} \xi^{[m]} - \theta(p\lambda z^j \xi^{[m - 1]})}{(j + 1 + pm)}
& 0 \leq j < p - 1, m \geq 1
\end{matrix}
$$
où, dans la dernière ligne, nous procédons par récurrence sur $m$ pour définir notre choix de
$\theta$. En développant, on obtient (pour $0 \leq j < p - 1$ et $1 \leq m$)
$$
\theta(z^j \xi^{[m]}) =
\textstyle{\frac{p z^{j + 1} \xi^{[m]}}{(j + 1 + pm)} -
\frac{p^2 \lambda z^{j + 1} \xi^{[m - 1]}}{(j + 1 + pm)(j + 1 + p(m - 1))} +
\ldots +
\frac{(-1)^m p^{m + 1} \lambda^m z^{j + 1}}
{(j + 1 + pm) \ldots (j + 1)}}
$$
bien que nous n'utilisions pas cette expression ci-dessous. Il est clair que $\theta$
se prolonge de manière unique en une application $p$-adiquement continue et $D$-linéaire sur $D'$.
Par construction, on a (\ref{item-theta-linear}) et (\ref{item-integrate}).
Il reste à démontrer (\ref{item-factor}) et (\ref{item-horizontal}).

\medskip\noindent
Démonstration de (\ref{item-factor}) et (\ref{item-horizontal}).
Puisque $\theta$ est $D$-linéaire et continue, il suffit de montrer que
$p - \theta \circ \partial_z$,
resp.\ $(\theta \otimes 1) \circ \text{d}_1 - \text{d}_1 \circ \theta$
prend ses valeurs dans $D$, resp.\ dans $\Omega_D$, lorsqu'on l'évalue sur les
éléments $z^i\xi^{[n]}$\footnote{Cela peut se vérifier par un calcul direct :
on constate que $p - \theta \circ \partial_z$, évalué en
$z^i\xi^{[n]}$, est nul, sauf en $1$, qui est envoyé sur $p$, et en
$\xi$, qui est envoyé sur $-p\lambda$. On constate que
$(\theta \otimes 1) \circ \text{d}_1 - \text{d}_1 \circ \theta$
évalué en $z^i\xi^{[n]}$ est nul, sauf en $z^{p - 1}\xi$,
qui est envoyé sur $-\lambda$.}.
Posons $D_0 = \mathbf{Z}_{(p)}[\lambda]$ et
$D_0' = \mathbf{Z}_{(p)}[z, \lambda]\langle \xi \rangle/(\xi - z^p + \lambda)$.
Observons que chacune des expressions ci-dessus est un élément de
$D_0'$ ou de $\Omega_{D_0'}$. Il suffit donc de démontrer le résultat
dans le cas de $D_0 \to D_0'$. Notons que $D_0$ et $D_0'$
sont des anneaux sans torsion et que $D_0 \otimes \mathbf{Q} = \mathbf{Q}[\lambda]$
et $D'_0 \otimes \mathbf{Q} = \mathbf{Q}[z, \lambda]$.
Ainsi, $D_0 \subset D'_0$ est le sous-anneau des éléments annulés
par $\partial_z$, et (\ref{item-factor})
résulte de (\ref{item-integrate}) ; voir la discussion qui précède immédiatement le
lemme \ref{lemma-find-homotopy}. De même, on a
$\text{d}_1(f) = \partial_\lambda(f)\text{d}\lambda$, donc
$$
\left((\theta \otimes 1) \circ \text{d}_1 - \text{d}_1 \circ \theta\right)(f)
=
\left(\theta(\partial_\lambda(f)) - \partial_\lambda(\theta(f))\right)
\text{d}\lambda
$$
En appliquant $\partial_z$ au coefficient, on obtient
\begin{align*}
\partial_z\left(
\theta(\partial_\lambda(f)) - \partial_\lambda(\theta(f))
\right)
& =
p \partial_\lambda(f) - \partial_z(\partial_\lambda(\theta(f))) \\
& =
p \partial_\lambda(f) - \partial_\lambda(\partial_z(\theta(f))) \\
& =
p \partial_\lambda(f) - \partial_\lambda(p f) = 0
\end{align*}
de sorte que le coefficient ne dépend pas de $z$, comme souhaité.
Cela achève la démonstration du lemme.
\end{proof}

\noindent
Notons qu'un revêtement $\alpha_p$ itéré $X' \to X$ (tel que défini dans
l'introduction de cette section) est fini localement libre. Ainsi, si $X$ est
connexe, le degré de $X' \to X$ est constant et est une puissance de $p$.

\begin{lemma}
\label{lemma-pullback-along-p-power-cover}
Soit $p$ un nombre premier. Soit $(S, \mathcal{I}, \gamma)$ un schéma à puissances
divisées sur $\mathbf{Z}_{(p)}$ tel que $p \in \mathcal{I}$. Posons
$S_0 = V(\mathcal{I}) \subset S$. Soit $f : X' \to X$ un revêtement
$\alpha_p$ itéré de schémas sur $S_0$, de degré constant $q$. Soit
$\mathcal{F}$ un cristal quelconque en faisceaux quasi-cohérents sur $X$, et posons
$\mathcal{F}' = f_{\text{cris}}^*\mathcal{F}$.
Dans le triangle distingué
$$
Ru_{X/S, *}\mathcal{F}
\longrightarrow
f_*Ru_{X'/S, *}\mathcal{F}'
\longrightarrow
E
\longrightarrow
Ru_{X/S, *}\mathcal{F}[1]
$$
l'objet $E$ a ses faisceaux de cohomologie annulés par $q$.
\end{lemma}

\begin{proof}
Notons que $X' \to X$ est un homéomorphisme ; on peut donc identifier les espaces
topologiques sous-jacents à $X$ et $X'$. La question est manifestement locale sur $X$ ;
on peut donc supposer $X$, $X'$ et $S$ affines, et $X' \to X$ donné comme une
composée
$$
X' = X_n \to X_{n - 1} \to X_{n - 2} \to \ldots \to X_0 = X
$$
où chaque morphisme $X_{i + 1} \to X_i$ est un revêtement $\alpha_p$.
Notons $\mathcal{F}_i$ l'image inverse de $\mathcal{F}$ sur $X_i$.
Il suffit de démontrer que chacune des applications
$$
R\Gamma(\text{Cris}(X_i/S), \mathcal{F}_i)
\longrightarrow
R\Gamma(\text{Cris}(X_{i + 1}/S), \mathcal{F}_{i + 1})
$$
s'insère dans un triangle dont le troisième terme a ses groupes de cohomologie annulés
par $p$. (On utilise ici l'axiome TR4 de la catégorie triangulée $D(X)$. Les détails sont
omis.)

\medskip\noindent
On peut donc supposer que $S = \Spec(A)$, $X = \Spec(C)$, $X' = \Spec(C')$
et $C' = C[z]/(z^p - c)$ pour un certain $c \in C$. Choisissons une algèbre polynomiale
$P$ sur $A$ et une surjection $P \to C$. Soit $D$ le complété $p$-adique de
l'enveloppe à puissances divisées de $\Ker(P \to C)$ dans $P$, comme dans (\ref{equation-D}).
Posons $P' = P[z]$, muni de la surjection $P' \to C'$ qui envoie $z$ sur la classe de $z$
dans $C'$. Choisissons un relèvement $\lambda \in D$ de $c \in C$. On voit alors que
le complété $p$-adique de l'enveloppe à puissances divisées $D'$ de
$\Ker(P' \to C')$ dans $P'$ est isomorphe au complété $p$-adique de
$D[z]\langle \xi \rangle/(\xi - (z^p - \lambda))$ ; voir le
lemme \ref{lemma-computation} et sa démonstration.
On voit donc que le résultat découle de ce lemme
par le calcul de la cohomologie des cristaux en modules quasi-cohérents de la
proposition \ref{proposition-compute-cohomology-crystal}.
\end{proof}

\noindent
La borne du lemme suivant n'est probablement pas optimale.

\begin{lemma}
\label{lemma-pullback-along-p-power-cover-cohomology}
Avec les notations et hypothèses du
lemme \ref{lemma-pullback-along-p-power-cover},
l'application
$$
f^* :
H^i(\text{Cris}(X/S), \mathcal{F})
\longrightarrow
H^i(\text{Cris}(X'/S), \mathcal{F}')
$$
a son noyau et son conoyau annulés par $q^{i + 1}$.
\end{lemma}

\begin{proof}
Cela résulte du fait que $E$ n'a de faisceaux de cohomologie non nuls qu'en
degrés $-1$ et supérieurs, de sorte que la suite spectrale
$H^a(\mathcal{H}^b(E)) \Rightarrow H^{a + b}(E)$ converge.
En combinant cela avec la suite exacte longue de cohomologie associée
à un triangle distingué, on obtient la borne.
\end{proof}

\noindent
Dans la situation \ref{situation-global}, supposons que $p \in \mathcal{I}$.
Posons
$$
X^{(1)} = X \times_{S_0, F_{S_0}} S_0.
$$
Notons $F_{X/S_0} : X \to X^{(1)}$ le morphisme de Frobenius relatif.

\begin{lemma}
\label{lemma-pullback-relative-frobenius}
Dans la situation ci-dessus, supposons que $X \to S_0$ soit lisse de dimension
relative $d$. Alors $F_{X/S_0}$ est un revêtement $\alpha_p$ itéré
de degré $p^d$. Les lemmes \ref{lemma-pullback-along-p-power-cover} et
\ref{lemma-pullback-along-p-power-cover-cohomology} s'appliquent donc à cette
situation. En particulier, pour tout cristal en modules quasi-cohérents
$\mathcal{G}$ sur $\text{Cris}(X^{(1)}/S)$, l'application
$$
F_{X/S_0}^* : H^i(\text{Cris}(X^{(1)}/S), \mathcal{G})
\longrightarrow
H^i(\text{Cris}(X/S), F_{X/S_0, \text{cris}}^*\mathcal{G})
$$
a son noyau et son conoyau annulés par $p^{d(i + 1)}$.
\end{lemma}

\begin{proof}
Il suffit de démontrer le premier énoncé. Pour cela, on peut supposer
que $X$ est \'etale sur $\mathbf{A}^d_{S_0}$ ; voir
Morphismes, lemme \ref{morphisms-lemma-smooth-etale-over-affine-space}.
Notons $\varphi : X \to \mathbf{A}^d_{S_0}$ ce morphisme \'etale.
Dans ce cas, le Frobenius relatif de $X/S_0$ s'insère dans un diagramme
$$
\xymatrix{
X \ar[d] \ar[r] & X^{(1)} \ar[d] \\
\mathbf{A}^d_{S_0} \ar[r] & \mathbf{A}^d_{S_0}
}
$$
où la flèche horizontale inférieure est le morphisme de Frobenius relatif
de $\mathbf{A}^d_{S_0}$ sur $S_0$. C'est le morphisme qui élève
toutes les coordonnées à la puissance $p$-ième ; c'est donc un revêtement
$\alpha_p$ itéré. La démonstration s'achève en observant que le diagramme
est un carré cartésien ; voir
Morphismes \'etales, lemme \ref{etale-lemma-relative-frobenius-etale}.
\end{proof}













\section{Action de Frobenius sur la cohomologie cristalline}
\label{section-frobenius}

\noindent
Dans cette section, nous démontrons que l'image inverse par Frobenius induit un quasi-isomorphisme
sur la cohomologie cristalline après inversion du nombre premier $p$. Mais, ne serait-ce que
pour formuler ce résultat, nous devons nous placer dans une situation particulière.

\begin{situation}
\label{situation-F-crystal}
Dans la situation \ref{situation-global}, supposons ce qui suit :
\begin{enumerate}
\item $S = \Spec(A)$ pour un certain anneau à puissances divisées $(A, I, \gamma)$
avec $p \in I$,
\item on se donne un homomorphisme d'anneaux à puissances divisées $\sigma : A \to A$
tel que $\sigma(x) = x^p \bmod pA$ pour tout $x \in A$.
\end{enumerate}
\end{situation}

\noindent
Dans la situation \ref{situation-F-crystal}, le morphisme
$\Spec(\sigma) : S \to S$ est un relèvement du Frobenius absolu
$F_{S_0} : S_0 \to S_0$, et le diagramme
$$
\xymatrix{
X \ar[d] \ar[r]_{F_X} & X \ar[d] \\
S_0 \ar[r]^{F_{S_0}} & S_0
}
$$
est commutatif, où $F_X : X \to X$ est le morphisme de Frobenius absolu
de $X$. On obtient donc un morphisme de topos cristallins
$$
(F_X)_{\text{cris}} :
(X/S)_{\text{cris}}
\longrightarrow
(X/S)_{\text{cris}}
$$
Ce morphisme est décrit dans la remarque \ref{remark-functoriality-cris}. Voici la terminologie relative aux
$F$-cristaux, suivant la notation de Saavedra ; voir
\cite{Saavedra}.

\begin{definition}
\label{definition-F-crystal}
Dans la situation \ref{situation-F-crystal}, un {\it $F$-cristal sur $X/S$
(relatif à $\sigma$)} est une paire $(\mathcal{E}, F_\mathcal{E})$
formée d'un cristal en $\mathcal{O}_{X/S}$-modules localement libres de type fini
$\mathcal{E}$ et d'un morphisme
$$
F_\mathcal{E} : (F_X)_{\text{cris}}^*\mathcal{E} \longrightarrow \mathcal{E}
$$
Un $F$-cristal est dit {\it non dégénéré} s'il existe un entier
$i \geq 0$ et un morphisme $V : \mathcal{E} \to (F_X)_{\text{cris}}^*\mathcal{E}$
tel que $V \circ F_{\mathcal{E}} = p^i \text{id}$.
\end{definition}

\begin{remark}
\label{remark-F-crystal-variants}
Soit $(\mathcal{E}, F)$ un $F$-cristal comme dans la
définition \ref{definition-F-crystal}.
Dans la littérature, la condition de non-dégénérescence fait souvent partie de la
définition d'un $F$-cristal. En outre, on suppose aussi souvent que
$F \circ V = p^n\text{id}$. Ce qui est nécessaire pour le résultat ci-dessous, c'est
qu'il existe un entier $j \geq 0$ tel que $\Ker(F)$ et
$\Coker(F)$ soient annulés par $p^j$. Si le rang de $\mathcal{E}$
est borné (par exemple si $X$ est quasi-compact), ces deux
conditions résultent de la condition de non-dégénérescence telle qu'elle est formulée dans
la définition. En effet, soient $R$ un anneau, $r \geq 1$ un entier et
$K, L \in \text{Mat}(r \times r, R)$ des matrices telles que
$K L = p^i 1_{r \times r}$. Alors $\det(K)\det(L) = p^{ri}$. 
Soit $L'$ la comatrice de $L$, c'est-à-dire
$L' L = L L' = \det(L)$. Posons $K' = p^{ri} K$ et $j = ri + i$.
On a alors $K' L = p^j 1_{r \times r}$ puisque $K L = p^i$, et
$$
L K' = L K \det(L) \det(M) = L K L L' \det(M) = L p^i L' \det(M) =
p^j 1_{r \times r}
$$
Il s'ensuit que, si $V$ est comme dans la définition \ref{definition-F-crystal},
alors, en posant $V' = p^N V$, où $N > i \cdot \text{rang}(\mathcal{E})$,
on obtient $V' \circ F = p^{N + i}$ et $F \circ V' = p^{N + i}$.
\end{remark}

\begin{theorem}
\label{theorem-cohomology-F-crystal}
Dans la situation \ref{situation-F-crystal}, soit $(\mathcal{E}, F_\mathcal{E})$
un $F$-cristal non dégénéré. Supposons que $A$ soit un anneau complet pour la topologie $p$-adique
et noethérien, et que $X \to S_0$ soit propre et lisse. Alors
le morphisme canonique
$$
F_\mathcal{E} \circ (F_X)_{\text{cris}}^* :
R\Gamma(\text{Cris}(X/S), \mathcal{E}) \otimes^\mathbf{L}_{A, \sigma} A
\longrightarrow
R\Gamma(\text{Cris}(X/S), \mathcal{E})
$$
devient un isomorphisme après inversion de $p$.
\end{theorem}

\begin{proof}
Écrivons d'abord la flèche comme composée de trois flèches.
Plus précisément, posons
$$
X^{(1)} = X \times_{S_0, F_{S_0}} S_0
$$
et notons $F_{X/S_0} : X \to X^{(1)}$ le morphisme de Frobenius relatif.
Notons $\mathcal{E}^{(1)}$ le changement de base de $\mathcal{E}$
par $\Spec(\sigma)$, autrement dit l'image inverse de $\mathcal{E}$
sur $\text{Cris}(X^{(1)}/S)$ par le morphisme de topos cristallins
associé au diagramme commutatif
$$
\xymatrix{
X^{(1)} \ar[r] \ar[d] & X \ar[d] \\
S \ar[r]^{\Spec(\sigma)} & S
}
$$
On a alors le morphisme de changement de base
\begin{equation}
\label{equation-base-change-sigma}
R\Gamma(\text{Cris}(X/S), \mathcal{E}) \otimes^\mathbf{L}_{A, \sigma} A
\longrightarrow
R\Gamma(\text{Cris}(X^{(1)}/S), \mathcal{E}^{(1)})
\end{equation}
Ce morphisme est celui de la remarque \ref{remark-base-change}. Notons que la composée
de $F_{X/S_0} : X \to X^{(1)}$ avec la projection $X^{(1)} \to X$
est le morphisme de Frobenius absolu $F_X$. On voit donc que
$F_{X/S_0}^*\mathcal{E}^{(1)} = (F_X)_{\text{cris}}^*\mathcal{E}$.
Ainsi, l'image inverse par $F_{X/S_0}$ donne un morphisme
\begin{equation}
\label{equation-to-prove}
F_{X/S_0}^* :
R\Gamma(\text{Cris}(X^{(1)}/S), \mathcal{E}^{(1)})
\longrightarrow
R\Gamma(\text{Cris}(X/S), (F_X)^*_{\text{cris}}\mathcal{E})
\end{equation}
Enfin, on peut utiliser $F_\mathcal{E}$ pour obtenir un morphisme
\begin{equation}
\label{equation-F-E}
R\Gamma(\text{Cris}(X/S), (F_X)^*_{\text{cris}}\mathcal{E})
\longrightarrow
R\Gamma(\text{Cris}(X/S), \mathcal{E})
\end{equation}
Le morphisme du théorème est la composée des trois morphismes
(\ref{equation-base-change-sigma}), (\ref{equation-to-prove}) et
(\ref{equation-F-E}) ci-dessus. Le premier est un
quasi-isomorphisme modulo toute puissance de $p$ d'après la
remarque \ref{remark-base-change-isomorphism}.
C'est donc un quasi-isomorphisme puisque les complexes en jeu sont parfaits
dans $D(A)$ ; voir la remarque \ref{remark-complete-perfect}.
Le troisième morphisme est un quasi-isomorphisme après inversion de $p$, simplement
parce que $F_\mathcal{E}$ admet un inverse à une puissance de $p$ près ; voir la
remarque \ref{remark-F-crystal-variants}.
Enfin, le deuxième est un isomorphisme après inversion de $p$
d'après le lemme \ref{lemma-pullback-relative-frobenius}.
\end{proof}






\begin{multicols}{2}[\section{Autres chapitres}]
\noindent
Préliminaires
\begin{enumerate}
\item \hyperref[introduction-section-phantom]{Introduction}
\item \hyperref[conventions-section-phantom]{Conventions}
\item \hyperref[sets-section-phantom]{Théorie des ensembles}
\item \hyperref[categories-section-phantom]{Catégories}
\item \hyperref[topology-section-phantom]{Topologie}
\item \hyperref[sheaves-section-phantom]{Faisceaux sur les espaces}
\item \hyperref[sites-section-phantom]{Sites et faisceaux}
\item \hyperref[stacks-section-phantom]{Champs}
\item \hyperref[fields-section-phantom]{Corps}
\item \hyperref[algebra-section-phantom]{Algèbre commutative}
\item \hyperref[brauer-section-phantom]{Groupes de Brauer}
\item \hyperref[homology-section-phantom]{Algèbre homologique}
\item \hyperref[derived-section-phantom]{Catégories dérivées}
\item \hyperref[simplicial-section-phantom]{Méthodes simpliciales}
\item \hyperref[more-algebra-section-phantom]{Compléments d'algèbre}
\item \hyperref[smoothing-section-phantom]{Lissage des morphismes d'anneaux}
\item \hyperref[modules-section-phantom]{Faisceaux de modules}
\item \hyperref[sites-modules-section-phantom]{Modules sur les sites}
\item \hyperref[injectives-section-phantom]{Objets injectifs}
\item \hyperref[cohomology-section-phantom]{Cohomologie des faisceaux}
\item \hyperref[sites-cohomology-section-phantom]{Cohomologie sur les sites}
\item \hyperref[dga-section-phantom]{Algèbre différentielle graduée}
\item \hyperref[dpa-section-phantom]{Algèbre à puissances divisées}
\item \hyperref[sdga-section-phantom]{Faisceaux différentiels gradués}
\item \hyperref[hypercovering-section-phantom]{Hyperrecouvrements}
\end{enumerate}
Schémas
\begin{enumerate}
\setcounter{enumi}{25}
\item \hyperref[schemes-section-phantom]{Schémas}
\item \hyperref[constructions-section-phantom]{Constructions de schémas}
\item \hyperref[properties-section-phantom]{Propriétés des schémas}
\item \hyperref[morphisms-section-phantom]{Morphismes de schémas}
\item \hyperref[coherent-section-phantom]{Cohomologie des schémas}
\item \hyperref[divisors-section-phantom]{Diviseurs}
\item \hyperref[limits-section-phantom]{Limites de schémas}
\item \hyperref[varieties-section-phantom]{Variétés}
\item \hyperref[topologies-section-phantom]{Topologies sur les schémas}
\item \hyperref[descent-section-phantom]{Descente}
\item \hyperref[perfect-section-phantom]{Catégories dérivées de schémas}
\item \hyperref[more-morphisms-section-phantom]{Compléments sur les morphismes}
\item \hyperref[flat-section-phantom]{Compléments sur la platitude}
\item \hyperref[groupoids-section-phantom]{Schémas en groupoïdes}
\item \hyperref[more-groupoids-section-phantom]{Compléments sur les schémas en groupoïdes}
\item \hyperref[etale-section-phantom]{Morphismes étales de schémas}
\end{enumerate}
Thèmes de théorie des schémas
\begin{enumerate}
\setcounter{enumi}{41}
\item \hyperref[chow-section-phantom]{Homologie de Chow}
\item \hyperref[intersection-section-phantom]{Théorie de l'intersection}
\item \hyperref[pic-section-phantom]{Schémas de Picard des courbes}
\item \hyperref[weil-section-phantom]{Théories cohomologiques de Weil}
\item \hyperref[adequate-section-phantom]{Modules adéquats}
\item \hyperref[dualizing-section-phantom]{Complexes dualisants}
\item \hyperref[duality-section-phantom]{Dualité pour les schémas}
\item \hyperref[discriminant-section-phantom]{Discriminants et différentes}
\item \hyperref[derham-section-phantom]{Cohomologie de de Rham}
\item \hyperref[local-cohomology-section-phantom]{Cohomologie locale}
\item \hyperref[algebraization-section-phantom]{Géométrie algébrique et formelle}
\item \hyperref[curves-section-phantom]{Courbes algébriques}
\item \hyperref[resolve-section-phantom]{Résolution des surfaces}
\item \hyperref[models-section-phantom]{Réduction semi-stable}
\item \hyperref[functors-section-phantom]{Foncteurs et morphismes}
\item \hyperref[equiv-section-phantom]{Catégories dérivées de variétés}
\item \hyperref[pione-section-phantom]{Groupes fondamentaux de schémas}
\item \hyperref[etale-cohomology-section-phantom]{Cohomologie étale}
\item \hyperref[crystalline-section-phantom]{Cohomologie cristalline}
\item \hyperref[proetale-section-phantom]{Cohomologie pro-étale}
\item \hyperref[relative-cycles-section-phantom]{Cycles relatifs}
\item \hyperref[more-etale-section-phantom]{Compléments de cohomologie étale}
\item \hyperref[trace-section-phantom]{La formule des traces}
\end{enumerate}
Espaces algébriques
\begin{enumerate}
\setcounter{enumi}{64}
\item \hyperref[spaces-section-phantom]{Espaces algébriques}
\item \hyperref[spaces-properties-section-phantom]{Propriétés des espaces algébriques}
\item \hyperref[spaces-morphisms-section-phantom]{Morphismes d'espaces algébriques}
\item \hyperref[decent-spaces-section-phantom]{Espaces algébriques décents}
\item \hyperref[spaces-cohomology-section-phantom]{Cohomologie des espaces algébriques}
\item \hyperref[spaces-limits-section-phantom]{Limites d'espaces algébriques}
\item \hyperref[spaces-divisors-section-phantom]{Diviseurs sur les espaces algébriques}
\item \hyperref[spaces-over-fields-section-phantom]{Espaces algébriques sur des corps}
\item \hyperref[spaces-topologies-section-phantom]{Topologies sur les espaces algébriques}
\item \hyperref[spaces-descent-section-phantom]{Descente et espaces algébriques}
\item \hyperref[spaces-perfect-section-phantom]{Catégories dérivées d'espaces}
\item \hyperref[spaces-more-morphisms-section-phantom]{Compléments sur les morphismes d'espaces}
\item \hyperref[spaces-flat-section-phantom]{Platitude sur les espaces algébriques}
\item \hyperref[spaces-groupoids-section-phantom]{Groupoïdes dans les espaces algébriques}
\item \hyperref[spaces-more-groupoids-section-phantom]{Compléments sur les groupoïdes dans les espaces}
\item \hyperref[bootstrap-section-phantom]{Amorçage}
\item \hyperref[spaces-pushouts-section-phantom]{Sommes amalgamées d'espaces algébriques}
\end{enumerate}
Thèmes de géométrie
\begin{enumerate}
\setcounter{enumi}{81}
\item \hyperref[spaces-chow-section-phantom]{Groupes de Chow des espaces}
\item \hyperref[groupoids-quotients-section-phantom]{Quotients de groupoïdes}
\item \hyperref[spaces-more-cohomology-section-phantom]{Compléments sur la cohomologie des espaces}
\item \hyperref[spaces-simplicial-section-phantom]{Espaces simpliciaux}
\item \hyperref[spaces-duality-section-phantom]{Dualité pour les espaces}
\item \hyperref[formal-spaces-section-phantom]{Espaces algébriques formels}
\item \hyperref[restricted-section-phantom]{Algébrisation des espaces formels}
\item \hyperref[spaces-resolve-section-phantom]{Résolution des surfaces revisitée}
\end{enumerate}
Théorie des déformations
\begin{enumerate}
\setcounter{enumi}{89}
\item \hyperref[formal-defos-section-phantom]{Théorie formelle des déformations}
\item \hyperref[defos-section-phantom]{Théorie des déformations}
\item \hyperref[cotangent-section-phantom]{Le complexe cotangent}
\item \hyperref[examples-defos-section-phantom]{Problèmes de déformation}
\end{enumerate}
Champs algébriques
\begin{enumerate}
\setcounter{enumi}{93}
\item \hyperref[algebraic-section-phantom]{Champs algébriques}
\item \hyperref[examples-stacks-section-phantom]{Exemples de champs}
\item \hyperref[stacks-sheaves-section-phantom]{Faisceaux sur les champs algébriques}
\item \hyperref[criteria-section-phantom]{Critères de représentabilité}
\item \hyperref[artin-section-phantom]{Axiomes d'Artin}
\item \hyperref[quot-section-phantom]{Espaces de Quot et de Hilbert}
\item \hyperref[stacks-properties-section-phantom]{Propriétés des champs algébriques}
\item \hyperref[stacks-morphisms-section-phantom]{Morphismes de champs algébriques}
\item \hyperref[stacks-limits-section-phantom]{Limites de champs algébriques}
\item \hyperref[stacks-cohomology-section-phantom]{Cohomologie des champs algébriques}
\item \hyperref[stacks-perfect-section-phantom]{Catégories dérivées de champs}
\item \hyperref[stacks-introduction-section-phantom]{Introduction aux champs algébriques}
\item \hyperref[stacks-more-morphisms-section-phantom]{Compléments sur les morphismes de champs}
\item \hyperref[stacks-geometry-section-phantom]{La géométrie des champs}
\end{enumerate}
Thèmes de théorie des espaces de modules
\begin{enumerate}
\setcounter{enumi}{107}
\item \hyperref[moduli-section-phantom]{Champs de modules}
\item \hyperref[moduli-curves-section-phantom]{Espaces de modules de courbes}
\end{enumerate}
Divers
\begin{enumerate}
\setcounter{enumi}{109}
\item \hyperref[examples-section-phantom]{Exemples}
\item \hyperref[exercises-section-phantom]{Exercices}
\item \hyperref[guide-section-phantom]{Guide bibliographique}
\item \hyperref[desirables-section-phantom]{Desiderata}
\item \hyperref[coding-section-phantom]{Style de programmation}
\item \hyperref[obsolete-section-phantom]{Obsolète}
\item \hyperref[fdl-section-phantom]{Licence de documentation libre GNU}
\item \hyperref[index-section-phantom]{Index généré automatiquement}
\end{enumerate}
\end{multicols}


\bibliography{my}
\bibliographystyle{amsalpha}

\end{document}
